\documentclass[11pt]{article}

\usepackage{amsmath,amssymb,amsthm,mathtools}
\usepackage{booktabs}
\usepackage{hyperref}

\title{Finite-core Volterra reductions for a Weyl-positive Riemann phase kernel}
\author{Draft research notes}
\date{June 1, 2026}

\newtheorem{theorem}{Theorem}[section]
\newtheorem{proposition}[theorem]{Proposition}
\newtheorem{lemma}[theorem]{Lemma}
\newtheorem{definition}[theorem]{Definition}
\newtheorem{conjecture}[theorem]{Conjecture}
\newtheorem{problem}[theorem]{Problem}
\newtheorem{remark}[theorem]{Remark}

\newcommand{\R}{\mathbb R}
\newcommand{\C}{\mathbb C}
\newcommand{\N}{\mathbb N}
\newcommand{\eps}{\varepsilon}
\newcommand{\ii}{\mathrm i}
\newcommand{\dd}{\,\mathrm d}
\newcommand{\sgn}{\operatorname{sgn}}
\newcommand{\supp}{\operatorname{supp}}
\newcommand{\rank}{\operatorname{rank}}
\newcommand{\Span}{\operatorname{span}}
\newcommand{\Kred}{K_{\mathrm{red}}}
\newcommand{\Cred}{C_{\mathrm{red}}}
\newcommand{\Tred}{T_{\mathrm{red}}}
\newcommand{\EL}{E_L}
\newcommand{\calH}{\mathcal H}
\newcommand{\calK}{\mathcal K}
\newcommand{\Beta}{\mathrm{Beta}}
\newcommand{\Mu}{\mathrm{M}}

\begin{document}
\maketitle

\begin{abstract}
We record a Weyl-positive approach to the Riemann phase kernel associated with
the even Riemann kernel $\Phi$.  The immediate analytic target is not a proof
of the Riemann hypothesis, but a concrete positivity theorem for a Weyl kernel
whose quantum characteristic function satisfies the Kastler--Loupias--
Miracle-Sole condition in all numerical tests performed so far.  Several
natural factorizations are ruled out.  In particular, the positive anti-Wick
density route fails by a local heat-deconvolution obstruction, and a number of
apparently plausible finite-core shortcuts fail by explicit counterexamples.

The surviving structure is a finite-core Volterra program upgraded to a
closed-trace quotient theorem for the full kernel.  We derive exact same-sign
finite-core formulae, a second-order theta identity
$\phi_n=(\partial_t^2-1/4)(e^{t/2}e^{-\pi n^2e^{2t}})$, a Volterra
boundary-plus-tail representation, and a quotient Schur factorization for the
normalized full-\(\Phi\) source/Volterra model.  The latest certificate closes
the active trace-range condition, the full-continuum source-inactive
domination, and the Douglas/Moore--Penrose Schur hypotheses.  What remains is
external to that normalized Schur certificate: the quotient-to-original Weyl
lift, uniform \(\omega\)-coverage for \(|\omega|<1/2\), and the final bridge
from Weyl/KLM positivity to the intended de Branges or RH-side formulation.
This draft separates proved identities, disproved routes, closed certificates,
and remaining equivalence gaps.
\end{abstract}

\tableofcontents

\section{Introduction}

The Riemann $\xi$-function admits a Fourier representation in terms of an even,
rapidly decreasing kernel $\Phi$.  A long-standing strategy, going back through
Polya and de Branges, is to replace the location of zeros of an entire function
by positivity of a Hilbert-space or operator kernel naturally attached to it.
The work recorded here follows that principle, but uses Weyl positivity and
the KLM quantum positive-type condition rather than a direct de Branges
Hermite-Biehler inequality.

The present draft is deliberately conservative.  It does not claim a proof of
the Riemann hypothesis.  It documents:
\begin{enumerate}
  \item exact formulae and reductions that have been proved;
  \item candidate positivity statements that have been disproved by explicit
        local or finite-dimensional tests;
  \item the closed normalized Weyl/Volterra quotient Schur certificate and the
        remaining external equivalence gaps.
\end{enumerate}

The main lesson is that the positive anti-Wick route is not the correct
factorization.  The viable route is Weyl/Moyal positivity: prove that the Weyl
operator kernel is positive, equivalently that the symplectic Fourier transform
of the Weyl symbol satisfies the KLM positivity condition.  After a parity
reduction this becomes a sharp contraction problem
\[
       A\ge 0,\qquad -A\le B\le A,
\]
or, in the reproducing-kernel Hilbert space generated by $A$,
\[
       |A^{-1/2}BA^{-1/2}|\le 1.
\]
The finite-core analysis and subsequent full-\(\Phi\) perturbation theory now
close the normalized Weyl/Volterra quotient Schur certificate.  The remaining
task is to identify that quotient certificate with the original Weyl kernel
quadratic form, prove the required uniform \(\omega\)-coverage, and then verify
the final de Branges/RH-facing bridge.

\section{The Riemann kernel and the Weyl symbol}

We use the positive-side Riemann kernel
\begin{equation}
  \Phi(x)=\sum_{n\ge 1}
  2\bigl(2\pi^2 n^4e^{9x/2}-3\pi n^2e^{5x/2}\bigr)
  e^{-\pi n^2e^{2x}},
  \qquad x\ge 0,
  \label{eq:Phi}
\end{equation}
and extend it evenly by $\Phi(-x)=\Phi(x)$.  Thus the corresponding Fourier
representation of $\xi$ may be written in the standard form
\[
  \xi\!\left(\frac12+\ii z\right)
   = \int_{-\infty}^{\infty}\Phi(t)e^{\ii zt}\dd t,
\]
up to the conventional normalization of $\xi$.

For $y\ge 0$ define
\begin{equation}
  \widehat C_y(\eta)
   =2\int_0^\infty \cos(\eta u)\Phi(y+u)\Phi(y-u)\dd u.
  \label{eq:Chat}
\end{equation}
The phase-space Weyl symbol under study is
\begin{equation}
  \sigma_\omega(x,\xi)
   =\int_{|x|}^{\infty}
      y\cosh(2\omega y)\widehat C_y(2\xi)\dd y,
  \qquad |\omega|<\frac12.
  \label{eq:sigma}
\end{equation}

The corresponding quantum characteristic function is
\begin{equation}
  Q_\omega(s,t)
  =
  \frac{
  2\pi\int_0^\infty
     y\cosh(2\omega y)
     \Phi(y+s/2)\Phi(y-s/2)
     \dfrac{\sin(ty)}{t}\dd y
  }{
  Q_\omega(0,0)
  },
  \label{eq:Q}
\end{equation}
where $\sin(ty)/t$ is interpreted as $y$ at $t=0$.

\begin{definition}[KLM test]
For finite sets $(s_j,t_j)\in\R^2$, the KLM matrix attached to
$Q_\omega$ is
\begin{equation}
   M_{jk}
   =
   Q_\omega(s_j-s_k,t_j-t_k)
   \exp\!\left\{\frac{\ii}{2}(t_js_k-s_jt_k)\right\}.
   \label{eq:KLM}
\end{equation}
The KLM condition is that every such matrix is positive semidefinite.
\end{definition}

In this normalization the KLM condition is equivalent to positivity of the
Weyl quantization of $\sigma_\omega$.  The associated coordinate kernel is
\begin{equation}
  K_\omega(a,b)
  =
  \frac12\int_{|(a+b)/2|}^{\infty}
   y\cosh(2\omega y)
   \Phi\!\left(y+\frac{a-b}{2}\right)
   \Phi\!\left(y-\frac{a-b}{2}\right)\dd y.
  \label{eq:WeylKernel}
\end{equation}

\begin{problem}[Main Weyl positivity target]
Prove that $K_\omega$ is a positive semidefinite kernel on $\R$ for
$|\omega|<1/2$.  Equivalently, prove that $Q_\omega$ satisfies the KLM
condition.
\end{problem}

\section{The anti-Wick route fails}

The first candidate factorization was anti-Wick positivity.  In the standard
Gaussian/Weyl normalization this asks for
\[
      m_\omega=e^{-c\Delta}\sigma_\omega\ge 0,
      \qquad c=\frac14.
\]
If true, $\sigma_\omega$ would be a Gaussian smoothing of a positive density,
giving immediate Weyl positivity.  This route fails already at the local
Taylor level.

\begin{proposition}[Numerical anti-Wick obstruction]
At $\omega=0.49$, near $x=0.3$, $\xi=0$, direct quadrature gives
\[
  \sigma_\omega(x,\xi)\approx 1.10125\cdot 10^{-3},
  \qquad
  \Delta\sigma_\omega(x,\xi)\approx 1.42523\cdot 10^{-1}.
\]
Hence even the mild deconvolution
\[
     \sigma_\omega-0.01\,\Delta\sigma_\omega
\]
is already negative, approximately $-3.24\cdot 10^{-4}$.
FFT deconvolution with moderate cutoffs also gives a negative anti-Wick
density near $(x,\xi)=(0.3125,0)$, with value about $-2.4$ in the tested
normalization.
\end{proposition}

\begin{remark}
This is not merely a failure of numerical resolution.  The sign of the first
heat-deconvolution correction is too large relative to $\sigma_\omega$ itself.
Thus positive anti-Wick density is not the right factorization for this
kernel.
\end{remark}

\section{Parity reduction and the contraction target}

On the half-line the full Weyl kernel separates into same-parity and reflected
cross terms.  The numerical evidence is strongest when formulated as a
contraction statement.

Let $A$ denote the same-sign half-line kernel and let $B$ denote the reflected
cross-kernel.  The Weyl positivity target is equivalent to
\begin{equation}
       A\ge 0,\qquad -A\le B\le A.
       \label{eq:ABtarget}
\end{equation}
When $A$ is not strictly positive, this is interpreted on the quotient by the
null space of $A$.  Equivalently,
\begin{equation}
       |A^{-1/2}BA^{-1/2}|\le 1.
       \label{eq:contraction}
\end{equation}
In reproducing-kernel language, the reflected cross-kernel must define a
self-adjoint contraction in the RKHS generated by $A$.

\begin{remark}[Why this formulation matters]
The ordinary kernel $B$ is not positive.  Direct finite-dimensional tests give
negative spectrum, for example minimum eigenvalues of order
$-5.2\cdot 10^{-4}$ at $\omega=0$ and $-7.0\cdot 10^{-4}$ at
$\omega=0.49$.  Thus the correct claim is not $B\ge 0$, but the sharper
relative inequality \eqref{eq:ABtarget}.
\end{remark}

For the full kernel $\Phi$, half-line parity tests on $x\in[0,8]$ with
$80$ sample points give minimum eigenvalues at roundoff:
\[
\begin{array}{c|cc}
\omega & \lambda_{\min}(\text{even}) & \lambda_{\min}(\text{odd})\\
\hline
0      & -9.6\cdot 10^{-29} & -6.6\cdot 10^{-34}\\
0.49   & 0\text{ to roundoff} & -1.3\cdot 10^{-39}.
\end{array}
\]
After projection away from the null space of $A$, the spectrum of
$A^{-1/2}BA^{-1/2}$ is numerically contained in $[-1,1]$ to roundoff.  At
$n=80$ the numerical rank of $A$ was $7$, and the observed spectral intervals
were approximately
\[
  [-0.9999999988,1.0000000000]\quad(\omega=0),
\]
and
\[
  [-0.9999999983,1.0000000000]\quad(\omega=0.49).
\]

\section{RKHS differentiation and the mixed kernel}

A useful formal reduction is that positivity of a sufficiently smooth kernel
implies positivity of its mixed derivative kernel.  This is elementary, but it
is the correct way to pass from a Gram formula for a Weyl kernel to the mixed
kernel appearing in the Volterra analysis.

\begin{lemma}[Derivative-kernel Gram lift]
Let $K$ be a $C^2$ positive semidefinite kernel on an interval $I$, and let
$\mathcal H_K$ be its RKHS.  Assume point-derivative evaluations are bounded.
Then
\[
     H(x,y)=\partial_x\partial_yK(x,y)
\]
is positive semidefinite on $I$.
\end{lemma}

\begin{proof}
Since $K(x,y)=\langle k_y,k_x\rangle_{\mathcal H_K}$, boundedness of
derivative evaluations gives vectors $\partial_x k_x$ in the RKHS, in the
weak derivative sense.  Therefore
\[
  \partial_x\partial_yK(x,y)
     =\langle \partial_y k_y,\partial_x k_x\rangle_{\mathcal H_K}.
\]
Finite matrices of $H$ are Gram matrices of the derivative kernel vectors.
\end{proof}

Consequently, a direct Gram formula for the same-sign Weyl kernel of a finite
core is enough to imply the mixed positivity needed in the contraction
argument.

\section{Finite theta cores and the zero-slope correction}

Write the positive-side theta modes as
\[
   \phi_n(t)
   =
   2\bigl(2\pi^2 n^4e^{9t/2}-3\pi n^2e^{5t/2}\bigr)
   e^{-\pi n^2e^{2t}},
   \qquad t\ge 0,
\]
so that $\Phi=\sum_{n\ge 1}\phi_n$ on the positive axis.  Let
\[
   \Phi_{\le N}=\sum_{n=1}^N\phi_n.
\]
The full even kernel has $\Phi'(0+)=0$, but finite truncations need not.
The first right derivatives at the origin are:
\[
\begin{array}{c|c}
N & \Phi_{\le N}'(0+)\\
\hline
1 & 3.949876526774993\cdot 10^{-2}\\
2 & 8.265279577727339\cdot 10^{-8}\\
3 & 1.391725863609305\cdot 10^{-16}\\
4 & 2.869736205910160\cdot 10^{-28}.
\end{array}
\]

The numerically natural finite core is therefore the zero-slope corrected
three-mode core
\begin{equation}
   \widetilde\Phi_3
   =
   \phi_1+\phi_2+\alpha_3\phi_3,
   \qquad
   \alpha_3
   =
   -\frac{\phi_1'(0+)+\phi_2'(0+)}{\phi_3'(0+)},
   \label{eq:tildePhi3}
\end{equation}
with
\[
   \alpha_3
   \approx
   1.000000001683821887529422431530891.
\]

Scaling tests with $\phi_1+\phi_2+\alpha\phi_3$ show that the zero-slope
constraint is active.  For $\omega=0.49$ on $x\in[-2.6,2.6]$ with $64$ sample
points, the minimum eigenvalue improves from roughly $-1.30\cdot 10^{-11}$ at
$\alpha=0$ to roundoff near $\alpha=1$, and deteriorates again when
$\alpha$ is moved beyond the zero-slope value.

The full theta tail beyond the three-mode core is extremely small on the
positive axis.  On $[0,8]$ the observed bounds were
\[
  \sup|\Phi-\widetilde\Phi_3|\approx 1.45\cdot 10^{-18},\qquad
  \sup|\Phi'-\widetilde\Phi_3'|\approx 1.39\cdot 10^{-16},
\]
and
\[
  \sup|\Phi''-\widetilde\Phi_3''|\approx 1.31\cdot 10^{-14}.
\]
This motivates the finite-core-plus-perturbation strategy.

\subsection{Endpoint cancellation in the mixed kernel}

The zero-slope condition is not cosmetic.  It removes a boundary-crossing
defect in the full-line mixed kernel.  Let $F$ denote a finite positive-side
core and set $f(t)=F(|t|)$.  The mixed source appearing after differentiating
the Weyl kernel has the form
\begin{equation}
\begin{aligned}
 S(s,t)
 &=
 W_\omega''(s+t) f(s)f(t)
 \\
 &\quad+
 W_\omega'(s+t)\{f'(s)f(t)+f(s)f'(t)\}
 +
 W_\omega(s+t)f'(s)f'(t),
\end{aligned}
\label{eq:mixed-source}
\end{equation}
where
\[
   W_\omega(u)=\frac14u\cosh(\omega u).
\]
For opposite-sign arguments, the reflected part contains integrals of the
form
\[
   R(x,y)=\int_0^\infty S(x+r,r-y)\dd r,
   \qquad x,y\ge 0,
\]
with the evident interchange when $y>x$.

If $F'(0+)=d\ne 0$, then $f'$ has a jump of size $2d$ at the origin.  Hence
the reflected source has a jump at the boundary-crossing point $r=y$:
\begin{equation}
 S(x+y,0+)-S(x+y,0-)
 =
 2d\{W_\omega'(x+y)F(x+y)+W_\omega(x+y)F'(x+y)\}.
 \label{eq:endpoint-jump}
\end{equation}
This jump vanishes exactly when $F'(0+)=0$.

More explicitly, put
\[
   A_F(s,t)
   =
   W_\omega'(s+t)F(s)+W_\omega(s+t)F'(s),
\]
and
\[
   S_{\mathrm{reg}}(s,t)
   =
   S(s,t)-d\,\operatorname{sgn}(t)A_F(s,t).
\]
For opposite-sign entries let $X=\max(|a|,|b|)$ and
$Y=\min(|a|,|b|)$.  Then the reflected mixed kernel decomposes exactly as
\begin{equation}
  \mathcal H_F(a,b)
  =
  \mathcal H_{F,\mathrm{reg}}(a,b)
  +
  d\,J_F(a,b),
  \label{eq:endpoint-decomp}
\end{equation}
where
\[
  \mathcal H_{F,\mathrm{reg}}(a,b)
  =
  \int_0^\infty S_{\mathrm{reg}}(X+r,r-Y)\dd r
\]
and
\[
  J_F(a,b)
  =
  \int_0^\infty
     \operatorname{sgn}(r-Y)A_F(X+r,r-Y)\dd r.
\]
As a distribution in the crossing variable $r$,
\[
  \frac{\dd}{\dd r}S(X+r,r-Y)
  =
  \left(\frac{\dd}{\dd r}S(X+r,r-Y)\right)_{\mathrm{classical}}
  +
  2d\,A_F(X+Y,0)\delta(r-Y).
\]
Thus $F'(0+)=0$ removes the endpoint delta and jump term.  For
$\widetilde\Phi_3$ this is exactly the role of the coefficient $\alpha_3$ in
\eqref{eq:tildePhi3}.

Numerically, this is the active finite-core mechanism.  At $\omega=0.49$ on
$[-2.6,2.6]$ with $64$ grid points, the two-mode mixed kernel has
\[
   \lambda_{\min}(\mathcal H_{\le 2})
   \approx -7.38288357\cdot 10^{-11}.
\]
On the corresponding witness vector, the true third-mode correction contributes
\[
  +7.45840143\cdot 10^{-11},
\]
leaving the three-mode value positive at about $7.55\cdot 10^{-13}$.  The
tail from $n\ge 4$ contributes only about $4.0\cdot 10^{-18}$ on this witness.
Thus the third mode should be understood as the endpoint-cancellation mode;
the far theta tail is perturbative only after the zero-slope core has been
formed.

\begin{conjecture}[Finite-core and tail strategy]
The proof of Weyl positivity can be split into:
\begin{enumerate}
  \item prove the Weyl kernel $K_{\widetilde\Phi_3,\omega}$ is positive;
  \item deduce the mixed positivity by the RKHS derivative-kernel lemma;
  \item bound $\Phi-\widetilde\Phi_3$ as a perturbative theta tail.
\end{enumerate}
\end{conjecture}

\section{Failed shortcuts}

Several plausible proof routes were tested and rejected.  We list them because
they sharply constrain the remaining proof.

\subsection{Local source positivity fails}

A local Green-source kernel $D_{\le 3}$ associated with the three-mode core was
tested on $[0,2.6]$ with $80$ points.  It is indefinite:
\[
\begin{array}{c|cc}
\omega & \lambda_{\min}(D_{\le 3}) & \min \operatorname{diag}(D_{\le 3})\\
\hline
0      & -3.45 & -0.507\\
0.49   & -3.39 & -0.491.
\end{array}
\]
Thus positivity cannot be proved by pointwise or layerwise positivity of this
local source.

\subsection{Hermite-Biehler shortcuts fail for finite cores}

The finite-core entire functions do not satisfy the needed
Hermite-Biehler inequality.  For $\widetilde\Phi_3$ at $\omega=0.49$, a
counterexample occurs near
\[
  z\approx 70+0.8854009689\,\ii,
\]
where
\[
  |\Xi_{\widetilde\Phi_3}(z+\ii\omega)|
   \approx 5.38\cdot 10^{-21},
  \qquad
  |\Xi_{\widetilde\Phi_3}(z-\ii\omega)|
   \approx 1.46\cdot 10^{-20}.
\]
Thus the expected inequality has the wrong sign at that point.  The raw
$\Phi_{\le 3}$ truncation also fails, for instance near
$z\approx 76.6667+4\ii$.  This is consistent with Polya's warning that finite
one-sided theta truncations do not inherit the full even entire-function
structure.

\subsection{Generic smooth-even positivity is false}

The target positivity is not a generic consequence of smoothness, evenness,
or Gaussian-like decay.  Model tests at $\omega=0.49$ on $x\in[-3,3]$ with
$81$ points showed that a Gaussian and low Hermite modifications pass to
roundoff, while a difference of Gaussians produces a negative eigenvalue of
order $-1.36\cdot 10^{-3}$.  The Riemann kernel requires its special theta
structure.

\subsection{First score integration by parts fails}

An attempted first-order score integration by parts, after splitting the
$\cosh$ weight, gives
\[
  K_\eps(x,y)
  =
  \frac1{16}\int_0^\infty
  [u+x+y]F_\eps(u+2x)F_\eps(u+2y)\dd u,
\]
with
\[
  a=-F'/F,\qquad Q=a(s)+a(t),\qquad
  C=\frac{s+t}{2Q},\qquad D=(\partial_s+\partial_t)C.
\]
The resulting $D$ is indefinite.  For $\widetilde\Phi_3$ at $\omega=0.49$,
the $\eps=+1$ branch has a minimum around $-2.17\cdot 10^3$ near
$(s,t)=(0.101,0)$, and the $\eps=-1$ branch has minimum around
$-3.32\cdot 10^{-2}$.  Therefore one cannot split the hyperbolic cosine and
prove positivity by this first-order score kernel.

\subsection{The finite anti-Loewner index guess is false}

The reduced one-mode anti-Loewner boundary kernel was expected to have exactly
two negative squares.  High-precision quadratic-grid tests show instead many
small negative eigenvalues.  For example, at $n=40$ and 100-digit precision,
the observed negative eigenvalues included
\[
 -3.56\cdot 10^{-1},\quad
 -6.39\cdot 10^{-6},\quad
 -1.96\cdot 10^{-10},\quad
 -5.01\cdot 10^{-15},
\]
continuing down to about $10^{-40}$.  The correct statement is not exact
finite negative index; it is effective dominance of the significant negative
directions by the Volterra tail.

\section{Exact same-sign formulae for finite cores}

The finite-core kernels admit exact integral formulae after expansion into
exponentials.  These formulae are important because naive quadrature can miss
endpoint layers where the diagonal is extremely small.

Let a finite positive-side core have the form
\[
   \phi(t)=\sum_i A_i e^{\lambda_i t-c_ie^{2t}},
   \qquad c_i>0.
\]
For $x,y\ge 0$, the same-sign Weyl kernel has the exact expansion
\begin{equation}
\begin{aligned}
 K_\phi(x,y)
 &=
 \frac1{16}\sum_{\eps=\pm1}\sum_{i,j}
 A_iA_j e^{(\lambda_i+\eps\omega)x}e^{(\lambda_j+\eps\omega)y}
 \\
 &\quad\cdot
 \left[
    I_{\log}(\alpha_{ij},p_{ij})
    +(x+y)I_0(\alpha_{ij},p_{ij})
 \right],
\end{aligned}
\label{eq:samesign}
\end{equation}
where
\[
   \alpha_{ij}=c_ie^{2x}+c_je^{2y},
   \qquad
   p_{ij}=\frac{\lambda_i+\lambda_j+2\eps\omega}{2},
\]
and
\[
   I_0(\alpha,p)=\int_1^\infty q^{p-1}e^{-\alpha q}\dd q
     =\alpha^{-p}\Gamma(p,\alpha),
\]
\[
   I_{\log}(\alpha,p)=\partial_p I_0(\alpha,p)
     =\int_1^\infty \log(q)q^{p-1}e^{-\alpha q}\dd q.
\]

\begin{proof}
On the same-sign half-line the lower integration endpoint becomes the boundary
corresponding to $u=0$.  Expanding both finite cores and the $\cosh$ factor
reduces the kernel to a sum of integrals of the form
\[
  \int_0^\infty (u+x+y)
  \exp\{(\lambda_i+\lambda_j+2\eps\omega)u/2\}
  \exp\{-\alpha_{ij}e^u\}\dd u,
\]
up to the displayed exponential prefactors and normalization.  The substitution
$q=e^u$ gives \eqref{eq:samesign}.
\end{proof}

\begin{remark}[Endpoint layer]
For $\widetilde\Phi_3$ the exact formula gives positive normalized minimum
eigenvalues around $4\cdot 10^{-9}$ at both $\omega=0$ and $\omega=0.49$ in
stress tests near $x\approx 2.18$.  Simpson quadrature had produced false
negative alarms because it missed endpoint layers with diagonal magnitude near
$10^{-208}$.
\end{remark}

\section{The second-order theta identity}

The strongest structure found so far is a second-order identity for each theta
mode.

\begin{lemma}[Second-order theta identity]
Let
\[
   b_n(t)=e^{t/2}e^{-\pi n^2e^{2t}}.
\]
Then
\begin{equation}
   \phi_n(t)=\left(\partial_t^2-\frac14\right)b_n(t).
   \label{eq:secondorder-t}
\end{equation}
Equivalently, with $v=2t$,
\[
   B_n(v)=e^{v/4}e^{-\pi n^2e^v},\qquad
   M=4\partial_v^2-\frac14,
\]
one has
\begin{equation}
   \phi_n(v/2)=MB_n(v).
   \label{eq:secondorder-v}
\end{equation}
\end{lemma}

\begin{proof}
Let $c=\pi n^2$ and $b(t)=e^{t/2-ce^{2t}}$.  Then
\[
  \frac{b'}b=\frac12-2ce^{2t},
\]
and
\[
  \frac{b''}b
  =
  \left(\frac12-2ce^{2t}\right)^2-4ce^{2t}
  =
  \frac14-6ce^{2t}+4c^2e^{4t}.
\]
Therefore
\[
  \left(\partial_t^2-\frac14\right)b(t)
  =
  (4c^2e^{4t}-6ce^{2t})e^{t/2-ce^{2t}},
\]
which is exactly
\[
  2(2\pi^2 n^4e^{9t/2}-3\pi n^2e^{5t/2})
  e^{-\pi n^2e^{2t}}.
\]
The $v$-form follows from $\partial_t=2\partial_v$.
\end{proof}

This identity changes the problem from a first-order score problem for
$\Phi$ to a second-order Volterra problem for the simpler base function
\[
   B(v)=B_1(v)+B_2(v)+\alpha_3B_3(v).
\]

\section{Volterra boundary-plus-tail representation}

Let
\[
   \Psi(v)=MB(v),
\]
so that $\Psi(v)=\widetilde\Phi_3(v/2)$ for the zero-slope corrected core.
For same-sign variables set
\[
   s=u+2x,\qquad t=u+2y,\qquad U=\frac{s+t}{2},
\]
and
\[
   W_\omega(s,t)=U\cosh(\omega U).
\]
Define the base score and normalized source by
\[
   a(v)=-\frac{B'(v)}{B(v)},\qquad
   \mu(v)=\frac{\Psi(v)}{B(v)},\qquad
   Q(s,t)=a(s)+a(t).
\]
Finally set
\begin{equation}
   C_\omega(s,t)
   =
   \frac{W_\omega(s,t)\mu(s)\mu(t)}{Q(s,t)},
   \qquad
   D_\omega(s,t)
   =
   (\partial_s+\partial_t)C_\omega(s,t).
   \label{eq:CD}
\end{equation}

\begin{proposition}[Exact Volterra identity]
For the same-sign kernel of the finite core,
\begin{equation}
\begin{aligned}
 K_\Psi(x,y)
 &=
 C_\omega(2x,2y)B(2x)B(2y)
 \\
 &\quad+
 \int_0^\infty
 D_\omega(u+2x,u+2y)
 B(u+2x)B(u+2y)\dd u.
\end{aligned}
\label{eq:Volterra}
\end{equation}
\end{proposition}

\begin{proof}
Along the Volterra ray $(s,t)=(u+2x,u+2y)$,
\[
   \frac{\dd}{\dd u}\{B(s)B(t)\}
   =
   -Q(s,t)B(s)B(t).
\]
By definition of $C_\omega$,
\[
   C_\omega(s,t)Q(s,t)B(s)B(t)
   =
   W_\omega(s,t)\Psi(s)\Psi(t).
\]
Also
\[
   \frac{\dd}{\dd u}
   \{C_\omega(s,t)B(s)B(t)\}
   =
   D_\omega(s,t)B(s)B(t)
   -
   W_\omega(s,t)\Psi(s)\Psi(t).
\]
Integrating from $u=0$ to $\infty$ and using the rapid decay at infinity gives
\eqref{eq:Volterra}.
\end{proof}

\begin{remark}
The individual boundary and fixed-$u$ source layers are indefinite.  For
$\widetilde\Phi_3$ on $x\in[0,2.6]$, tests gave boundary minimum eigenvalues
near $-1.29\cdot 10^{-2}$ and worst fixed-layer source eigenvalues near
$-2.65\cdot 10^{-2}$ at $\omega=0.49$.  Nevertheless the integrated Volterra
object and the direct same-sign kernel are positive to numerical roundoff.
This is the main reason the problem must be treated at operator level rather
than by pointwise source positivity.
\end{remark}

\section{Reduced boundary and the anti-Loewner form}

The endpoint-normalized kernel is
\[
   \Kred(s,t)=\frac{K(s/2,t/2)}{\Psi(s)\Psi(t)}.
\]
The reduced boundary term is
\begin{equation}
   \Cred(s,t)=\frac{W_\omega(s,t)}{a(s)+a(t)}.
   \label{eq:Cred}
\end{equation}
The Volterra remainder is
\[
   \Tred=\Kred-\Cred.
\]

After the change of variables $X=a(s)$, $Y=a(t)$, the boundary has the
anti-Loewner shape
\[
   \frac{f(X)+f(Y)}{X+Y},
   \qquad f=a^{-1},
\]
up to the smooth positive and hyperbolic-cosine weights.  This observation is
structural but not sufficient: as noted above, the reduced boundary has many
small negative eigenvalues.  The correct use of the anti-Loewner structure is
to isolate a low-dimensional macroscopic negative space and then prove that
the Volterra tail dominates the remaining negative boundary energy.

\subsection{Exact reduced finite-core formula}

For the finite core write
\[
   \Psi(v)=\sum_i A_i e^{\beta_i v-c_ie^v},
   \qquad
   r_i(s)=\frac{A_i e^{\beta_i s-c_ie^s}}{\Psi(s)}.
\]
Let $S=(s+t)/2$ and $\alpha_{ij}=c_ie^s+c_je^t$.  Then
\begin{equation}
\begin{aligned}
 \Kred(s,t)
 &=
 \sum_{i,j}r_i(s)r_j(t)
 \int_0^\infty
 (S+u)\cosh(\omega(S+u))
 \\
 &\qquad\qquad\cdot
 e^{(\beta_i+\beta_j)u}
 e^{-\alpha_{ij}(e^u-1)}\dd u.
\end{aligned}
\label{eq:Kred-u}
\end{equation}
Equivalently,
\begin{equation}
\begin{aligned}
 \Kred(s,t)
 &=
 \frac12\sum_{i,j}\sum_{\sigma=\pm1}
 r_i(s)r_j(t)e^{\sigma\omega S}e^{\alpha_{ij}}
 \\
 &\quad\cdot
 \left[
   S I_0(\alpha_{ij},p_{ij,\sigma})
   + I_{\log}(\alpha_{ij},p_{ij,\sigma})
 \right],
\end{aligned}
\label{eq:Kred-q}
\end{equation}
where
\[
  p_{ij,\sigma}=\beta_i+\beta_j-1+\sigma\omega.
\]
A stable Laguerre form is obtained from
\[
  e^\alpha\int_1^\infty F(q)e^{-\alpha q}\dd q
  =
  \alpha^{-1}\int_0^\infty F(1+r/\alpha)e^{-r}\dd r.
\]

\section{Finite-interval split certificate}

The current proof target is an operator split on a fixed interval
$[0,L]$, with $L=8$ in the numerical tests.  Let
\[
   \EL=\Span\{s-L/2,\ e^{-5s}\}.
\]
This two-dimensional space models the two macroscopic negative boundary modes.
The observed finite-dimensional certificate is:
\[
   \Kred|_{\EL^\perp}\ge 0,
   \qquad
   \operatorname{Schur}_{\EL}(\Kred)\ge 0,
\]
and on the negative spectral subspace of $\Cred|_{\EL^\perp}$,
\begin{equation}
   \langle v,\Tred v\rangle
   \ge
   \eta_L\langle v,-\Cred v\rangle,
   \qquad
   \eta_L>1.
   \label{eq:tail-dominates}
\end{equation}

\begin{conjecture}[Reduced finite-interval split lemma]
For $\widetilde\Phi_3$, $L=8$, and $|\omega|\le 0.49$, the decomposition
$\Kred=\Cred+\Tred$ satisfies the split positivity statement above.  Moreover
\eqref{eq:tail-dominates} holds with a uniform margin $\eta_L>1$.
\end{conjecture}

The numerical evidence is stable across several Galerkin discretizations.
For Gauss-Legendre Galerkin order up to $70$ with Laguerre order $160$, the
results were:
\[
\begin{array}{c|ccc}
\omega & \lambda_{\min}(\Kred|_{\EL^\perp})
       & \lambda_{\min}(\Cred|_{\EL^\perp})
       & \min \Tred/(-\Cred)\\
\hline
0.49 & \text{roundoff} & -3.45\cdot 10^{-4} & 3.9\text{ to }5.0\\
0    & \text{roundoff} & -5.45\cdot 10^{-4} & 3.7\text{ to }3.9.
\end{array}
\]
Legendre basis tests on $[0,8]$ gave the same qualitative picture.  For basis
orders $12$ and $24$, the reduced kernel was positive to roundoff, while the
tail-to-residual-boundary ratio remained above $4$ in the tested cases.

\begin{remark}[Why a crude finite matrix tail bound is not enough]
The reduced boundary $\Cred$ has fast Legendre coefficient decay.  At
$\omega=0.49$, the outside operator norm after cuts $16,24,32,40$ was roughly
\[
  1.92\cdot 10^{-6},\quad
  5.21\cdot 10^{-9},\quad
  1.54\cdot 10^{-11},\quad
  8.03\cdot 10^{-13}.
\]
By contrast $\Kred$ and $\Tred$ decay much more slowly, with outside norms
still around $10^{-3}$ to $10^{-2}$ at comparable cuts.  Therefore the final
proof should use fast decay only to localize the negative boundary subspace,
and should use the Volterra structure to prove positivity on the remaining
space.
\end{remark}

\section{Perturbing from the core to the full theta kernel}

The perturbation from $\widetilde\Phi_3$ to $\Phi$ is analytically small on
the positive axis.  However, positivity is delicate near null spaces, so a
naive operator-norm comparison is not enough.  The correct perturbative target
should be relative to the positive directions generated by the finite core and
to the contraction formulation \eqref{eq:contraction}.

Numerical core/tail energy tests on $x\in[-2.6,2.6]$ with $80$ points and
core rank $13$ were favorable.  At $\omega=0.49$, the theta tail restricted to
the core range had minimum around $-6.6\cdot 10^{-36}$ and maximum around
$7.2\cdot 10^{-29}$; the tail on the core null space was of order
$10^{-44}$ to $10^{-39}$, and the cross norm was about $4.2\cdot 10^{-34}$.
The $\omega=0$ tests were similar.

\begin{problem}[Tail perturbation lemma]
Prove that the kernel induced by $\Phi-\widetilde\Phi_3$ is form-small
relative to the positive finite-core kernel on the relevant quotient spaces,
uniformly for $|\omega|\le 0.49$.
\end{problem}

\section{What remains to prove}

The current proof program has narrowed to the following concrete lemmas.

\begin{problem}[Finite-core Weyl positivity]
Prove that the same-sign Weyl kernel of $\widetilde\Phi_3$ is positive on the
half-line.  Equivalently, prove the reduced split certificate for
$\Kred=\Cred+\Tred$.
\end{problem}

\begin{problem}[Boundary negative-space localization]
Prove that the negative spectral subspace of $\Cred$ outside
$\EL=\Span\{s-L/2,e^{-5s}\}$ is controlled by a rapidly decaying tail, in a
form compatible with the Volterra domination estimate.
\end{problem}

\begin{problem}[Volterra domination]
On the negative spectral subspace of $\Cred|_{\EL^\perp}$, prove
\[
  \langle v,\Tred v\rangle
   \ge \eta_L\langle v,-\Cred v\rangle
\]
with $\eta_L>1$, uniformly in the target range of $\omega$.
\end{problem}

\begin{problem}[Finite-dimensional Schur complement]
Prove the positivity of the Schur complement on $\EL$:
\[
   \operatorname{Schur}_{\EL}(\Kred)\ge 0.
\]
The numerical evidence now indicates that this should be treated as a singular
Schur-complement problem.  If
\[
   D=\Kred|_{\EL^\perp},
\]
then the needed statement is
\[
   D\ge 0,\qquad
   P_{\ker D}\Kred \EL=0,
\]
together with
\[
   K_{\EL\EL}-K_{\EL D}D^+K_{D\EL}\ge 0
\]
in the Moore-Penrose sense.  The observed margin is cutoff-sensitive because
the complement block has eigenvalues at and below $10^{-12}$ in Legendre
coordinates.  Thus the proof should not rely on a strictly positive numerical
gap in this two-dimensional block.
\end{problem}

\begin{lemma}[Range condition from a common Gram map]
Let $H=L^2([0,L])$ and write $H=E\oplus M$, with $M=E^\perp$.  Let $T$ be a
symmetric integral operator with kernel $K$.  Suppose there is a Hilbert space
$\mathcal G$ and a linear map $\Gamma:H\to\mathcal G$ such that
\[
   \langle f,Tg\rangle_H
   =
   \langle \Gamma f,\Gamma g\rangle_{\mathcal G}
\]
for all test functions in $E+M$.  Let $D=P_MT|_M$ be the compression to $M$.
Then
\[
   P_{\ker D}TE=0.
\]
Moreover the Moore-Penrose Schur form satisfies
\[
   K_{EE}-K_{EM}D^+K_{ME}\ge 0.
\]
\end{lemma}

\begin{proof}
If $v\in\ker D$, then
\[
   0=\langle v,Dv\rangle_H=\langle v,Tv\rangle_H
    =\|\Gamma v\|_{\mathcal G}^2.
\]
Hence $\Gamma v=0$.  For every $e\in E$,
\[
   \langle v,Te\rangle_H
   =
   \langle \Gamma v,\Gamma e\rangle_{\mathcal G}
   =
   0,
\]
which is the range condition.

For the Schur form, let $\Gamma_M$ denote the restriction of $\Gamma$ to $M$,
and let $R=\overline{\operatorname{ran}\Gamma_M}$.  For fixed $e\in E$,
\[
   \inf_{m\in M}\langle e+m,T(e+m)\rangle_H
   =
   \inf_{m\in M}\|\Gamma e+\Gamma_Mm\|_{\mathcal G}^2
   =
   \|(I-P_R)\Gamma e\|_{\mathcal G}^2\ge 0.
\]
This infimum is exactly the Moore-Penrose Schur quadratic form.  Thus the
singular Schur complement is positive in the quotient sense.
\end{proof}

\begin{remark}[What the Volterra Gram must prove]
The direct reduced Volterra formula has the schematic form
\[
  \Kred(s,t)
  =
  \int_0^\infty
  \bigl(u+(s+t)/2\bigr)
  \cosh(\omega(u+(s+t)/2))
  A_s(u)A_t(u)\dd u,
  \qquad
  A_s(u)=\frac{\Psi(s+u)}{\Psi(s)}.
\]
At $\omega=0$ one may split
\[
  u+\frac{s+t}{2}
  =
  u+\min(s,t)+\frac{|s-t|}{2}.
\]
The first term is a genuine Brownian Volterra Gram kernel, since
\[
  u+\min(s,t)
  =
  \int_0^\infty
  {\bf 1}_{r\le u+s}{\bf 1}_{r\le u+t}\dd r.
\]
The residual $|s-t|A_s(u)A_t(u)/2$ is strongly indefinite in finite tests.
Similarly, at $\omega>0$, splitting the hyperbolic cosine into exponential
branches destroys a cancellation: the $e^{+\omega}$ branch is positive in the
tested finite core, while the $e^{-\omega}$ branch is not.  Thus the remaining
Gram theorem cannot be proved layer-by-layer or branch-by-branch; it must keep
the Volterra integral and the full cosh weight intact.

Once the Volterra domination theorem proves that the quadratic form
\[
  Q(f,g)=\int_0^L\int_0^L f(s)\Kred(s,t)g(t)\dd s\dd t
\]
is nonnegative on $E+M$, a common Gram map follows by the standard
Kolmogorov construction:
\[
  \mathcal G=\overline{(E+M)/\{f:Q(f,f)=0\}},
  \qquad
  \Gamma f=[f],
\]
with
\[
  \langle \Gamma f,\Gamma g\rangle_{\mathcal G}=Q(f,g).
\]
Therefore the hard analytic statement is not the abstract existence of the
Gram map; it is the Volterra domination inequality that makes $Q$ positive.
\end{remark}

\begin{remark}[Moment-transform form of the residual domination]
The domination of the $|s-t|$ residual can be written as a one-dimensional
moment-transform inequality.  For each cosh branch $\sigma=\pm1$, define
\[
  B_\sigma(s,u)
  =
  e^{\sigma\omega(s+u)/2}A_s(u),
\]
and, for a test function $f$,
\[
  m_\sigma(u)=\int_0^L f(s)B_\sigma(s,u)\dd s,
  \qquad
  n_\sigma(u)=\int_0^L (s+u)f(s)B_\sigma(s,u)\dd s.
\]
Then the reduced Volterra quadratic form is exactly
\[
  Q(f,f)
  =
  \frac12\sum_{\sigma=\pm1}
  \int_0^\infty m_\sigma(u)n_\sigma(u)\dd u,
\]
with the obvious single-branch interpretation when $\omega=0$.

Indeed, expanding $m_\sigma n_\sigma$ and symmetrizing in $s,t$ gives the
factor $u+(s+t)/2$; summing $\sigma=\pm1$ gives the full hyperbolic cosine
weight.

Equivalently, for fixed $u$ and branch $\sigma$, let
\[
  d\nu_u(s)=f(s)B_\sigma(s,u)\dd s,\qquad
  F_u(r)=\nu_u([r,L]).
\]
Then
\[
  \int\!\!\int \bigl(u+\min(s,t)\bigr)\dd\nu_u(s)\dd\nu_u(t)
  =
  uF_u(0)^2+\int_0^L F_u(r)^2\dd r,
\]
while
\[
  \frac12\int\!\!\int |s-t|\dd\nu_u(s)\dd\nu_u(t)
  =
  F_u(0)\int_0^L F_u(r)\dd r
  -
  \int_0^L F_u(r)^2\dd r.
\]
The Brownian square cancels exactly against the negative part of the
$|s-t|$ residual, leaving
\[
  m_\sigma(u)n_\sigma(u).
\]
Thus the remaining Volterra domination theorem is precisely positivity of the
combined moment form
\[
  \frac12\sum_{\sigma=\pm1}
  \int_0^\infty m_\sigma(u)n_\sigma(u)\dd u\ge 0.
\]
For $\omega>0$ this positivity is not branchwise: numerical tests show the
$\sigma=-1$ branch is indefinite, while the sum over both branches is positive
to roundoff.
\end{remark}

\begin{proposition}[Dilation monotonicity implies the moment inequality]
For each branch set
\[
  M_{\sigma,\lambda}f(u)
  =
  \int_0^L e^{\lambda(s+u)}f(s)B_\sigma(s,u)\dd s
\]
and
\[
  \mathcal N_\lambda(f)
  =
  \frac12\sum_{\sigma=\pm1}
  \int_0^\infty |M_{\sigma,\lambda}f(u)|^2\dd u .
\]
Assume that $\mathcal N_{\lambda_2}-\mathcal N_{\lambda_1}$ is nonnegative
on the test space whenever $0\le \lambda_1\le \lambda_2$.  Then
\[
  \frac12\sum_{\sigma=\pm1}
  \int_0^\infty m_\sigma(u)n_\sigma(u)\dd u\ge0 .
\]
\end{proposition}

\begin{proof}
Differentiating under the integral gives
\[
  \left.\partial_\lambda M_{\sigma,\lambda}f(u)\right|_{\lambda=0}
  =
  \int_0^L (s+u)f(s)B_\sigma(s,u)\dd s
  =
  n_\sigma(u),
\]
while $M_{\sigma,0}f=m_\sigma$.  Hence
\[
  \mathcal N'_0(f)
  =
  \sum_{\sigma=\pm1}
  \int_0^\infty m_\sigma(u)n_\sigma(u)\dd u .
\]
Right monotonicity of $\mathcal N_\lambda$ at $\lambda=0$ gives
$\mathcal N'_0(f)\ge0$, which is the claimed moment inequality.  The
$\omega=0$ convention is obtained either by using the duplicated identical
branches or by replacing the factor $1/2\sum_{\sigma=\pm1}$ with the single
branch.
\end{proof}

\begin{problem}[Dilation monotonicity of the finite-core transform]
Prove, for the zero-slope three-mode core and $|\omega|\le0.49$, that
\[
  \mathcal N_{\lambda_2}-\mathcal N_{\lambda_1}\ge0
  \qquad(0\le\lambda_1\le\lambda_2)
\]
on $E+M$.  Numerically this Loewner monotonicity holds at the stress values
$\omega=0$ and $\omega=0.49$, and the finite-difference identity
$(\mathcal N_h-\mathcal N_0)/h\to2Q$ recovers the moment matrix.
\end{problem}

\begin{remark}[A false Fubini proof]
Changing variables to the center
\[
  x=u+\frac{s+t}{2}
\]
gives the layer representation
\[
  \int_{(s+t)/2}^\infty
  W(x)
  \frac{\Psi(x+(s-t)/2)\Psi(x-(s-t)/2)}
       {\Psi(s)\Psi(t)}
  \dd x .
\]
It would be enough for arbitrary positive weights $W$ if the fixed-$x$
triangular layer were positive.  Numerically it is not: for
$\widetilde\Phi_3$ on $[0,8]$, the fixed-center layer has negative eigenvalues
of size $10^{-2}$ to $10^{-1}$.  Hence the proof must use the special tilted
weights, or an equivalent global Volterra cancellation, not positivity of
center layers.
\end{remark}

\begin{problem}[Branch domination]
Let $a=\omega/2$ and write the derivative form as
\[
  \frac12\mathcal N'_0=\frac12(Q_a+Q_{-a}),
\]
where
\[
  Q_a(f)=\int_0^\infty m_a(u)n_a(u)\dd u,
  \quad
  m_a(u)=\int_0^L f(s)e^{a(s+u)}A_s(u)\dd s .
\]
Since $\Psi(s)>0$ on the positive half-line, the reduced normalization is only
a positive diagonal conjugation.  Thus $Q_a\ge0$ is equivalent to positivity of
the unnormalised Hankel--Volterra kernel
\[
  L_{g_a}(s,t)=
  \int_0^\infty
  \left(u+\frac{s+t}{2}\right)g_a(s+u)g_a(t+u)\dd u,
  \qquad
  g_a(v)=e^{av}\Psi(v).
\]
The numerical evidence indicates that $Q_a\ge0$ for $a\ge0$, while $Q_{-a}$
has only two visible negative directions and is dominated there by $Q_a$.
For $\widetilde\Phi_3$, $a=0.245$, the generalized ratio
$Q_a/(-Q_{-a})$ on the negative subspace of $Q_{-a}$ is about $3.68$ at
both $L=8$ and $L=12$.  Proving this positive-branch theorem plus reflected
branch domination would prove the base derivative $\mathcal N'_0\ge0$.
With a spectral floor $10^{-10}$, the reflected obstruction reduces to one
endpoint-localized negative direction and the Schur block of the combined form
has margin about $2.7\cdot10^{-7}$ in the current Galerkin tests.

The one-mode positive-branch theorem cannot be proved by extending the
multiplicative variable to the full Mellin line.  For
$g_a(v)=2c e^{pv}(2ce^v-3)e^{-ce^v}$, $p=a+5/4$, the full-line Mellin
multiplier is proportional to
\[
  c^{-p+i\tau}\Gamma(p-i\tau)\bigl(2(p-i\tau)-3\bigr),
\]
and the derivative symbol is negative near $\tau=0$.  Thus the endpoint
$e^s\ge1$ is essential in the one-mode proof.

Equivalently, with $S=e^s$, $r=e^u$, and
\[
  k_p(x)=2cx^p(2cx-3)e^{-cx},
\]
the one-mode theorem is monotonicity in $p$ of the truncated Mellin transform
\[
  (T_pF)(r)=\int_1^\infty k_p(rS)F(S)\frac{\dd S}{S},
  \qquad r\ge1,
\]
namely
\[
  \frac12\partial_p\|T_pF\|_{L^2([1,\infty),\dd r/r)}^2\ge0,
  \qquad p\ge5/4 .
\]
The derivative splits as
\[
  \frac12\partial_p\|T_pF\|^2
  =
  \langle T_pF,(\log r)T_pF\rangle
  +
  \Re\langle T_pF,T_p((\log S)F)\rangle .
\]
The first term is positive from the output endpoint $r\ge1$; the second is an
indefinite input anti-commutator.  Current Galerkin tests show the positive
output term dominates the negative input part with ratio at least about $1.34$
for the one-mode branch at $a=0.245$.

At the gamma-kernel level the endpoint appears through the exact commutator
\[
  D H_p + H_pD = -b_p(r)\operatorname{ev}_0,
  \qquad
  b_p(x)=e^{px-ce^x},
\]
where $D=d/dx$ and
\[
  (H_pf)(r)=\int_0^\infty b_p(r+s)f(s)\dd s .
\]
The one-mode kernel is
\[
  k_p=2c\bigl[(2p-3)-2D\bigr]b_p.
\]
Thus the endpoint Hardy inequality should follow from a boundary-square
identity for this commutator; this final algebraic closure is the remaining
open step in the one-mode proof.

Expanding with $A=2p-3$, $y=H_pF$, $\eta=H_p(LF)$, and
$z=(A-2D)y$, one obtains, up to the harmless factor $(2c)^2$,
\[
  E(F)=\langle z,Lz\rangle+\langle z,(A-2D)\eta\rangle .
\]
The half-line integrations by parts
\[
  \langle y,LDy\rangle+\langle Dy,Ly\rangle=-\|y\|^2,
  \qquad
  \langle y,D\eta\rangle+\langle Dy,\eta\rangle=-y(0)\eta(0)
\]
give
\[
  E(F)
  =
  A^2\bigl(\langle y,Ly\rangle+\langle y,\eta\rangle\bigr)
  +2A\|y\|^2
  +2Ay(0)\eta(0)
  +4\bigl(\langle Dy,LDy\rangle+\langle Dy,D\eta\rangle\bigr).
\]
The remaining obstruction is that $\eta=H_p(LF)$ is not eliminated by the
boundary commutator.  In particular, for $p<3/2$ the two terms linear in $A$
can have the wrong sign, so the final proof needs an additional Hardy identity
or Schur estimate controlling the constrained $\eta$ terms.

The corrected endpoint target can be stated without hiding this constraint.
Put
\[
  \mathfrak H_p(F)
   =\langle Dy,LDy\rangle+\langle Dy,D\eta\rangle
\]
and
\[
  \mathfrak L_p(F)
   =
   A^2\bigl(\langle y,Ly\rangle+\langle y,\eta\rangle\bigr)
   +2A\|y\|^2+2Ay(0)\eta(0).
\]
Then
\[
  E(F)=\mathfrak L_p(F)+4\mathfrak H_p(F).
\]
Moreover
\[
  \langle y,Ly\rangle+\langle y,\eta\rangle
  =\frac12\partial_p\|y\|^2
\]
and
\[
  \mathfrak H_p(F)
  =
  \frac12\partial_p\|D H_pF\|^2
  +\frac12|H_pF(0)|^2.
\]
Indeed,
$\partial_p Dy=y+LDy+D\eta$ and
$\langle Dy,y\rangle=-|y(0)|^2/2$.  Thus the boundary square is present, but
it is not sufficient by itself.  Finite Galerkin tests with
\texttt{dy\_hardy\_split.py} show that $\mathfrak H_p$ has a stable small
negative direction at the hard endpoint $p=5/4$, while
$\mathfrak L_p+4\mathfrak H_p$ remains positive to roundoff.  On $[0,12]$ with
output cutoff $20$, the ratio of $4\mathfrak H_p$ to $-\mathfrak L_p$ on the
negative spectral subspace of $\mathfrak L_p$ is about $3.65$ at $p=5/4$,
$7.89$ at $p=1.375$, and $2.75\cdot10^2$ at $p=1.495$.  The remaining Schur
floor at $p=5/4$ is of order $10^{-13}$, indicating again that the proof
should use a singular Schur/range condition rather than a determinant gap.
Scanning $p$ also shows that the threshold is sharp for this reduction:
below $p=5/4$ the full expanded form develops genuine negative directions
($\min E_p\approx -3.6\cdot10^{-10}$ at $p=1.20$ and
$-2.0\cdot10^{-7}$ at $p=1.00$ in the same discretization), while on
$[5/4,2.25]$ the worst ratio and Schur floor occur at $p=5/4$.

There is also an exact square completion.  With
\[
  W_p=(2D-A)y,\qquad \Theta_p=(2D-A)\eta,
\]
one has
\[
  E(F)
  =
  \int_0^\infty r\,W_p(r)^2\,\dd r
  +\int_0^\infty W_p(r)\Theta_p(r)\,\dd r .
\]
This follows from
\[
 \int r(2Dy-Ay)^2
 =4\langle Dy,LDy\rangle+A^2\langle y,Ly\rangle+2A\|y\|^2
\]
and
\[
 \int (2Dy-Ay)(2D\eta-A\eta)
 =4\langle Dy,D\eta\rangle+A^2\langle y,\eta\rangle
  +2Ay(0)\eta(0).
\]
Since $W_p=-T_pF$ and $\Theta_p=-T_p(LF)$ up to the harmless constant, the
one-mode theorem is equivalently the singular range inequality
\[
  \langle T_pF,LT_pF\rangle+\langle T_pF,T_pLF\rangle\ge0,
\]
or formally
\[
  L+T_pLT_p^{-1}\ge0
\]
on $\operatorname{Ran}T_p$.  The compactness of $T_p$ makes this a genuine
range/Hardy condition at the endpoint $r=0$, not a bounded inverse statement.

At $p=5/4$ this range condition has a cleaner multiplicative normal form.  Let
$r,S\ge1$,
\[
  a(x)=x^{3/2}e^{-cx},\qquad A(S)=S^{-1/4}F(S).
\]
Since
\[
  k_{5/4}(x)
  =
  2cx^{5/4}(2cx-3)e^{-cx}
  =
  -4cx^{3/4}\frac{d}{dx}\bigl(x^{3/2}e^{-cx}\bigr),
\]
one obtains
\[
  T_{5/4}F(r)
  =
  4cr^{-1/4}\int_1^\infty h_r(S)A(S)\,\dd S,
\]
where
\[
  h_r(S)
  =
  -\frac{d}{dS}a(rS)
  =
  r(crS-3/2)(rS)^{1/2}e^{-crS}.
\]
For $c=\pi$ this kernel is strictly positive on $r,S\ge1$.  Moreover
\[
  \left.\partial_pT_pF(r)\right|_{p=5/4}
  =
  4cr^{-1/4}\int_1^\infty \log(rS)h_r(S)A(S)\,\dd S.
\]
Thus the endpoint theorem is equivalent to
\[
 \int_1^\infty r^{-3/2}
 \left(\int_1^\infty h_r(S)A(S)\,\dd S\right)
 \left(\int_1^\infty \log(rS)h_r(S)A(S)\,\dd S\right)\dd r
 \ge0.
\]
The extension to $p=5/4+\delta$ is the same inequality with
$A$ replaced by $S^\delta A$ and the outside weight replaced by
$r^{2\delta-3/2}\dd r$.  This is the precise positive-kernel Hardy/Chebyshev
lemma still to be proved.

A direct layerwise proof is false: decomposing the positive weight into
threshold layers gives indefinite fixed layers.  Numerically the layer
$\alpha=0.5$ has minimum eigenvalue about $-1.96\cdot10^{-3}$ in the endpoint
one-mode model.  Hence the proof must use the integrated $r$-structure, not
pointwise layer positivity.

Finally, the endpoint form can be symmetrized.  On
$L^2([1,\infty),x^{-1/2}\dd x)$ define
\[
  \kappa(x,y)=xy(cxy-3/2)e^{-cxy}.
\]
Let $K$ be the associated self-adjoint integral operator and let
$L$ denote multiplication by $\log x$.  Then the endpoint inequality is
equivalent to
\[
  \langle KA,LKA\rangle+\langle KA,KLA\rangle\ge0,
\]
or, as a symmetric quadratic form,
\[
  KLK+\frac12(K^2L+LK^2)\ge0.
\]
The operator $K$ itself is indefinite; numerical tests on $[1,12]$ give seven
negative eigenvalues for $K$, while the anti-commutator form is positive to
roundoff.  This is the final self-adjoint form of the endpoint range problem.

The endpoint boundary also has an exact commutator expression.  Let
\[
  A_0=x\frac{d}{dx}+\frac14
\]
on $L^2(x^{-1/2}\dd x)$.  Then $[A_0,L]=1$, and integration by parts gives
\[
  A_0K+KA_0=-\kappa(\cdot,1)\operatorname{ev}_1 .
\]
On the full Mellin line the right side is absent, which is exactly the false
unrestricted argument.  Thus the endpoint theorem reduces to converting this
rank-one dilation-commutator defect into
\[
  KLK+\frac12(K^2L+LK^2)\ge0.
\]

In logarithmic variables this becomes a Hankel transport problem.  The unitary
map $(Uf)(t)=e^{t/4}f(e^t)$ sends $K$ to
\[
  (Hf)(t)=\int_0^\infty h(t+s)f(s)\,\dd s,\qquad
  h(u)=e^{5u/4}(ce^u-3/2)e^{-ce^u},
\]
while $A_0$ becomes $\partial_t$ and $L$ becomes multiplication by $t$.  The
anti-commutator kernel is
\[
  P(t,s)=\int_0^\infty
  \left(r+\frac{t+s}{2}\right)h(t+r)h(s+r)\,\dd r
\]
and obeys the exact transport equation
\[
  (\partial_t+\partial_s)P(t,s)
  =
  -\frac{t+s}{2}h(t)h(s).
\]
The source kernel on the right is indefinite; its quadratic form is the product
of the two endpoint moments
\[
  \left(\int h(t)f(t)\,\dd t\right)
  \left(\int th(t)f(t)\,\dd t\right).
\]
Thus the endpoint defect does not directly factor as a positive square.  The
remaining plausible proof is a Green-kernel total-positivity theorem: the
Hankel kernel $h(t+s)$ is numerically sign-regular with signature
$(-1)^{n(n-1)/2}$, while $P(t,s)$ is numerically totally positive.  Proving
this sign-regular Green theorem would imply $P\ge0$ from nonnegative principal
minors and close the endpoint case.

The determinant mechanism is slightly sharper than plain sign-regularity of
$h(t+s)$.  Define
\[
  f_0(t,r)=h(t+r),\qquad f_1(t,r)=(t+r)h(t+r).
\]
Then
\[
  P(t,s)=\frac12\int_0^\infty
  \{f_1(t,r)f_0(s,r)+f_0(t,r)f_1(s,r)\}\,\dd r .
\]
Thus $P$ is a Gram kernel over the disjoint union
$[0,\infty)\times\{0,1\}$.  By the vector-valued Andreief--Cauchy--Binet
formula, $P$ is totally positive if, for every type word
$\alpha\in\{0,1\}^n$,
\[
  \det\bigl[f_{\alpha_j}(t_i,r_j)\bigr]_{i,j=1}^n
\]
has sign $(-1)^{n(n-1)/2}$ whenever
$t_1<\cdots<t_n$ and $r_1<\cdots<r_n$.  In that case the two determinants in
the Cauchy--Binet integrand have the same sign, so their product is
nonnegative pointwise, and all minors of $P$ are nonnegative.

Numerically this extended sign-regularity holds for every mixed type word
through order $5$ on the tested grids.  However, the hoped-for all-order
Laguerre sign-regularity theorem is false.

Equivalently, with $X=e^t$, $Y=ce^r$, $a=3/2$, and $\beta=5/4$,
\[
  h(t+r)=c^{-\beta}X^\beta Y^\beta (XY-a)e^{-XY}.
\]
Positive row and column factors do not affect signs, while
\[
  (t+r)h(t+r)
  =
  c^{-\beta}X^\beta Y^\beta
  \bigl(\log X+\log(Y/c)\bigr)(XY-a)e^{-XY}.
\]
For the scalar kernel $K(X,Y)=(XY-a)e^{-XY}$ one has
\[
  \partial_X^k K(X,Y)=(-1)^kY^k(XY-a-k)e^{-XY}.
\]
After coalescing both node sets at $z=XY$, the local $n$th order reverse
sign-regularity condition is controlled by
\[
  W_{n-1}(z)=
  \det\left[\partial_z^i\{z^j(z-a-j)\}\right]_{i,j=0}^{n-1}.
\]
Reverse sign-regularity would require $W_{n-1}(z)>0$ throughout the rectangle.
For the endpoint value $a=3/2$ and $n=7$,
\[
\begin{aligned}
W_6(z)=&-394053660000+367783416000z-169746192000z^2\\
&+51438240000z^3-11430720000z^4+1959552000z^5\\
&-261273600z^6+24883200z^7,
\end{aligned}
\]
and
\[
  W_6(16/5)=-18572426057952/3125<0 .
\]
Since $16/5>\pi$, the scalar Laguerre kernel already fails the required
reverse sign-regularity inside the endpoint rectangle $X\ge1$, $Y\ge\pi$.
Thus no citation can close the endpoint proof through separate sign control of
all mixed Cauchy--Binet type words.  The next pointwise variant, pairing each
type word with its complement before integration, also fails numerically: the
order-two paired density takes negative values.  Positivity of $P$, if true,
must therefore use cancellation after the Volterra integration, equivalently an
integrated Green identity or endpoint range Hardy inequality for
\[
  Q(f)=\int_0^\infty A(r)B(r)\,\dd r,\quad
  A(r)=\sum_i c_i h(t_i+r),\quad
  B(r)=\sum_i c_i(t_i+r)h(t_i+r).
\]

The endpoint Laguerre equation gives such a Green identity, but not a local
positive square.  With $z=ce^u$,
\[
  h''(u)+(z-2)h'(u)+\left(\frac54 z+\frac{15}{16}\right)h(u)=0.
\]
Equivalently, if
\[
  \phi(u)=e^{3u/2}e^{-ce^u},
\]
then $h(u)=-e^{-u/4}\phi'(u)$.  For a finite range vector set
\[
  d_i=c_i e^{-t_i/4},\quad
  \Phi(r)=\sum_i d_i\phi(t_i+r),\quad
  \Psi(r)=\sum_i d_i(t_i+r)\phi(t_i+r).
\]
Then
\[
  A(r)=-e^{-r/4}\Phi'(r),\qquad
  B(r)=e^{-r/4}\{\Phi(r)-\Psi'(r)\},
\]
and hence
\[
  A(r)B(r)
  =
  \frac{d}{dr}\left(-\frac12e^{-r/2}\Phi(r)^2\right)
  +e^{-r/2}\left(\Phi'(r)\Psi'(r)-\frac14\Phi(r)^2\right).
\]
Thus
\[
  Q(f)=\frac12\Phi(0)^2+
  \int_0^\infty e^{-r/2}
  \left(\Phi'(r)\Psi'(r)-\frac14\Phi(r)^2\right)\dd r .
\]
The residual density can be negative pointwise, so the remaining theorem is the
constrained integrated Hardy inequality for the pair $(\Phi,\Psi)$, where
$\Psi$ is the logarithmic derivative of the same Laplace range vector.

This constraint is equivalently a parameter derivative.  Set
\[
  \Phi_\beta(r)=\sum_i d_i e^{\beta(t_i+r)}e^{-ce^{t_i+r}}.
\]
At the endpoint $\beta=3/2$ one has
\[
  \Phi=\Phi_\beta,\qquad \Psi=\partial_\beta\Phi_\beta .
\]
Thus
\[
  2Q=
  \Phi(0)^2+
  \partial_\beta\int_0^\infty e^{-r/2}|\Phi_\beta'(r)|^2\,\dd r
  -\frac12\int_0^\infty e^{-r/2}|\Phi_\beta(r)|^2\,\dd r
  \quad(\beta=3/2).
\]
Numerically, the interior form
\[
  R_{3/2}:=
  \partial_\beta\|\Phi_\beta'\|_{e^{-r/2}\dd r}^2
  -\frac12\|\Phi_\beta\|_{e^{-r/2}\dd r}^2
\]
has exactly one negative direction on the tested Galerkin spaces, is positive
to roundoff on the boundary-null subspace $\Phi(0)=0$, and the rank-one
boundary form $\Phi(0)^2$ dominates the remaining direction.  The endpoint
proof is now the boundary-null lemma plus rank-one Schur closure
\[
  \Phi(0)=0\Rightarrow R_{3/2}(\Phi)\ge0,\qquad
  R_{3/2}(\Phi)+\Phi(0)^2\ge0
\]
for truncated Laplace range functions with spectral support $y\ge1$.

The boundary condition can be normalized away.  Put
\[
  \lambda=ce^t\ge c,\qquad \tau=e^r-1,\qquad x=1+\tau.
\]
Since
\[
  \phi(t+r)=\phi(t)x^{3/2}e^{-\lambda\tau},
\]
the condition $\Phi(0)=0$ becomes $\sum_i\alpha_i=0$ after setting
$\alpha_i=d_i\phi(t_i)$.  Hence the boundary-null range is spanned by
\[
  F_\lambda(\tau)=e^{-\lambda\tau}-e^{-c\tau},\qquad
  H_\lambda(\tau)=\log(\lambda/c)e^{-\lambda\tau}.
\]
Let
\[
  q_\lambda(x)=\left(\frac32-\lambda x\right)e^{-\lambda(x-1)}.
\]
Then the boundary-null form is equivalent to positivity of the kernel
\[
\begin{aligned}
K(\lambda,\mu)=\int_0^\infty x^{3/2}\{&
\log x\,[q_\lambda-q_c][q_\mu-q_c]\\
&+\frac12[q_\lambda-q_c]\log(\mu/c)q_\mu\\
&+\frac12[q_\mu-q_c]\log(\lambda/c)q_\lambda\}\,\dd\tau ,
\end{aligned}
\]
on $\lambda,\mu\ge c$.  This is the final concrete form of the boundary-null
lemma.

Two direct proof attempts clarify the remaining structure.  First, the mixed
spectral derivative is too strong: if $K_0$ denotes the unconditioned kernel,
then positivity of $\partial_\lambda\partial_\mu K_0$ would imply positivity of
$K$ by double integration from $c$, but numerical tests show a stable small
negative mode.  Second, the split $K=L+C$ is robust.  Here
\[
  L(\lambda,\mu)=
  \int_0^\infty x^{3/2}\log x\,[q_\lambda-q_c][q_\mu-q_c]\,d\tau
\]
is a Gram kernel, while $C$ is the remaining symmetrized logarithmic
cross-kernel.  On the tested spectral windows $C$ has two negative directions,
and $L$ dominates $-C$ on that subspace with margin about $1.62$.  The Schur
complement of $L+C$ relative to this negative subspace is positive to the
resolved numerical precision.  Thus the remaining self-contained proof of
$K\ge0$ is a two-dimensional Schur comparison for this explicit split.

The cross-kernel also has an exact Sturm--Liouville form.  With
\[
  F_\lambda=e^{-\lambda\tau}-e^{-c\tau},\qquad
  H_\lambda=\log(\lambda/c)e^{-\lambda\tau},\qquad
  A=\frac32+x\frac{d}{d\tau},
\]
one has $F_\lambda(0)=0$ and
\[
  C(F,H)=\int_0^\infty x^{3/2}(AF)(AH)\,d\tau
  =
  \int_0^\infty x^{7/2}F'H'\,d\tau
  -\frac32\int_0^\infty x^{3/2}FH\,d\tau .
\]
Equivalently,
\[
  C(F,H)=
  \int_0^\infty x^{7/2}\left(F'+\frac{F}{x}\right)
  \left(H'+\frac{H}{x}\right)d\tau
  =
  \int_0^\infty x^{3/2}(xF)'(xH)'\,d\tau .
\]
The negative potential is therefore removable; the two negative squares come
from the logarithmic spectral generator on this derivative range.
The endpoint boundary term vanishes because $F(0)=0$.  Thus the negative-index
part of the proof is equivalently a two-negative-square theorem for this
paired Laplace-range Sturm form.  Simple endpoint moment restrictions remove
the dominant negative direction but do not yet remove the second one, so the
rank-two defect is not merely the first two obvious Taylor moments.

The logarithmic companion is scalar:
\[
  H(\tau)=\int_0^\infty
  \frac{e^{-cu}F(\tau)-F(\tau+u)}{u}\,du
  =
  \log\!\left(\frac{-\partial_\tau}{c}\right)F(\tau)
\]
on the truncated Laplace range.  Shiftwise positivity of the integrand is too
strong; the negative-index theorem must use cancellation across this logarithmic
integral.

There is, however, an exact rank-two source split.  Define
\[
  g_\lambda(\tau)=\bigl((1+\tau)e^{-\lambda\tau}\bigr)'
  =(1-\lambda(1+\tau))e^{-\lambda\tau},\qquad
  u_\lambda=g_\lambda-g_c,\qquad
  \ell_\lambda=\log(\lambda/c),
\]
and use the inner product
\[
  \langle a,b\rangle_E=\int_0^\infty x^{3/2}a(\tau)b(\tau)\,d\tau .
\]
Then
\[
  E(\lambda,\mu)=\langle u_\lambda,u_\mu\rangle_E
\]
is the positive Dirichlet Gram kernel, while
\[
  C(\lambda,\mu)=\frac12\{\langle u_\lambda,\ell_\mu g_\mu\rangle_E
  +\langle \ell_\lambda g_\lambda,u_\mu\rangle_E\}.
\]
Since $g_\lambda=u_\lambda+g_c$, this gives
\[
  C=P_0+R_2,
  \qquad
  P_0(\lambda,\mu)=\frac12(\ell_\lambda+\ell_\mu)E(\lambda,\mu),
\]
with
\[
  R_2(\lambda,\mu)
  =\frac12\ell_\mu\langle u_\lambda,g_c\rangle_E
  +\frac12\ell_\lambda\langle g_c,u_\mu\rangle_E .
\]
Thus $R_2$ has rank at most two.  Numerically, on
$\lambda\in[c,ce^{12}]$ with order $90$ quadrature,
\[
  \|C-P_0-R_2\|\simeq 4.1\cdot 10^{-9},\qquad
  \operatorname{rank}R_2=2,
\]
while $P_0$ itself has two negative directions.  Therefore the naive
``positive minus rank-two'' split is not the natural one.  The remaining
finite-index theorem should instead prove that $P_0$ has at most two negative
squares, that the rank-two source term does not increase the index of $C$, and
then use the already observed $L/(-C)>1$ Schur domination to prove $L+C\ge0$.

Equivalently, whitening by the positive Dirichlet Gram form gives the
generalized problem
\[
  C v=\rho E v.
\]
On the same stress window the stable generalized eigenvalues begin
\[
  -0.65118,\quad -0.10356,\quad 0.42653,\quad 0.82960,\ldots,
\]
so the expected defect is exactly two-dimensional after quotienting the
near-null closure of the Dirichlet range.

The rank-two source term has now been isolated from the remaining index
question.  Let $N_0$ denote the negative spectral subspace of $P_0$.  Then
$R_2$ does not increase the index of $C$ precisely if
\[
  C|_{N_0^\perp}\ge0
\]
with the Moore--Penrose range condition on the nullspace of this block.  The
Schur block over $N_0$ itself is allowed to be negative; those are exactly the
two retained negative directions.  Numerically, for
$\lambda\in[c,ce^{12}]$ with order $90$ quadrature,
\[
  \operatorname{ind}_-(P_0)=2,\qquad
  \operatorname{ind}_-(C)=2,\qquad
  \lambda_{\min}(C|_{N_0^\perp})\simeq -4.5\cdot10^{-10},
\]
which is at the quadrature/nullspace floor.  Simple moment quotients such as
$(1,s)$ or $(1,s^2)$ do not remove both negative directions.  Thus the
remaining analytic theorem is not a low-order endpoint moment estimate, but a
finite-index theorem for the anti-commutator kernel
\[
  P_0(s,t)=\frac12(s+t)E(s,t),\qquad s,t\ge0.
\]

This anti-commutator has an exact projected-log interpretation.  With
\[
  F_\lambda=e^{-\lambda\tau}-e^{-c\tau},\qquad
  B=(1+\tau)\partial_\tau+1,\qquad
  u_\lambda=BF_\lambda,
\]
and with $S_0$ the logarithmic operator projected onto the boundary-null span,
\[
  S_0F_\lambda=\log(\lambda/c)F_\lambda,
\]
one has
\[
  P_0(\lambda,\mu)
  =\frac12\langle BF_\lambda,BS_0F_\mu\rangle
   +\frac12\langle BS_0F_\lambda,BF_\mu\rangle.
\]
Replacing $S_0$ by the unprojected companion
$SF_\lambda=\log(\lambda/c)e^{-\lambda\tau}$ gives $C=P_0+R_2$ with
$\operatorname{rank}R_2\le2$.  The tempting comparison with the unconditioned
anti-commutator
\[
  A_g(\lambda,\mu)=\frac12(\log(\lambda/c)+\log(\mu/c))
  \langle g_\lambda,g_\mu\rangle
\]
does not close the proof: $P_0-A_g$ is a rank-four boundary correction in the
tested finite sections.  Thus the two-negative-square theorem must be proved
directly for the projected-log anti-commutator $P_0$.

The direct quotient certificate suggested by the finite sections is now
explicit.  For boundary-null functions the two first-order completions
\[
  x\partial_\tau+1,\qquad x\partial_\tau+\frac32
\]
give the same Dirichlet form, since the cross term integrates by parts and
the endpoint value $F(0)$ vanishes.  Let
\[
  \beta(\lambda)=\langle u_\lambda,g_c\rangle,\qquad
  \zeta(\lambda)=1-\frac{c}{\lambda}.
\]
Here
\[
  \zeta(\lambda)=-c\int_0^\infty F_\lambda(\tau)\,d\tau.
\]
Numerically, imposing
\[
  \sum_i a_i\beta(\lambda_i)=0,\qquad
  \sum_i a_i\zeta(\lambda_i)=0
\]
makes $P_0$ nonnegative to the quadrature/nullspace floor.  On
$\lambda\in[c,ce^{10}]$ with order $100$ quadrature the quotient minimum is
about $-10^{-10}$ while the two unrestricted negative eigenvalues are
$-4.18\cdot10^{-3}$ and $-7.10\cdot10^{-8}$.  Thus the remaining proof of
Lemma A is the concrete quotient theorem: $P_0\ge0$ on the simultaneous
nullspace of $\beta$ and $\zeta$.

Equivalently, it suffices to prove the finite-rank lower bound
\[
  P_0(\alpha)+14\,\beta(\alpha)^2+14\,\zeta(\alpha)^2\ge0 .
\]
The numerical threshold for the scalar coefficient is approximately $13.04$
on the stable windows.  The Stieltjes representation of
$\log(\lambda/c)$ does not yield a layerwise proof: the individual resolvent
layers become negative on the same quotient for large layer parameter.  The
remaining proof therefore has to use the integrated endpoint Hardy/Green
structure.

The same two rows also certify the finite index of the cross-kernel $C$:
finite sections satisfy
\[
  C+14(\beta^2+\zeta^2)\ge0
\]
to the same numerical floor, with scalar threshold about $11.78$ on the stable
windows.  This gives a direct route to $\operatorname{ind}_-(C)\le2$.
The rows $(\beta,\zeta)$ are not, however, the right final Schur coordinates
for $K=L+C$: the quotient block has a large near-null space and the
Moore--Penrose Schur form in this arbitrary row basis is ill-conditioned.  The
stable final closure is still the spectral one relative to the actual negative
subspace of $C$, where numerically $L/(-C)>1$ and the Schur complement of
$L+C$ is positive.

The current bottleneck has an exact closed-form version.  If
\[
  G(\lambda,\mu)=\int_1^\infty x^{3/2}g_\lambda(x)g_\mu(x)\,dx,\qquad
  a=\lambda+\mu,
\]
and
\[
  J_p(a)=\int_1^\infty x^p e^{-a(x-1)}\,dx,
\]
then
\[
  G(\lambda,\mu)
  =J_{3/2}(a)-aJ_{5/2}(a)+\lambda\mu J_{7/2}(a),
\]
with
\[
  J_{1/2}(a)=\frac1a+
  \frac{\sqrt\pi\,e^a\operatorname{erfc}(\sqrt a)}{2a^{3/2}},
  \qquad
  J_p(a)=\frac1a+\frac{p}{a}J_{p-1}(a).
\]
Consequently $E,\beta,\zeta,P_0$, and $C$ are explicit half-integer
incomplete-gamma kernels.  Rebuilding the finite sections from these formulas,
without the inner $x$-quadrature, reproduces the same two negative eigenvalues
and the same scalar thresholds: about $13.04$ for $P_0$ and $11.78$ for $C$.

The same formula splits $G$ into a rational part plus one erfc term:
\[
\begin{aligned}
  G_{\rm rat}
  &=-1+\frac{\lambda\mu-3/2}{a}
    +\frac{(7/2)\lambda\mu-9/4}{a^2}
    +\frac{(35/4)\lambda\mu}{a^3}
    +\frac{(105/8)\lambda\mu}{a^4},\\
  G_{\rm erfc}
  &=
  \frac{\sqrt\pi\,(105\lambda\mu-18a^2)e^a
  \operatorname{erfc}(\sqrt a)}{16a^{9/2}}.
\end{aligned}
\]
This split is useful diagnostically but not by itself a proof: in the tested
sections the corrected rational piece is nonnegative to the numerical floor,
whereas the corrected erfc piece still has two negative directions.  Thus the
finite-rank Hardy theorem must keep the rational and erfc boundary rows coupled
inside the full kernel.  Assigning the $\beta_{\rm rat}\beta_{\rm erfc}$ cross
term to the erfc block makes that block more negative in the same finite
sections, so this is not merely a matter of choosing a better componentwise
allocation of the finite-rank correction.

There is also a useful endpoint normalization.  Since $u_\lambda=g_\lambda-g_c$
vanishes at $\lambda=c$, divide the kernel by $(\lambda-c)(\mu-c)$.  This
congruence preserves inertia on $(c,\infty)$ and turns the two boundary rows
into
\[
  e(\lambda)=\frac1\lambda,\qquad
  d(\lambda)=\frac{\beta(\lambda)}{\lambda-c}.
\]
Thus the finite-rank Hardy target is equivalently
\[
  \widehat P_0+14(ee^T+dd^T)\ge0,\qquad
  \widehat P_0(\lambda,\mu)
  =\frac{P_0(\lambda,\mu)}{(\lambda-c)(\mu-c)}.
\]
The analogous statement for $\widehat C$ has the same two rows.  Numerically
this normalized certificate has the same two negative directions before the
row correction and is nonnegative after the row correction to the spectral
floor.

This normalization also rules out two plausible shortcuts.  The row
$\beta/\zeta$ is not monotone near the endpoint, so the pair $(\beta,\zeta)$
does not reduce to a Chebyshev-system quotient after scalar normalization.
The row $d(\lambda)$ is positive, decreasing, and convex in
$s=\log(\lambda/c)$ on the tested windows, but finite-difference tests fail at
the third complete-monotonicity sign.  The remaining theorem is therefore the
coupled conditional positivity of $\widehat P_0$ relative to the two rows
$1/\lambda$ and $\beta/(\lambda-c)$, not a standalone monotonicity theorem for
the rows.

The normalized theorem also has a direct range formulation.  Let
\[
  F_\lambda(\tau)=\frac{e^{-\lambda\tau}-e^{-c\tau}}{\lambda-c},
  \qquad B=x\partial_\tau+1,\qquad x=1+\tau .
\]
Then $v_\lambda=BF_\lambda=(g_\lambda-g_c)/(\lambda-c)$ and
\[
  \widehat P_0(\alpha)
  =\langle BF,BH\rangle_{x^{3/2}d\tau},\qquad
  H=\sum_i \alpha_i\log(\lambda_i/c)F_{\lambda_i}.
\]
The two rows are source moments:
\[
  e(\lambda)=-c\int_0^\infty F_\lambda(\tau)\,d\tau=\frac1\lambda,
\]
and, using $B^*h=-xh'-3h/2$,
\[
  d(\lambda)=
  \int_0^\infty x^{3/2}F_\lambda(\tau)
  \left[-c^2x^2+\frac{7c}{2}x-\frac32\right]e^{-c(x-1)}\,d\tau .
\]
Thus the row-null theorem is a concrete divided-difference Laplace range
inequality for $\langle BF,BH\rangle$ under two source moment constraints.

The mixed-derivative route remains false after normalization: finite sections
of $\partial_s\partial_t\widehat P_0$ still have stable negative spectrum.
The promising formulation is instead a bordered determinant theorem.  For
$0<s_1<\cdots<s_n$, define the bordered matrix
\[
  \mathcal B_n=
  \begin{pmatrix}
  0&0&e(s_j)\\
  0&0&d(s_j)\\
  e(s_i)&d(s_i)&\widehat P_0(s_i,s_j)
  \end{pmatrix}.
\]
Numerically, all tested consecutive and random bordered determinants are
positive.  High-precision checks at the first double-precision failure orders
give positive determinants as small as $10^{-86}$ and $10^{-101}$.  Moreover
$\widehat C-\widehat P_0$ is a symmetric rank term containing the included row
$d$, so the same bordered determinants control both $P_0$ and $C$.
The next finite-index lemma is therefore the bordered total-positivity of
$\mathcal B_n$ for all $n$.

Subsequent high-precision confluent tests show that this bordered
total-positivity statement is still too strong.  At local nodes
$s_i=s_0+ih$, the first genuine failure occurs at order $n=8$; for example
near $s_0=0.5$, $h=0.05$, the row-null form for $\widehat P_0$ has a negative
eigenvalue of size about $4.35\cdot10^{-20}$.  The associated coefficients are
an alternating high-order finite-difference pattern, explaining why ordinary
Galerkin sections saw only the two macroscopic negative modes.

This corrects the proof program: $\widehat P_0$ and $\widehat C$ do not have
only two negative squares as point kernels.  There is a tiny confluent negative
tail.  The full kernel remains viable because the positive $L$ term dominates
that tail strongly; on the same local witnesses the observed ratio
$L/(-\widehat P_0)$ ranges from about $21$ to $58$.  Thus the exact theorem to
prove is not finite index for $C$ alone, but direct positivity of
\[
  \widehat K=\widehat L+\widehat C,
\]
with $\widehat L$ controlling both the macroscopic and confluent negative
directions.

High-precision local matrices for $\widehat K$ support this corrected target.
On $n=8$ close-node tests near $s_0=0.45$ and $s_0=0.50$, $\widehat C$ has
negative local eigenvalues of size about $10^{-5}$ and $10^{-13}$, while
$\widehat K$ has nonnegative spectrum down to the confluent numerical floor
near $10^{-19}$.  The remaining proof should therefore split the problem into
a macroscopic Schur comparison and a local Hardy/Wronskian domination of the
confluent negative tail by $\widehat L$.

There is a cleaner transfer.  Let $B=x\partial_\tau+1$ and
$A=B+1/2$.  For the normalized boundary-null range
\[
  F_\lambda=\frac{e^{-\lambda\tau}-e^{-c\tau}}{\lambda-c},\qquad
  H_\lambda=\frac{\log(\lambda/c)e^{-\lambda\tau}}{\lambda-c},
\]
write $U=BF$ and $W=BH$.  The $B$-branch kernel is
\[
  K_B=\int x^{3/2}\{\log x\,U^2+UW\}\,d\tau ,
\]
whereas the older endpoint $A$-branch kernel is
\[
  K_A=\int x^{3/2}\{\log x\,(AF)^2+(AF)(AH)\}\,d\tau .
\]
Since $F(0)=0$,
\[
  \langle BF,G\rangle+\langle F,BG\rangle=-\frac12\langle F,G\rangle.
\]
Applying this once with $G=H$ and once with $G=(\log x)F$, using
$B((\log x)F)=(\log x)BF+F$, gives the exact identity
\[
  K_B-K_A=\frac12\int_0^\infty x^{3/2}F(\tau)^2\,d\tau .
\]
Thus the corrected $B$-branch endpoint theorem follows from the older
$A$-branch endpoint theorem by adding a positive Gram term.  Numerically this
identity holds to $10^{-15}$ in endpoint quadrature and to about $10^{-61}$ in
the close-node high-precision test.

There is one further correction.  The older $A$-branch theorem is too strong if
read as point-kernel positivity.  A closed-form high-precision evaluation of
the endpoint Green kernel $P(t,s)$, using incomplete-gamma log moments rather
than quadrature, finds a genuine local confluent negative mode at order
$n=10$:
\[
  \lambda_{\min}\simeq -2.38\cdot10^{-27}
\]
for nodes $t_i=0.5+0.05i$.  The boundary condition $\Phi(0)=0$ does not remove
this mode.  Thus the total-positivity or positive point-kernel formulation of
the endpoint Hardy theorem is false.

The transfer identity suggested a sharper target because the corrected
$B$-branch kernel remained positive in the first close-node local test.  The
tempting endpoint problem was therefore
\[
  K_B=K_A+\frac12\int_0^\infty x^{3/2}F(\tau)^2\,d\tau\ge0,
\]
not $K_A\ge0$ alone.

This point-kernel target is also false.  The close-node test at spacing
$h=0.05$ was pre-asymptotic.  A formal Taylor-jet computation at the integrand
level gives the confluent matrices
\[
  \left[\frac{\partial_s^i\partial_t^jK(s_0,t_0)}{i!\,j!}\right]_{i,j=0}^{n-1},
  \qquad t_0=s_0,
\]
which must be positive for any analytic positive kernel.  At $s_0=0.5$, the
$B$-branch jet is positive through $n=8$ but has
\[
  \lambda_{\min}(n=9)\simeq -6.05\cdot10^{-6}.
\]
The obstruction persists nearby; at $n=9$ the minima are about
$-1.72\cdot10^{-5}$ for $s_0=0.45$ and $-2.22\cdot10^{-7}$ for $s_0=0.55$.
Thus the exact identity $K_B=K_A+\frac12F{\rm Gram}$ is useful but insufficient:
the positive divided-difference Gram term does not dominate all high-order
endpoint jets.  The next problem is to determine whether the full three-mode
kernel contributes an additional positive jet term, whether a genuine
range/quotient condition removes this endpoint jet, or whether endpoint
positivity only appears after pairing with the reflected Schur block.

The first alternative is supported by the direct finite-core jet test.  For the
full reduced Volterra kernel
\[
 K_{\rm red}(s,t)=
 \int_0^\infty
   \Bigl[u+\frac{s+t}{2}\Bigr]
   \cosh\!\Bigl(\omega\bigl[u+\frac{s+t}{2}\bigr]\Bigr)
   \frac{\Psi(s+u)}{\Psi(s)}
   \frac{\Psi(t+u)}{\Psi(t)}\,du,
\]
with $\Psi(v)=\widetilde\Phi_3(v/2)$, the same Taylor-jet computation at
$s_0=0.5$ gives, for $\omega=0.49$,
\[
  \lambda_{\min}(n=9)\simeq 7.03\cdot10^{-6},
  \qquad
  \lambda_{\min}(n=10)\simeq 1.95\cdot10^{-6}.
\]
The ninth-order test also remains positive near the obstruction:
$7.42\cdot10^{-6}$ at $s_0=0.45$ and $6.67\cdot10^{-6}$ at $s_0=0.55$.
Therefore the endpoint defect is not fatal for the finite-core Weyl kernel; it
is removed by the full three-mode Volterra structure.  The next analytic lemma
should identify this finite-core correction and prove that it is a positive
Gram or Schur-positive perturbation of the endpoint model.

More precisely, the correction
\[
  \Delta=K_{\rm red}(\widetilde\Phi_3)-K_{\rm endpoint,B}
\]
is not positive as a kernel.  In the ninth-order jet at $s_0=0.5$,
$\omega=0.49$, it has negative eigenvalues.  However, if $e$ denotes the unique
negative vector of the endpoint jet, then
\[
  \frac{\langle e,\Delta e\rangle}{-\langle e,K_{\rm endpoint,B}e\rangle}
  \simeq 11.14.
\]
Moreover, the full finite-core jet is positive on the endpoint-positive
complement, with minimum about $2.04\cdot10^{-5}$, and its Schur complement
relative to $e$ is about $1.06\cdot10^{-5}$.  Thus the correct local statement
is a Schur comparison, not $\Delta\ge0$.

If $\Psi=\Psi_1+Z$ and
\[
 A_s(u)=\frac{\Psi(s+u)}{\Psi(s)},\qquad
 A_s^1(u)=\frac{\Psi_1(s+u)}{\Psi_1(s)},\qquad
 \Theta_s(u)=
 \frac{1+Z(s+u)/\Psi_1(s+u)}{1+Z(s)/\Psi_1(s)},
\]
then the higher-mode finite-core correction is exactly
\[
 K_{\rm red}(\Psi)-K_{\rm red}(\Psi_1)
 =
 \int_0^\infty W(s,t,u)A_s^1(u)A_t^1(u)
      \{\Theta_s(u)\Theta_t(u)-1\}\,du,
\]
where
\[
 W(s,t,u)=
 \Bigl[u+\frac{s+t}{2}\Bigr]
 \cosh\!\Bigl(\omega\bigl[u+\frac{s+t}{2}\bigr]\Bigr).
\]
The global proof should promote this local Schur comparison into a Volterra
moment inequality on the endpoint-defect space and its complement.

A global Galerkin version supports the same interpretation.  On smooth
quadrature spaces over $[0,2]$, the endpoint $B$-model has only near-null
negative modes, for instance
\[
  -1.80\cdot10^{-13},\qquad -1.65\cdot10^{-18},
\]
while $K_{\rm red}(\widetilde\Phi_3)$ remains nonnegative to the Galerkin
floor.  Relative to these endpoint-negative modes the correction ratios are
about $12.97$ and $1.09$, and the Schur complement of the full finite-core
form is positive to the same floor.  On $[0,4]$ the endpoint model has no
negative smooth Galerkin mode above $10^{-18}$ and the full finite-core form is
positive.

Thus the endpoint obstruction should be treated as a distributional
Taylor-jet defect, not as a robust negative eigenspace in smooth $L^2[0,L]$.
The global theorem should split the test space as
\[
  \mathcal H_{\rm test}=J_{\rm endpoint}\oplus M_{\rm smooth}
\]
and prove the singular Schur/moment statement:
\[
  Q_\Psi\ge0\ \text{on }M_{\rm smooth},\qquad
  Q_\Psi(J_{\rm endpoint},M_{\rm smooth})\in{\rm Ran}(Q_\Psi|_{M_{\rm smooth}}),
\]
with nonnegative Moore--Penrose Schur form on $J_{\rm endpoint}$.  Here
\[
 Q_\Psi(f)=
 \frac12\sum_{\sigma=\pm1}\int_0^\infty m_\sigma(u)n_\sigma(u)\,du
\]
is the Volterra moment form.

The endpoint-defect space can now be made concrete.  Let
\[
 J_9(s_0)=
 \left[
 \frac{\partial_s^i\partial_t^jK_{\rm endpoint,B}(s_0,s_0)}{i!\,j!}
 \right]_{i,j=0}^{8}.
\]
Its lowest eigendirection \(e(s_0)\) defines the Taylor-normalized functional
\[
  \Lambda_{s_0}(f)=\sum_{k=0}^8 e_k(s_0)\frac{f^{(k)}(s_0)}{k!}.
\]
Numerically this lowest eigendirection is simple and smooth through the local
window, with spectral gap at least about \(1.12\cdot10^{-5}\).  The lowest
eigenvalue crosses zero at
\[
  s_* \simeq 0.5530736.
\]
It is negative for \(s_0<s_*\) and positive for \(s_0>s_*\), with no second
negative window detected up to \(s_0=1\).  Thus
\[
  J_{\rm endpoint}
  =
  \overline{\operatorname{span}}\{\Lambda_{s_0}:0\le s_0<s_*\},
\]
and the smooth complement should be imposed by
\[
  M_{\rm smooth}=\{f:\Lambda_{s_0}(f)=0\text{ for }0\le s_0<s_*\},
\]
or equivalently by quotienting this endpoint-jet range in the singular
Schur/GNS construction.

The proposed complement has also been tested directly.  In a shifted Legendre
Galerkin space on $[0,8]$, imposing sampled constraints
\[
  \Lambda_a(f)=0,\qquad a\in[0.05,0.54],
\]
removes the endpoint obstruction.  For $\widetilde\Phi_3$, $\omega=0.49$,
basis order $24$, and eight active constraints, the constrained minimum is
only $-2.45\cdot10^{-17}$, while the first clearly positive eigenvalues are
about $10^{-15},10^{-12},10^{-9},10^{-7}$.  With twelve constraints and basis
order $30$ the minimum remains at the numerical floor.  The same test at
$\omega=0$ is also floor-level.  As a control, the one-mode core has a visible
local endpoint defect, but after imposing the same sampled endpoint constraints
its smooth complement is again positive to the floor.

Thus the next theorem is the continuum constrained moment inequality:
\[
  \Lambda_a(f)=0\ (0\le a<s_*)\quad\Longrightarrow\quad
  Q_\Psi(f)\ge0.
\]
The finite-core correction is needed for the singular Schur block over
$J_{\rm endpoint}$; the smooth complement should be controlled by this
endpoint-Hardy/Volterra moment theorem.

There is one important technical correction to the most naive formulation of
this theorem.  The condition $\Lambda_a(f)=0$ should not be treated globally as
a nonsingular eighth-order ODE obtained by solving for $f^{(8)}(a)$.  With the
stable sign convention $e_0(a)<0$, the coefficient $e_8(a)$ changes sign in the
active interval; numerically,
\[
  e_8(a)=0\quad\text{at}\quad a\simeq 0.1646167.
\]
At the same time the lowest eigendirection remains simple, and the lower
coefficients remain nonzero in the tested window.  For example, on
$[0.02,0.545]$ the diagnostic scan gives
\[
  \min |e_0|\simeq 6.99\cdot 10^{-3},\qquad
  \min |e_7|\simeq 3.67\cdot 10^{-2}.
\]
Thus the correct coordinate-free object is the smooth hyperplane field in the
nine-dimensional jet bundle,
\[
  E_a=\left\{j_a^8 f:
  \sum_{k=0}^8 e_k(a){f^{(k)}(a)\over k!}=0\right\},
  \qquad 0\le a<s_*.
\]
On subintervals where $e_8\ne0$ this may be written as an eighth-order
equation, but the proof should be phrased in the jet-hyperplane form.  The
continuum theorem is therefore the positivity of $Q_\Psi$ on sections whose
eighth jets lie in $E_a$ throughout the active endpoint window.

The corresponding quotient certificate is simple.  Let
\[
  (Rf)(a)=\Lambda_a(f),\qquad 0\le a<s_*.
\]
It is enough to prove a factorization
\[
  Q_\Psi(f)=\|Gf\|_{\mathcal H}^2-\|SRf\|_{\mathcal X}^2
\]
with $G$ a Volterra Gram map and $S$ bounded on the completed trace range.
Then $Rf=0$ implies $Q_\Psi(f)=\|Gf\|_{\mathcal H}^2\ge0$.  Equivalently, the
negative part of the Volterra form must factor through the closed span of the
endpoint jet functionals $\Lambda_a$.  This is the coordinate-free replacement
for the failed global ODE chart.

The finite-dimensional certificate has a clean Schur form.  In a Galerkin
space write \(V=\ker R\oplus U\), where \(U\) is the row space of \(R\), and
block
\[
  K=\begin{pmatrix} A&B\\ B^*&C\end{pmatrix}.
\]
Then
\[
  K=P-R^*DR,\qquad P\ge0,\quad D\ge0,
\]
exists provided
\[
  A\ge0,\qquad {\rm ran}\,B\subset{\rm ran}\,A.
\]
Indeed, choose \(M\ge B^*A^+B-C\) with \(M\ge0\) and set
\[
  P=\begin{pmatrix} A&B\\ B^*&C+M\end{pmatrix}.
\]
The Schur complement gives \(P\ge0\), and because \(R|_U\) is injective onto
its trace range, \(M\) can be represented as \(U^*R^*DRU\) with \(D\ge0\).

Numerically, the endpoint \(B\)-model on \([0,2]\), basis order \(16\), and
eight sampled active traces gives
\[
  \lambda_{\min}(K)\simeq -1.80\cdot10^{-13},\qquad
  \lambda_{\min}(K|_{\ker R})\simeq 1.85\cdot10^{-10},
\]
with zero range residual at the displayed precision.  The constructed repair
has
\[
  \lambda_{\min}(K+R^*DR)\simeq 9.8\cdot10^{-19}.
\]
Thus the sampled endpoint negative direction factors through \(R\), exactly as
the quotient theorem predicts.

The continuum form of the same statement uses the Moore--Penrose range
condition rather than literal range inclusion.  Let \(V\) be the Hilbert form
space for the truncated problem, for instance \(H^8(0,L)\), and define
\[
  (Rf)(a)=\Lambda_a(f)=
  \sum_{k=0}^8 e_k(a){f^{(k)}(a)\over k!},\qquad 0\le a<s_*.
\]
Once the eigenline \(e(a)\) is simple and bounded on the active interval,
\[
  \|Rf\|_{L^2(0,s_*)}
  \le C\sum_{k=0}^8\|f^{(k)}\|_{L^2(0,s_*)}
  \le C\|f\|_{H^8(0,L)}.
\]
Thus \(R\) is closed and \(N=\ker R\) is closed.  Write
\[
  V=N\oplus U,\qquad U=N^\perp,
\]
and give the trace range \(X_R=R(U)\) the transported norm
\[
  \|Ru\|_{X_R}:=\|u\|_V.
\]
Then \(R_U:U\to X_R\) is unitary by construction, so no closed-range
assumption in \(L^2(0,s_*)\) is needed.

Let \(q\) be the Hermitian form for \(Q_\Psi\), blocked as
\[
  q(n+u,m+v)
  =a(n,m)+b(n,v)+\overline{b(m,u)}+c(u,v).
\]
The continuum Schur hypotheses are
\[
  a\ge0
\]
on \(N\), and the Douglas range condition
\[
  b(n,u)=\langle A^{1/2}n,\Gamma u\rangle_{\mathcal H_A}
\]
for a bounded operator \(\Gamma:U\to\mathcal H_A\), where \(\mathcal H_A\) is
the completion of \(N/\ker A\) in the norm \(a(n,n)^{1/2}\).  This is the
closed-form replacement for \({\rm ran}\,B\subset{\rm ran}\,A\).

Set
\[
  T=\Gamma^*\Gamma-C,\qquad M=T_+.
\]
Define \(S:X_R\to U\) by
\[
  S(Ru)=M^{1/2}u.
\]
Since \(R_U\) is unitary, \(S\) is bounded.  Now put
\[
  p(n+u)
  =
  \|A^{1/2}n+\Gamma u\|_{\mathcal H_A}^2
  +\langle (C+M-\Gamma^*\Gamma)u,u\rangle .
\]
Because \(M\ge \Gamma^*\Gamma-C\), \(p\ge0\).  Expanding gives
\[
  q(n+u)=p(n+u)-\|M^{1/2}u\|^2
        =p(n+u)-\|SR(n+u)\|^2.
\]
Thus
\[
  Q_\Psi(f)=\|Gf\|^2-\|SRf\|^2,\qquad \|Gf\|^2:=p(f).
\]
In particular \(Rf=0\Rightarrow Q_\Psi(f)\ge0\).

Conversely, any such quotient factorization forces \(a\ge0\) by restriction
to \(\ker R\), and the positive block form \(p\) gives the same
Moore--Penrose/Douglas range condition for \(b\).  Hence the continuum
Schur/range problem is exactly the positivity of \(Q_\Psi\) on \(\ker R\)
together with this bounded \(A^{1/2}\)-factorization of the cross form.

The finite diagnostics now measure the corresponding Douglas constant directly
through
\[
  \Gamma^*\Gamma=B^*A^+B.
\]
For the endpoint \(B\)-model on \([0,2]\), basis order \(16\), and eight
active traces, the finite quotient has zero normalized range residual and
\[
  \|\Gamma\|^2\simeq 6.90\cdot10^{-8}.
\]
For the full \(\widetilde\Phi_3\) reduced kernel on \([0,8]\), basis order
\(12\), and six active traces, the corresponding value is
\[
  \|\Gamma\|^2\simeq 5.89\cdot10^{-2},
\]
again with zero normalized range residual in the finite split.  Thus the next
analytic lemma is the uniform Douglas estimate
\[
  |b(n,u)|^2\le C_D\,a(n,n)\,\|u\|_U^2,
  \qquad n\in\ker R,\ u\in U,
\]
with \(C_D\) stable under the Galerkin and truncation limits.

An endpoint stress refinement gives no sign of blow-up.  For the endpoint
\(B\)-model on \([0,2]\), \(\omega=0.49\), and
\(\#\{\Lambda_a\}\simeq{\rm basis}/2\), the values
\[
\begin{array}{c|c|c|c}
{\rm basis} & \lambda_{\min}(K) & \lambda_{\min}(K|_{\ker R})
& \|\Gamma\|^2 \\
\hline
10 &  2.30\cdot10^{-12} & 2.65\cdot10^{-6}  & 1.36\cdot10^{-6}\\
12 & -1.72\cdot10^{-13} & 7.96\cdot10^{-8}  & 5.21\cdot10^{-7}\\
14 & -1.80\cdot10^{-13} & 6.84\cdot10^{-9}  & 3.09\cdot10^{-7}\\
16 & -1.80\cdot10^{-13} & 1.85\cdot10^{-10} & 6.90\cdot10^{-8}
\end{array}
\]
show a decreasing finite Douglas constant, with zero normalized range residual
in each case.  This is evidence for the uniform estimate, not a substitute for
it.

The top eigenvector of \(B^*A^+B\) at basis order \(16\) has
\[
  \lambda_{\max}(B^*A^+B)\simeq 6.90\cdot10^{-8},
  \qquad \|\Gamma\|\simeq 2.63\cdot10^{-4}.
\]
Its physical-space representative is oscillatory on \([0,2]\), while the
paired witness \(A^+Bu\in\ker R\) is concentrated toward the right edge of the
local endpoint interval.  This points to an endpoint Hardy/trace mechanism for
the Douglas estimate.

At low order the extremizing shape itself is not yet stable.  A top-vector scan
over basis orders \(10,12,14\) gives consecutive normalized correlations about
\(0.46\) and \(0.44\), while
\[
  \lambda_{\max}(B^*A^+B)
  =
  1.36\cdot10^{-6},\ 5.21\cdot10^{-7},\ 3.09\cdot10^{-7}.
\]
Thus the constants move in the right direction, but the proof should target
the operator inequality directly rather than a presumed limiting extremizer.

The Riesz witness profile gives the concrete Hardy target.  For the basis
\(16\) top vector, \(A^+Bu\) has about \(59\%\) of its sampled \(L^2\) mass in
the last \(0.1\) of \([0,2]\), and about \(86\%\) in the last \(0.4\).  Hence
the expected analytic input is a right-end Hardy inequality for the source
functional \(L_u(n)=b(n,u)\) on \(n\in\ker R\):
\[
  |L_u(n)|^2\le C\,E_{\rm Hardy,right}(n)\,\|u\|_U^2,
\]
with the moving jet constraints \(\Lambda_a(n)=0\) removing the singular
endpoint defect.

The cross source \(h_u=NBu\) is broader than the witness: in the same test it
peaks near \(s\simeq1.48\), has about \(6.7\%\) of its mass in the last \(0.1\),
and about \(64\%\) in the last \(0.8\).  Thus the boundary layer is produced by
the Green/Riesz solve \(Aw=h_u\) on \(\ker R\).  The Hardy estimate should
control this endpoint Green response.

The spectral profile of this Riesz solve is useful because it rules out a
rank-one extremizer proof.  Diagonalizing the positive block \(A\) on the
finite \(\ker R\) gives, for the same basis \(16\) endpoint stress test,
\[
\begin{array}{c|c|c|c}
{\rm mode} & \lambda_j & {\rm energy\ share} & {\rm witness\ share}\\
\hline
3 & 3.69\cdot10^{-6} & 34.3\% & 0.12\%\\
6 & 9.54\cdot10^{-3} & 24.1\% & 0.00\%\\
2 & 2.28\cdot10^{-7} & 20.7\% & 1.16\%\\
5 & 7.50\cdot10^{-4} &  7.7\% & 0.00\%\\
1 & 1.09\cdot10^{-8} &  7.3\% & 8.59\%
\end{array}
\]
while the witness norm itself is dominated by the two smallest modes:
mode \(0\) carries about \(90\%\) of the sampled witness norm but only about
\(1.3\%\) of the Douglas energy.  Thus the analytic estimate should be a
multi-mode endpoint Green inequality of the schematic form
\[
  \sum_j {|\langle h_u,e_j\rangle|^2\over \lambda_j}
  \le C_D \|u\|_U^2,
\]
for the constrained eigenfunctions \(e_j\in\ker R\), rather than a proof based
on a single limiting profile.

A refinement scan makes the same point more sharply.  With
\(\#\{\Lambda_a\}\simeq{\rm basis}/2\), the share of Douglas energy in the two
smallest \(A\)-eigenmodes, the share of ordinary witness norm in those modes,
and the effective energy rank are
\[
\begin{array}{c|c|c|c}
{\rm basis} & E_{\rm low2} & W_{\rm low2} & {\rm effective\ rank}\\
\hline
10 & 91.9\% & 100.0\% & 1.74\\
12 & 92.0\% & 100.0\% & 2.07\\
14 &  4.2\% &  96.7\% & 2.89\\
16 &  8.6\% &  98.7\% & 4.32
\end{array}
\]
Thus the visible endpoint boundary layer persists in the low modes, but the
energy estimate migrates into a moving mid-mode spectral window.  The analytic
lemma should therefore be a uniform windowed estimate
\[
  \sum_{j\in I}{|\langle h_u,e_j\rangle|^2\over \lambda_j}
  \le C_I\|u\|_U^2,\qquad \sup_I C_I<\infty,
\]
with the bound supplied by the moving endpoint trace constraints.

The direct window scan confirms where the estimate is stressed.  For a window
\(I\), set
\[
  C_I=\lambda_{\max}\bigl(B^*P_I A_I^{-1}P_I B\bigr).
\]
The full window gives the Douglas constant.  In the endpoint stress sections
the largest single mode shifts from mode \(0\) to modes \(2,3\), while broad
proper windows still carry almost all of the full constant:
\[
\begin{array}{c|c|c|c}
{\rm basis} & {\rm best\ single} & {\rm single/full}
& {\rm large\ proper\ windows}\\
\hline
10 & 0 & 73.4\% & [0,4]\ {\rm carries}\ 99.5\%\\
12 & 1 & 63.3\% & [0,5]\ {\rm carries}\ 99.99\%\\
14 & 2 & 46.6\% & [1,6]\ {\rm carries}\ 97.3\%\\
16 & 3 & 34.7\% & [2,8]\ {\rm carries}\ 91.4\%
\end{array}
\]
At basis \(16\), the tails beginning at modes \(3,4,6,7\) carry approximately
\(70.9\%,36.4\%,28.0\%,3.9\%\) of the full constant.  This suggests a proof
by a finite low/mid spectral Schur comparison, followed by a genuine
high-frequency Hardy tail estimate using the smoothing of the source range
\(B(U)\).

The three-part finite certificate makes this division explicit.  In the same
endpoint stress sections, the smallest cutoff \(M\) for which
\(C_{\rm tail}(M)\le \alpha C_{\rm full}\) is
\[
\begin{array}{c|cccc}
{\rm basis} & \alpha=.5 & \alpha=.25 & \alpha=.1 & \alpha=.05\\
\hline
10 & 1 & 2 & 2 & 3\\
12 & 2 & 2 & 2 & 3\\
14 & 3 & 4 & 5 & 5\\
16 & 4 & 7 & 7 & 7
\end{array}
\]
Thus by basis \(16\) the finite block must include modes \(0,\ldots,6\)
before the tail is below \(5\%\).  At that cutoff
\[
  C_{\rm head}(7)\simeq .961\,C_{\rm full},\qquad
  C_{\rm tail}(7)\simeq .0386\,C_{\rm full}.
\]
The endpoint trace part is also visible numerically.  The sampled trace
constraints annihilate the sample rows to roundoff, but the first two
constrained modes still have dense moving-trace leakage of size about
\(8.4\cdot10^{-2}\) and \(1.0\cdot10^{-1}\) in the basis \(16\) test.  Their
source fractions are only \(1.5\%\) and \(7.4\%\), however; the dominant source
modes are the transition modes \(3\) and \(6\), with source fractions about
\(34.7\%\) and \(24.1\%\).  This is consistent with the intended theorem:
the closed continuum trace condition should remove the low singular defect,
while a finite Schur block controls the transition modes and a Hardy tail
controls the rest.

Finally, a fixed-basis trace-refinement test distinguishes a real endpoint
obstruction from a sampling artifact.  Holding the basis at \(16\) and
increasing the sampled trace rows from \(5\) to \(12\), while measuring leakage
against a denser \(32\)-point moving trace grid, gives
\[
\begin{array}{c|c|c|c}
\#\Lambda_a & \dim\ker R & \|\Gamma\|^2 & {\rm dense\ trace\ leakage}\\
\hline
5  & 11 & 9.41\cdot10^{-12} & 4.01\cdot10^{1}\\
6  & 10 & 1.55\cdot10^{-10} & 3.69\cdot10^{0}\\
7  &  9 & 1.36\cdot10^{-9}  & 1.48\cdot10^{-1}\\
8  &  8 & 6.90\cdot10^{-8}  & 9.58\cdot10^{-2}\\
9  &  7 & 8.94\cdot10^{-8}  & 1.67\cdot10^{-2}\\
10 &  6 & 1.02\cdot10^{-6}  & 1.28\cdot10^{-3}\\
11 &  5 & 4.84\cdot10^{-6}  & 7.89\cdot10^{-5}\\
12 &  4 & 5.03\cdot10^{-5}  & 3.02\cdot10^{-6}
\end{array}
\]
Thus sparse sampled traces can create spurious low-mode leakage, but the
leakage collapses under trace refinement.  The remaining analytic input should
be a trace-resolution lemma for the moving defect field:
\[
  \sup_a |\Lambda_a f|
  \le C\left(\max_j|\Lambda_{a_j}f|+h^{m-1/2}\|f\|_{W^{m,2}}\right),
\]
with constants compatible with the \(A\)-energy norm on the low/mid spectral
block.  This would justify passing from sampled finite certificates to the
closed continuum trace space \(N=\ker(\Lambda_a)_{0\le a<s_*}\).

The coefficient field is numerically smooth enough for this route.  On
\([0.02,0.545]\), using the sign convention \(e_0(a)<0\), the endpoint defect
eigenline remains one-dimensional with minimum spectral gap about
\(1.16\cdot10^{-5}\) and consecutive sampled eigenvector cosine at least
\(0.991\).  Finite differences give
\[
  \sup_a\|e(a)\|\simeq1,\qquad
  \sup_a\|e'(a)\|\simeq5.80,\qquad
  \sup_a\|e''(a)\|\simeq52.7.
\]
For \(F(a)=\Lambda_a(f)\),
\[
  F'(a)=
  \sum_k e_k'(a){f^{(k)}(a)\over k!}
  +\sum_k e_k(a){f^{(k+1)}(a)\over k!}.
\]
Thus a standard mesh estimate reduces trace resolution to energy control of
the jets \(f^{(k)}(a)\) on the finite transition block, together with the
separate high-frequency Hardy tail estimate.

The finite operator certificate for this lemma is direct.  For
\(T(a)f=\Lambda_a(f)\), cutoff \(M\), and fill distance \(\delta\),
\[
  \|T_{\rm dense}\|_{A\to\ell^\infty,E_M}
  \le
  \|T_{\rm sample}\|_{A\to\ell^\infty,E_M}
  +\delta\|\partial_aT\|_{A\to\ell^\infty,E_M}.
\]
At basis \(16\), using dense \(33\)-point trace evaluation, the measured
operator constants are
\[
\begin{array}{c|c|c|c|c|c}
\#\Lambda_a & M & \delta & \|T_{\rm dense}\| &
\|\partial_aT\| & {\rm bound}\\
\hline
8  & 7 & 3.75\cdot10^{-2} & 5.86\cdot10^3 & 3.57\cdot10^5 & 1.34\cdot10^4\\
10 & 6 & 2.92\cdot10^{-2} & 5.49          & 3.35\cdot10^2 & 9.76\\
12 & 4 & 2.39\cdot10^{-2} & 5.58\cdot10^{-4} & 3.40\cdot10^{-2} & 8.11\cdot10^{-4}
\end{array}
\]
The sampled trace operator itself is at roundoff on \(\ker R\).  This confirms
the mechanism: dense trace is controlled by fill distance times the derivative
trace operator, and the derivative trace operator collapses as the moving trace
condition is resolved.

The corresponding coupled high-frequency scan increases the basis and trace
sample count together, with approximately \(0.625\,{\rm basis}\) trace rows:
\[
\begin{array}{c|c|c|c|c|c}
{\rm basis} & \#\Lambda_a & \dim\ker R & {\rm leakage} &
C_{\rm tail}\le .1 C_{\rm full} & C_{\rm tail}\le .05 C_{\rm full}\\
\hline
12 &  8 & 4 & 4.48\cdot10^{-4} & 3 & 4\\
14 &  9 & 5 & 1.15\cdot10^{-3} & 4 & 4\\
16 & 10 & 6 & 1.27\cdot10^{-3} & 4 & 4\\
18 & 11 & 7 & 7.12\cdot10^{-4} & 6 & 6
\end{array}
\]
Thus the source tail decays once the transition modes are included, but the
transition window itself moves with refinement.  The high-frequency theorem
should therefore be stated with a moving cutoff or spectral threshold:
\[
  C_{\rm tail}(M)\le \epsilon(M)C_{\rm full},\qquad \epsilon(M)\to0,
\]
after the endpoint trace-resolution theorem has removed the low-mode sampling
artifact.

A sharper sufficient condition uses the scalar source envelope.  If
\(A e_j=\lambda_j e_j\) and \(b_j(u)=\langle Bu,e_j\rangle\), then
\[
 C_{\rm tail}(M)
 =
 \left\|\sum_{j\ge M}\lambda_j^{-1}b_j^*b_j\right\|
 \le
 \sum_{j\ge M}\lambda_j^{-1}\|b_j\|^2.
\]
Thus it is enough to prove a summable majorant for
\[
  s_j={\lambda_j^{-1}\|b_j\|^2\over C_{\rm full}}.
\]
In the coupled sections the scalar envelope is nearly sharp after the
transition window:
\[
\begin{array}{c|cc|cc}
{\rm basis} & C_{\rm tail}(4) & \sum_{j\ge4}s_j
& C_{\rm tail}(6) & \sum_{j\ge6}s_j\\
\hline
12 & 0      & 0      & 0      & 0\\
14 & .0236 & .0236 & 0      & 0\\
16 & .0382 & .0383 & 0      & 0\\
18 & .1914 & .1918 & .0137 & .0137\\
20 & .4382 & .6155 & .1733 & .1738
\end{array}
\]
The basis \(20\) row shows that a fixed mode-index cutoff is not the invariant
object: new very small eigenvalues are inserted at the low end.  The source
packet shifts in index but is far more stable spectrally.  For basis \(20\),
the dominant normalized envelope entries are
\[
  s_2=.101,\quad s_3=.640,\quad s_4=.301,\quad
  s_5=.141,\quad s_6=.158,\quad s_7=.016,
\]
with peak eigenvalue \(\lambda_3\simeq3.62\cdot10^{-6}\).  Across bases
\(14,16,18,20\) the peak remains near \(3\cdot10^{-6}\), while its mode index
moves from \(0\) to \(3\).  The tail above the fixed spectral cutoff
\(\lambda\ge10^{-2}\) is \(2.36\%,.57\%,1.37\%,1.61\%\) for bases
\(14,16,18,20\), respectively.

The analytic target is therefore a spectral source-envelope theorem:
\[
  \lambda_j^{-1}\|b_j\|^2\le a(\lambda_j) C_{\rm full},
  \qquad \sum_{\lambda_j\ge\Lambda}a(\lambda_j)\to0
  \quad(\Lambda\to\infty).
\]
The finite Schur block then covers the low/mid spectral source packet, while
the Hardy/source-envelope estimate proves the true high-frequency tail.
\end{problem}

\begin{lemma}[Abstract spectral-tail reduction]
Let \(A\ge0\) be a self-adjoint energy operator, let \(P_{\ge\Lambda}\) be its
spectral projection, and let \(B:U\to H\) be a source operator.  For \(m>0\),
if \(A^{m/2}B\) is bounded, then
\[
  \|A^{-1/2}P_{\ge\Lambda}B\|^2
  \le
  \Lambda^{-(m+1)}\|A^{m/2}B\|^2 .
\]
If \(A e_j=\lambda_j e_j\), \(b_j(u)=\langle Bu,e_j\rangle\), and
\(A^{m/2}B\) is Hilbert--Schmidt, then
\[
  \sum_{\lambda_j\ge\Lambda}\lambda_j^{-1}\|b_j\|^2
  \le
  \Lambda^{-(m+1)}\|A^{m/2}B\|_{\rm HS}^2 .
\]
\end{lemma}

\begin{proof}
On \(\operatorname{Ran}P_{\ge\Lambda}\), the spectral theorem gives
\(A^{-(m+1)/2}\le\Lambda^{-(m+1)/2}\).  Hence
\[
  A^{-1/2}P_{\ge\Lambda}B
  =
  A^{-(m+1)/2}P_{\ge\Lambda}A^{m/2}B,
\]
which proves the operator estimate after taking norms.  The scalar estimate is
the same inequality summed in the eigenbasis of \(A\).
\end{proof}

Thus the remaining analytic task is to prove a commuted source estimate:
after imposing the closed continuum trace condition, the Volterra source
operator \(B\) must gain a positive number of \(A\)-derivatives.  In concrete
terms this should follow by integrating the source functional by parts through
the endpoint Sturm operator; the boundary terms must be exactly the trace
functionals killed by the range condition.

\begin{lemma}[Green criterion for the commuted source estimate]
Let \(N=\ker R\), where \(R f=(\Lambda_a f)_{0\le a<s_*}\), and let
\(h_u=Bu\) represent the cross form \(L_u(n)=\langle h_u,n\rangle\) on \(N\).
Assume that for some integer \(r\ge1\) there are \(g_u\) and boundary
distributions \(\eta_u\) such that the Green identity
\[
  \langle h_u,A^r n\rangle
  =
  \langle g_u,n\rangle+\langle \eta_u,Rn\rangle
\]
holds for smooth test functions, with
\(\|g_u\|\le M_r\|u\|_U\).  Then \(A^rB\) is bounded on the closed trace space:
\[
  \sum_j \lambda_j^{2r}|\langle h_u,e_j\rangle|^2
  \le M_r^2\|u\|_U^2 .
\]
Consequently the spectral tail satisfies
\[
  C_{\rm tail}(\Lambda)
  \le M_r^2\Lambda^{-(2r+1)}
\]
after the same normalization as in the abstract spectral-tail lemma.
\end{lemma}

\begin{proof}
For \(e_j\in N\) with \(A e_j=\lambda_j e_j\), the trace term vanishes and
\[
  \lambda_j^r\langle h_u,e_j\rangle=\langle g_u,e_j\rangle .
\]
Bessel's inequality gives the stated bound.
\end{proof}

The finite commuted-source diagnostic supports this target.  On the endpoint
\(B\)-model, \(\widetilde\Phi_3\), \(\omega=.49\), and the same coupled
trace-refined sections as above, the normalized \(m=1\) commuted source norm
\(\|A^{1/2}B\|^2/C_{\rm full}\) is
\[
\begin{array}{c|ccccc}
{\rm basis} & 12 & 14 & 16 & 18 & 20\\
\hline
\|A^{1/2}B\|^2/C_{\rm full}
&1.31\cdot10^{-3}&5.93\cdot10^{-4}&1.40\cdot10^{-4}
&3.35\cdot10^{-4}&3.97\cdot10^{-4}.
\end{array}
\]
The values remain small even when the raw source packet shifts in mode index.
Because the finite spectra here still have \(\lambda_{\max}\simeq.154<1\),
the numerical high-\(\Lambda\) bound itself is not sharp; the important
evidence is boundedness of the commuted source norm.

The concrete endpoint Sturm identity can be written without ambiguity.  Put
\[
  x=1+\tau,\qquad \lambda=ce^s,\qquad
  F_s(\tau)={e^{-\lambda\tau}-e^{-c\tau}\over \lambda-c},
  \qquad
  H_s(\tau)={s e^{-\lambda\tau}\over \lambda-c},
\]
and \(D Y=(xY)'\).  The endpoint \(B\)-model kernel is
\[
\begin{aligned}
K_B(s,t)=\int_0^\infty x^{3/2}\{&
\log x\,D F_sD F_t
+\frac12 D F_sD H_t
+\frac12 D H_sD F_t\}\,d\tau .
\end{aligned}
\]
For \(a(\tau)\), define
\[
  L_aG=-x{d\over d\tau}\{a x^{3/2}DG\}.
\]
The Green identity
\[
  \int a x^{3/2}DP\,DQ\,d\tau
  =
  [a x^{3/2}DQ\cdot xP]_0^\infty
  +\int P\,L_aQ\,d\tau
\]
is then oriented with the boundary-null feature first.  Since \(F_s(0)=0\) and
\(\log x=0\) at \(\tau=0\),
\[
  K_B(f,u)
  =
  \int F_f L_{\log}F_u\,d\tau
  +{1\over2}\int F_f L_1H_u\,d\tau
  +{1\over2}\int F_u L_1H_f\,d\tau .
\]
Thus the physical Sturm endpoint contributes no boundary defect.  The source
functional is explicitly
\[
  h_u(s)=
  \int F_s\bigl(L_{\log}F_u+\tfrac12L_1H_u\bigr)\,d\tau
  +{1\over2}\int (L_1H_s)F_u\,d\tau .
\]
The remaining commutation problem is therefore in the spectral variable \(s\):
one must show that integration by parts against the moving trace operator
\(\Lambda_s(f)=\sum_{k=0}^8 e_k(s)f^{(k)}(s)/k!\) produces only boundary
concomitants generated by \(\Lambda_s f\) and its \(s\)-derivatives.

The algebraic closure of this moving trace is triangular.  If
\(\Lambda_s(f)=0\) for every active \(s\), then
\[
  D_s^q\Lambda_s(f)=
  \sum_{k=0}^8\sum_{m=0}^q {q\choose m}
  {e_k^{(q-m)}(s)\over k!}f^{(k+m)}(s)=0 .
\]
The coefficient of the newest jet \(f^{(8+q)}(s)\) is \(e_8(s)/8!\).  The
computed moving defect field has no zero of \(e_8\) on the active interval:
on \([.02,.545]\), a 17-point high-precision scan gives
\[
  \min |e_8(s)|\simeq 1.52\cdot10^{-2},
  \qquad
  \operatorname{rank}\{D_s^q\Lambda_s:0\le q\le4\}=5.
\]
Thus the closed trace condition recursively eliminates
\(f^{(8)},f^{(9)},\ldots\).  The remaining analytic task is a membership
statement: the \(s\)-boundary concomitant obtained from commuting \(h_u(s)\)
must be expressible as
\[
  \sum_{q=0}^Q \alpha_q(s,u)D_s^q\Lambda_s
\]
with bounded coefficients.

The naive scalar version of this membership statement is false.  If
\(h_u(s)=K_B(s,u)\) is treated as a scalar source and \(D_s^Q\) is integrated by
parts, the endpoint row is
\[
  B_Q[h_u](f)=\sum_{j=0}^{Q-1}(-1)^j h_u^{(Q-1-j)}(s_0)f^{(j)}(s_0).
\]
Since \(\Lambda_s\) has order \(8\), this row would have to lie in
\(\operatorname{span}\{D_s^q\Lambda_s:0\le q\le Q-9\}\).  A high-precision
Taylor-jet test at \(s_0=.2825\), with \(u=.08,.24,.40,.52\) and
\(Q=9,\ldots,13\), gives relative residuals ranging from \(1.38\cdot10^{-1}\)
to \(9.95\cdot10^{-1}\).  Thus raw repeated \(s\)-integration by parts is not
the source concomitant.  The concomitant must instead be computed before
scalarization, from the Sturm-adjoint commutator acting on the feature pair
\((F_s,H_s)\) in the Green formula for \(h_u\).

The first viable scalar remnant is the Green concomitant of the trace
differential expression itself.  Write
\[
  P f=\Lambda_s(f)=\sum_{k=0}^8 a_k(s)f^{(k)}(s),
  \qquad a_k(s)={e_k(s)\over k!}.
\]
The corresponding Lagrange boundary row is
\[
  B_P[h_u](f)=
  \sum_{k=1}^8\sum_{j=0}^{k-1}
    (-1)^{k-1-j} D_s^{k-1-j}(a_kh_u)(s_0)f^{(j)}(s_0).
\]
A high-precision membership test against
\(\operatorname{span}\{D_s^q\Lambda_s:0\le q\le q_{\max}\}\) gives small
residuals.  Replacing finite differences by analytic derivatives of \(e(s)\)
from the confluent eigenvalue equation \(J(s)e(s)=\lambda(s)e(s)\) improves the
residual to approximately \(2.9\cdot10^{-5}\) at \(q_{\max}=10\).

Nevertheless, exact finite ideal membership is impossible in the ordinary
differential-operator sense.  Since \(a_8(s)=e_8(s)/8!\ne0\), any finite
identity
\[
  B_P[h_u]=\sum_{q=0}^Q\alpha_qD_s^q\Lambda_s
\]
has highest-jet coefficient \(\alpha_Qa_8\) at \(f^{(8+Q)}\), forcing
\(\alpha_Q=0\), and then descending forces all \(\alpha_q=0\).  But the
\(f^{(7)}\)-coefficient of \(B_P[h_u]\) is \(-(e_8/8!)h_u(s_0)\), which is
nonzero for the source points tested.  Thus the small residual is a high-jet
near-dependence, not a literal local ideal identity.  The corrected target is
to use the Lagrange identity
\[
  D_s B_P[h_u,f]=h_u\,P f-f\,P^*h_u
\]
on the closed trace space \(P f=0\), and then prove that the remaining adjoint
source term and endpoint pairing are controlled by the commuted Sturm energy.
This identity has been checked at the row level using analytic derivatives of
the confluent eigenvector \(e(s)\), with defect below \(3\cdot10^{-79}\) in the
test above.  The term \(P^*h_u\) is not small: at \(s_0=.2825\), for
\(u=.08,.24,.40,.52\), it ranges from roughly \(-107.4\) to \(-93.1\), about
an order of magnitude larger than the Euclidean norm of the boundary row.
Therefore the next estimate must control the adjoint source itself, not merely
discard the endpoint concomitant.

The finite Galerkin control constants confirm the same point.  On the sampled
closed trace quotient with basis \(16\) and \(10\) trace constraints, the
smallest positive endpoint-energy eigenvalue is \(4.69\cdot10^{-8}\).  The
sharp \(A^{-1}\)-norms of \(B_P[h_u,\cdot]\) are about \(7.1\cdot10^3\), while
the norms of the local adjoint-source rows \(P^*h_u(s_0)\operatorname{ev}_{s_0}\)
are about \(6.5\cdot10^4\) to \(7.5\cdot10^4\).  Thus the residual terms are
finite in the Galerkin quotient but are dominated by the tiny low positive
mode.  The continuum theorem should therefore be formulated with the commuted
Sturm/Hardy graph norm, or by first removing the low Schur block and then
proving a high-frequency Hardy trace estimate.

The resulting theorem has the following form.  Let \(N=\ker R\) be the closed
trace space and \(A=K|_N\ge0\).  Let \(E_u\) denote the two-component residual
operator consisting of the endpoint concomitant and the adjoint source term
left by the Lagrange identity.  If \(L_M\) is the span of the first \(M\)
positive \(A\)-modes, then one must prove:
\[
  \text{Schur}_{L_M}(K)\ge0,
  \qquad
  \|E_u f\|^2\le C_H(M,u)\langle Af,f\rangle
  \quad(f\perp_A L_M),
\]
with \(C_H(M,u)\) absorbable by the positive commuted Sturm energy.  The high
estimate should follow from
\[
  \|A^{-1/2}P_{\lambda\ge\Lambda}E_u^*\|^2
  \le
  \Lambda^{-(m+1)}\|A^{m/2}E_u^*\|^2 .
\]
In the finite endpoint model, the corresponding residual control matrix has
high-tail fractions about \(6.2\cdot10^{-3}\) after removing two modes,
\(6.5\cdot10^{-4}\) after three modes, \(1.7\cdot10^{-4}\) after four modes,
and \(8.8\cdot10^{-6}\) after five modes.  Thus the data support precisely a
finite low Schur block plus a high-frequency commuted Hardy estimate.

This split is stable under a coupled basis/trace refinement.  For bases
\(12,14,16,18\), with approximately \(0.625\) trace constraints per basis
function, removing two low modes puts the high tail below \(10^{-2}\) from
basis \(14\) onward; removing three modes puts it below \(10^{-3}\) from basis
\(16\) onward; and removing four or five modes puts it near or below
\(10^{-4}\).  Numerically, the low transition block does not drift upward with
the Galerkin dimension.  A conservative analytic statement should therefore
use an unspecified finite \(M\), with the endpoint model suggesting \(M\le5\).

The high block also satisfies the expected finite Hardy graph estimate.  In the
basis \(18\), constraints \(11\) quotient, the sharp constants in
\[
  \|E_u f\|^2\le C_{m,M}(u)\|f\|_{W^{m,2}(0,L)}^2,
  \qquad f\in H_M,
\]
drop rapidly with \(m\).  At \(M=5\), the largest constants over the four tested
source points are \(C_{6,5}<9.7\cdot10^{-2}\),
\(C_{8,5}<3.9\cdot10^{-6}\), and \(C_{10,5}<2.4\cdot10^{-8}\).  Thus the
remaining continuum input is a high-order endpoint trace theorem together with
the commuted Sturm estimate controlling the \(W^{m,2}\) graph norm on the high
closed-trace block.

The trace theorem itself is standard: for \(m\ge8\), point evaluations of
\(f,\ldots,f^{(7)}\) are bounded by \(\|f\|_{W^{m,2}(0,L)}\), and the
coefficients in \(E_u\) are analytic on compact \(s_0,u\)-windows.  The hard
part is the elliptic estimate on the high closed-trace block.  The uncommuted
endpoint energy is not sufficient: in the same basis \(18\) model, the
constants in
\[
  \|f\|_{W^{m,2}}^2\le G_{m,M}\langle Af,f\rangle
  \qquad(f\in H_M)
\]
are already \(G_{6,5}\simeq 1.1\cdot10^{10}\),
\(G_{8,5}\simeq 1.8\cdot10^{13}\), and
\(G_{10,5}\simeq 3.3\cdot10^{17}\).  Thus the final analytic ingredient must
be a genuinely commuted Sturm energy \(S_m\) satisfying
\[
  \langle S_m f,f\rangle \ge c_m\|f\|_{W^{m,2}}^2
  \qquad (f\in H_M\cap\ker R),
\]
with the boundary terms controlled by the finite Lagrange residual block.

The finite commuted-kernel analogue
\[
  S_m^{\rm comm}=\sum_{r=0}^m (D_s^r)^*K(D_s^r)
\]
does close the numerical high-block estimate.  On the basis \(18\),
constraints \(11\) quotient, at cutoff \(M=5\),
\[
  \|f\|_{W^{8,2}}^2\le 2.66\cdot10^3
  \langle S_8^{\rm comm}f,f\rangle,\qquad
  \|f\|_{W^{10,2}}^2\le 1.24\cdot10^3
  \langle S_{10}^{\rm comm}f,f\rangle.
\]
Composed with the finite Hardy graph estimate for the residual row \(E_u\),
this gives
\[
  \|E_u f\|^2\le 1.03\cdot10^{-2}
  \langle S_8^{\rm comm}f,f\rangle,\qquad
  \|E_u f\|^2\le 2.89\cdot10^{-5}
  \langle S_{10}^{\rm comm}f,f\rangle .
\]
Thus the finite high-frequency theorem closes with margin.  The continuum
missing step is precisely the commuted Sturm elliptic estimate for the true
Green identity, followed by Galerkin compactness.

A basis-refinement scan changes the conservative formulation of this target.
At basis \(20\), \(M=5,m=10\) is still subunit but only gives
\[
  D_{10,5}=6.51\cdot 10^{-1}.
\]
Removing one more finite mode gives a stable high-block estimate:
\[
  D_{10,6}=2.51\cdot10^{-7},\qquad
  D_{12,6}=8.85\cdot10^{-11}.
\]
Accordingly, the analytic theorem to prove should be stated on
\(H_6\cap\ker R\), with at least ten commutations:
\[
  \|E_u f\|^2\le \eta\,\langle S_{10}f,f\rangle,\qquad
  \eta<1 .
\]
The \(m=12\) estimate is a high-margin fallback if endpoint constants are lost
in the continuum passage.

One must not interpret \(S_m^{\rm comm}\) as the whole continuum elliptic
form.  The endpoint kernel \(K\), treated as an integral operator on
\([0,L]\), is compact on interior packets.  Hence the compact quadratic form
\[
  \sum_{r=0}^m \langle K D^r f,D^r f\rangle
\]
cannot dominate \(\|f\|_{W^{m,2}}^2\) on any infinite-dimensional
finite-codimension subspace.  Indeed, if
\(\chi\in C_c^\infty((s_\ast,L))\) and
\[
  f_n(s)=n^{-m}\chi(s)e^{ins},
\]
then \(Rf_n=0\), \(\|f_n\|_{W^{m,2}}\asymp 1\), the lower derivative terms
vanish, and \(D^m f_n\) is bounded and weakly null in \(L^2\).  Compactness of
\(K\) gives
\[
  \langle K D^m f_n,D^m f_n\rangle\to0.
\]
Finite Schur-block orthogonality can be imposed by an \(o(1)\) correction.
Thus the continuum estimate must come from the local positive principal part
produced by the true commuted Sturm Green identity, not from the compact
commuted kernel alone.  The precise analytic target is therefore a Garding
estimate for the ten-times commuted Sturm form:
\[
  S_{10}(f)=\int_0^L a_{10}(s)|D^{10}f(s)|^2\,ds
  +\hbox{lower order terms}+\hbox{trace squares},
  \qquad a_{10}(s)\ge a_0>0,
\]
followed by removal of the six finite Schur modes.

This last formulation is a target, not yet a theorem.  In fact, the current
endpoint identity does not by itself produce such an \(a_{10}\).  If one had an
exact identity
\[
  K_B(f,f)=\int a_{10}(s)|D^{10}f(s)|^2\,ds
  +\hbox{lower order local terms}+\hbox{trace squares}
\]
with \(a_{10}>0\) on an interior interval, the oscillatory packets
\(f_n=n^{-10}\chi e^{ins}\), with
\(\chi\in C_c^\infty((s_\ast,L))\), would give a contradiction:
the trace terms vanish, \(K_B(f_n,f_n)\to0\), and lower-order local terms
vanish after \(W^{10,2}\)-normalization, while the top term has a positive
limit.  Thus no positive local principal coefficient can be extracted from
the compact endpoint form itself.

The finite feature-commutator route also fails as an exact algebraic identity.
For fixed \(s\), finite \(s\)-derivatives of \((F_s,H_s)\) span functions of
the form
\[
  P(\tau)e^{-\lambda\tau}+C e^{-c\tau},
\]
with \(P\) a finite polynomial.  The Sturm sources \(L_1H_s\) and
\(L_{\log}F_s\) contain the non-polynomial factors \(x^{3/2}\) and
\(x^{3/2}\log x\), respectively.  Numerically, order-ten \(s\)-jets approximate
these sources very well on finite weighted windows, but this is an
approximation phenomenon and does not define an intrinsic principal coefficient
\(a_{10}(s)\).

The corrected analytic route is therefore not to prove a local \(a_{10}\)
identity for \(K_B\).  It is to introduce an explicit auxiliary high-block
regularizer, prove the Hardy trace residual is controlled by it on
\(H_6\cap\ker R\), and then prove that regularizer is absorbable by the
positive Volterra/Schur block.

The auxiliary regularizer must be source-range, not Sobolev.  On the sampled
high block define
\[
  S_{\rm src}(f)=\sum_{t\in T}\|E_t f\|^2,\qquad
  E_t f=(B_P[h_t,f](s_0),\,P^\ast h_t(s_0)f(s_0)).
\]
At basis \(20\), constraints \(12\), and cutoff \(M=6\), the Sobolev route
fails:
\[
  \|E f\|^2\le 2.94\cdot10^{-10}\|f\|_{W^{10,2}}^2,\qquad
  \|f\|_{W^{10,2}}^2\le 4.98\cdot10^{20}\langle Af,f\rangle,
\]
whose product is \(1.46\cdot10^{11}\).  In contrast the source-range
regularizer is directly absorbable:
\[
  S_{\rm src}(f)\le 3.28\cdot10^6\langle Af,f\rangle
  \quad\hbox{on }H_6\cap\ker R,
\]
and this is only \(1.19\cdot10^{-6}\) of the full Schur budget.  The continuum
target is therefore a Hardy/Green estimate for the source rows themselves:
\[
  E_t^\ast E_t\le \eta_t A\qquad (f\in H_6\cap\ker R),
\]
uniformly in the source window.  The finite low block carries the large Schur
mass; the high source-range tail should be proved small directly, without
passing through a full Sobolev norm.

A 9-node Gauss scan over the source window \(u\in[.08,.52]\) supports the
uniform version.  At basis \(20\), constraints \(12\), and \(M=6\),
\[
  \frac{\eta_{\rm high}}{\eta_{\rm full}}
  =1.193754965\cdot10^{-6}
\]
for the quadrature-integrated source regularizer
\(\int_{.08}^{.52}\|E_u f\|^2\,du\).  The worst single-node fraction is
\[
  1.194356816\cdot10^{-6},
\]
attained near \(u=.0870\).  The fraction varies only in the fourth decimal
place across the source window, so the remaining continuum theorem should be a
uniform source-window Hardy/Green estimate for \(E_u^\ast E_u\) on
\(H_6\cap\ker R\).

A refinement scan over both the Galerkin basis and the source quadrature
supports the same conclusion.  With the trace rule
\(\operatorname{constraints}=\lfloor .625\,\operatorname{basis}\rfloor\),
source counts \(5,9,13\), and the same cutoff \(M=6\), the integrated
high/full fractions are
\[
\begin{array}{c|c|c|c}
\operatorname{basis} & \operatorname{constraints} &
\operatorname{positive\ modes} &
\eta_{\rm high}/\eta_{\rm full} \\\hline
18 & 11 & 7  & 2.245269967\cdot 10^{-7}\\
20 & 12 & 8  & 1.193754965\cdot 10^{-6}\\
22 & 13 & 10 & 5.502913863\cdot 10^{-7}.
\end{array}
\]
For each fixed basis the value is unchanged, to the displayed precision, when
the source quadrature is refined from 5 to 13 nodes.  Thus the present finite
obstruction is not the source discretization.  The analytic target is the
continuum source-window Hardy/Green estimate for the smooth family \(E_u\),
followed by the finite low Schur comparison.

The finite-net version is an operator-Lipschitz estimate in the \(A\)-energy
metric.  If
\[
  E_u^\ast E_u\le \eta(u)A,\qquad
  \left({E_v-E_u\over v-u}\right)^\ast
  \left({E_v-E_u\over v-u}\right)\le L(u,v)^2 A,
\]
then a source mesh of radius \(h/2\) gives the conservative covering bound
\[
  E_u^\ast E_u
  \le 2\max_j \eta(u_j)A+2(h/2)^2 L^2 A .
\]
At basis \(20\), constraints \(12\), cutoff \(M=6\), and a 17-point uniform
mesh on \([.08,.52]\), the sampled maximum high/full fraction is
\[
  1.194375338\cdot10^{-6},
\]
the maximum high-block Lipschitz \(A\)-constant is \(1.291568103\cdot10^5\),
and the resulting finite-net covering fraction, measured against the minimum
sampled full block, is
\[
  3.172132313\cdot10^{-6}.
\]
Thus even a crude continuum covering remains at the \(10^{-6}\) scale.  The
proof target is to obtain this Lipschitz estimate analytically from the
Green/source formula for \(h_u\) and the closed trace condition.

The analytic derivative is explicit.  For the endpoint-B feature pair \(V,W\),
\[
h_u^{(k)}(s_0)=\int_0^\infty e^{5r/2}
\left[
r\,\partial_s^kV(s_0,r)V(u,r)
+{1\over2}\{\partial_s^kV(s_0,r)W(u,r)
+\partial_s^kW(s_0,r)V(u,r)\}
\right]\,dr,
\]
and therefore
\[
\partial_u h_u^{(k)}(s_0)=\int_0^\infty e^{5r/2}
\left[
r\,\partial_s^kV(s_0,r)\partial_uV(u,r)
+{1\over2}\{\partial_s^kV(s_0,r)\partial_uW(u,r)
+\partial_s^kW(s_0,r)\partial_uV(u,r)\}
\right]\,dr.
\]
Consequently
\[
  \partial_u E_u f
  =
  \bigl(B_P[\partial_u h_u,f](s_0),
        P^\ast(\partial_u h_u)(s_0)f(s_0)\bigr).
\]
Numerically, the analytic derivative agrees with centered finite differences
at \(u=.08,.30,.52\) to relative error \(<4\cdot10^{-9}\).  On the same
17-point source grid the maximum analytic derivative \(A\)-constant on the
high block is \(1.311726997\cdot10^5\), the maximum derivative high/full
fraction is \(1.201677229\cdot10^{-6}\), and the resulting analytic covering
fraction is
\[
  3.172133606\cdot10^{-6}.
\]
Thus the finite evidence now tests exactly the derivative estimate required by
the continuum Hardy/Green theorem.

The Hardy/Green theorem itself is a representer statement.  On the closed
trace high block \(H_{\rm hi}\), with inner product
\[
  \langle f,g\rangle_A=\langle Af,g\rangle,
\]
each scalar component \(\ell\) of \(E_u\) and \(\partial_uE_u\) should have a
Green representer \(g_\ell\in H_{\rm hi}\) satisfying
\[
  \ell(f)=\langle g_\ell,f\rangle_A,\qquad f\in H_{\rm hi}.
\]
Then Cauchy's inequality gives
\[
  |\ell(f)|^2\le \|g_\ell\|_A^2\,\langle Af,f\rangle.
\]
For a two-row operator the sharp constant is the largest eigenvalue of the
two-by-two Gram matrix \((\langle g_i,g_j\rangle_A)\).

The finite Galerkin representers are explicit.  If \(Ae_j=\lambda_j e_j\) and
\(\ell_j=\ell(e_j)\), then on the high block
\[
  g_\ell=\sum_{j\ge6}{\ell_j\over\lambda_j}e_j,\qquad
  \|g_\ell\|_A^2=\sum_{j\ge6}{|\ell_j|^2\over\lambda_j}.
\]
At basis \(20\), constraints \(12\), cutoff \(M=6\), the finite range identity
has relative defect \(1.68\cdot10^{-80}\).  The sampled constants are
\[
\begin{array}{c|c|c}
u & E_u & \partial_uE_u\\\hline
.08 & 9.346947046\cdot10^5 & 3.863008587\cdot10^4\\
.30 & 8.313998380\cdot10^5 & 8.945221434\cdot10^4\\
.52 & 7.030785145\cdot10^5 & 1.311726997\cdot10^5 .
\end{array}
\]
The dominant row is the adjoint-evaluation component; the boundary component,
and especially its \(u\)-derivative, is much smaller.  Thus the remaining
continuum proof is to construct these Green representers analytically and
bound their \(A\)-norms uniformly in \(u\).

The dominant row factorizes.  Write
\[
  p(u)=P^\ast h_u(s_0).
\]
Then the adjoint-evaluation component is
\[
  \ell_u^{\rm eval}(f)=p(u)f(s_0).
\]
Consequently its Green representer is
\[
  g_u^{\rm eval}=p(u)\,k_{s_0}^{\rm hi},\qquad
  \partial_u g_u^{\rm eval}=p'(u)\,k_{s_0}^{\rm hi},
\]
where \(k_{s_0}^{\rm hi}\) is the \(A\)-Green representer of point evaluation
on the closed-trace high block.  In RKHS notation,
\[
  k_{s_0}^{\rm hi}=P_{\rm hi}P_{\ker R}K_A(\cdot,s_0),
  \qquad
  P_{\ker R}=I-R^\ast(RR^\ast)^+R,
\]
and \(P_{\rm hi}\) removes the first six closed-trace \(A\)-modes.  Thus
\[
  |p(u)f(s_0)|^2
  \le |p(u)|^2\,\|k_{s_0}^{\rm hi}\|_A^2\,\langle Af,f\rangle,
\]
with the same statement for \(p'(u)\).

At basis \(20\), constraints \(12\), cutoff \(M=6\), the finite projected
evaluation representer has
\[
  \|k_{s_0}^{\rm hi}\|_A^2=80.1690618322
\]
and the factorization defect is \(5.20\cdot10^{-81}\).  A 33-point source
mesh, using \(p,p'\), and \(p''\), gives
\[
\sup |p|\lesssim 107.695130153,\qquad
\sup |p'|\lesssim 40.8158694613,
\]
and hence the adjoint-evaluation constants
\[
  9.298201046\cdot10^5,\qquad 1.335564620\cdot10^5
\]
for \(E_u\) and \(\partial_uE_u\), respectively.  The continuum proof of this
piece is therefore reduced to boundedness of the projected evaluation
representer and compact-window bounds for the explicit scalar Volterra
amplitude \(p(u)\).

The boundary row has an analogous finite-jet form.  There are coefficient
functions \(b_j(u)\), \(0\le j\le7\), such that
\[
  B_P[h_u,f](s_0)=\sum_{j=0}^7 b_j(u)f^{(j)}(s_0),
\]
with
\[
  b_j(u)=\sum_{k=j+1}^8(-1)^{k-1-j}
  \partial_s^{k-1-j}(a_k(s)h_u(s))\big|_{s=s_0}.
\]
Let \(k_j^{\rm hi}\) be the \(A\)-Green representer of
\(f\mapsto f^{(j)}(s_0)\) on the closed-trace high block.  Then
\[
  g_u^{\rm bdry}=\sum_{j=0}^7b_j(u)k_j^{\rm hi},
  \qquad
  \partial_u g_u^{\rm bdry}=\sum_{j=0}^7b'_j(u)k_j^{\rm hi}.
\]
If \(G_{ij}=\langle k_i^{\rm hi},k_j^{\rm hi}\rangle_A\), then
\[
  \|g_u^{\rm bdry}\|_A^2=b(u)^TG b(u),\qquad
  \|\partial_u g_u^{\rm bdry}\|_A^2=b'(u)^TG b'(u).
\]
At basis \(20\), constraints \(12\), cutoff \(M=6\), the finite jet
representers satisfy the range equations with relative defect
\(2.21\cdot10^{-80}\), and the factorization defect is
\(1.97\cdot10^{-80}\).  The mesh-covered norms are
\[
  \|b\|_G\lesssim 98.4429389340,\qquad
  \|b'\|_G\lesssim 18.8407728129,
\]
which give
\[
  \|g_u^{\rm bdry}\|_A^2\lesssim 9.691012226\cdot10^3,\qquad
  \|\partial_u g_u^{\rm bdry}\|_A^2\lesssim 3.549747202\cdot10^2.
\]
Thus the boundary row, too, is reduced to fixed jet representers and compact
Volterra coefficient bounds.

This gives the finite fixed-representer version of the closed-trace
Hardy/Green theorem.  Let \(H_{\rm hi}\) be the closed-trace high RKHS with
\(A\)-inner product.  Suppose the fixed representers
\[
  k_{s_0}^{\rm hi},\qquad k_j^{\rm hi}\quad(0\le j\le7)
\]
exist in \(H_{\rm hi}\), where \(k_{s_0}^{\rm hi}\) represents point
evaluation at \(s_0\), and \(k_j^{\rm hi}\) represents
\(f\mapsto f^{(j)}(s_0)\).  If
\[
  E_u f=
  \left(
    \sum_{j=0}^7 b_j(u)f^{(j)}(s_0),
    p(u)f(s_0)
  \right),
\]
then, with \(G_{ij}=\langle k_i^{\rm hi},k_j^{\rm hi}\rangle_A\),
\[
  \|E_u f\|^2
  \le
  \left(|p(u)|^2\|k_{s_0}^{\rm hi}\|_A^2+b(u)^TG b(u)\right)
  \langle Af,f\rangle,
\]
and the same bound holds for \(\partial_uE_u\) after replacing \(p,b\) by
\(p',b'\).  Thus the continuum source-window Hardy/Green theorem follows from
boundedness of these fixed representers and compact-window bounds on the
scalar coefficient functions.

The finite scan \texttt{fixed\_representer\_theorem\_scan.py} checks this
reduction across nearby Galerkin sections.  On \(u\in[.08,.52]\), a 17-point
mesh with derivative covering gives
\[
  \sup |p|\lesssim 107.972859867,\qquad
  \sup |p'|\lesssim 41.2346895930.
\]
With the trace rule
\(\operatorname{constraints}=\lfloor .625\,\operatorname{basis}\rfloor\) and
cutoff \(M=6\), the fixed-representer constants are
\[
\begin{array}{c|c|c|c|c|c|c}
\operatorname{basis}&\operatorname{constraints}&\operatorname{pos}&
\lambda_{\rm hi,min}&\|k_{s_0}^{\rm hi}\|_A^2&
\operatorname{eval\ cover}&\operatorname{bdry\ cover}\\\hline
18&11&7&1.543899343\cdot10^{-1}&4.235662880&
4.937994436\cdot10^4&4.653506259\cdot10^2\\
20&12&8&9.537782084\cdot10^{-3}&80.16906183&
9.346220237\cdot10^5&9.715522428\cdot10^3\\
22&13&10&5.262714063\cdot10^{-5}&2029.620478&
2.366159657\cdot10^7&1.485357499\cdot10^5 .
\end{array}
\]
The finite range defects are at most \(1.29\cdot10^{-78}\).  The growth of the
constants as \(\lambda_{\rm hi,min}\) decreases shows that the remaining
analytic issue is not the coefficient algebra, but the continuum
closed-trace endpoint Sobolev/RKHS estimate for the fixed jet representers.

We then tested whether that estimate follows from the local continuum trace
tower alone.  The closed trace condition implies
\[
  \partial_s^q\Lambda_s(f)\big|_{s=s_0}=0,\qquad q=0,1,2,\ldots ,
\]
where
\[
  \Lambda_s(f)=\sum_{k=0}^8 e_k(s){f^{(k)}(s)\over k!}.
\]
The script \texttt{local\_trace\_tower\_representer\_scan.py} imposes these
exact tower rows at \(s_0=.2825\), using eigenderivatives of \(e(s)\) computed
from the confluent equation.  With tower constraints equal to
\(\operatorname{basis}-8\) and cutoff \(M=6\), it gives
\[
\begin{array}{c|c|c|c|c|c}
\operatorname{basis}&\operatorname{tower}&\operatorname{pos}&
\operatorname{high}&\|\operatorname{ev}_{s_0}\|_A^2&
\|D^7\operatorname{ev}_{s_0}\|_A^2\\\hline
18&10&11&5&1.969126974\cdot10^5&4.411328604\cdot10^{14}\\
20&12&12&6&3.761728684\cdot10^6&4.418566287\cdot10^{17}\\
22&14&14&8&3.938406606\cdot10^7&1.649733096\cdot10^{22}.
\end{array}
\]
The range equations are still solved to working precision, but the constants
grow rapidly.  Thus the individual fixed jet representer theorem is too
strong as a standalone target; the local trace tower does not control
arbitrary endpoint jet directions in the compact \(A\)-energy.

The correct formulation keeps the special source rows together:
\[
  E_u f=
  \left(
    \sum_{j=0}^7 b_j(u)f^{(j)}(s_0),
    p(u)f(s_0)
  \right).
\]
The remaining theorem should be a direct source-row Schur/observability
estimate
\[
  E_u^\ast E_u\le C(u)A
\]
on the closed continuum trace space after the finite Schur block has been
removed.  The individual jet Gram matrix remains useful bookkeeping, but the
analytic proof must exploit cancellation in the coefficient family
\((p(u),b(u))\).

The direct source-row test confirms this correction, but also shows that the
trace hypothesis must be global.  The script
\texttt{local\_trace\_tower\_source\_row\_scan.py} imposes the exact local
trace tower and computes the sharp high-block constants for \(E_u\) itself.
With 9 source nodes on \([.08,.52]\), tower constraints
\(\operatorname{basis}-8\), and cutoff \(M=6\), it gives
\[
\begin{array}{c|c|c|c|c|c}
\operatorname{basis}&\operatorname{tower}&\operatorname{pos}&
\operatorname{high}&\max C_{\rm high}&
\max(C_{\rm high}/C_{\rm full})\\\hline
18&10&11&5&2.298068079\cdot10^9&2.233323914\cdot10^{-7}\\
20&12&12&6&4.409786636\cdot10^{10}&2.804065742\cdot10^{-8}\\
22&14&14&8&4.605309936\cdot10^{11}&1.855528363\cdot10^{-10}.
\end{array}
\]
The differentiated source row has the same qualitative behavior; its
high/full fractions are \(2.21\cdot10^{-7}\), \(4.99\cdot10^{-7}\), and
\(9.08\cdot10^{-10}\) for the three bases.  Thus the special source-row
combination has tiny high Schur mass even though arbitrary jets do not.

The absolute constants still grow under the local tower.  Therefore the
continuum theorem should be stated on the global closed trace space
\[
  Rf=(\Lambda_a f)_{0\le a<s_\ast}=0,
\]
not merely under the local tower at \(s_0\).  The remaining analytic input is
an interval observability/range theorem combining the Lagrange identity with
global trace resolution over the active endpoint interval.

The global sampled-trace scan tests exactly this formulation.  The script
\texttt{global\_trace\_source\_observability\_scan.py} uses the sampled
interval trace operator on \([.02,.545]\), with
\(\operatorname{constraints}=\lfloor .625\,\operatorname{basis}\rfloor\), and
keeps the corrected source row \(E_u\) intact.  With 9 source nodes on
\([.08,.52]\) and cutoff \(M=6\), it gives
\[
\begin{array}{c|c|c|c|c|c}
\operatorname{basis}&\operatorname{constraints}&\operatorname{pos}&
\operatorname{high}&\max C_{\rm high}&
\max(C_{\rm high}/C_{\rm full})\\\hline
18&11&7&1&4.933603300\cdot10^4&2.246582581\cdot10^{-7}\\
20&12&8&2&9.346947046\cdot10^5&1.194375338\cdot10^{-6}\\
22&13&10&4&2.355734776\cdot10^7&5.505833523\cdot10^{-7}.
\end{array}
\]
For \(\partial_uE_u\), the corresponding high/full fractions are
\(2.27\cdot10^{-7}\), \(1.20\cdot10^{-6}\), and \(5.54\cdot10^{-7}\).
Compared with the local-tower constants
\[
  2.30\cdot10^9,\qquad 4.41\cdot10^{10},\qquad 4.61\cdot10^{11},
\]
the global trace operator supplies the missing observability.  The theorem to
prove is now:
\[
  E_u^\ast E_u\le C_{\rm window}(u)A,\qquad
  f\in H_M\cap\ker R_{\rm global},
\]
uniformly for \(u\in[.08,.52]\), with the finite low Schur block handled
separately.

A second diagnostic verifies that this is genuinely an interval-trace
phenomenon, rather than a better-tuned local tower.  The local tower at \(s_0\)
is used only to generate the source-heavy high directions; those same
directions are then measured by the sampled global trace operator on
\([.02,.545]\).  The largest source direction in the local high block gives
\[
\begin{array}{c|c|c|c|c|c|c}
\operatorname{basis}&\operatorname{local}&\operatorname{global}&
\operatorname{high}&\lambda_{\rm source}&\|Rf\|_{\ell^2}&
\|Rf\|_{\ell^2}/\|E f\|\\\hline
18&10&11&5&1.828513540\cdot10^{10}&7.797463823\cdot10^3&
5.766393536\cdot10^{-2}\\
20&12&12&6&3.509554015\cdot10^{11}&3.087183713\cdot10^5&
5.211185214\cdot10^{-1}\\
22&14&13&8&3.661655940\cdot10^{12}&1.533555834\cdot10^8&
8.014211977\cdot10^{1}.
\end{array}
\]
At the chosen numerical threshold the source constant on the sampled
global-trace kernel inside this local high block is zero in all three cases.
This supports the compactness form of the missing theorem: an \(A\)-normalized
high-block sequence with \(R_{\rm global}f_n\to0\) cannot retain a nonzero
source row \(E_u f_n\).  The Lagrange identity should convert such a sequence
into a closed-trace solution, and the interval trace-resolution identity should
force that solution to vanish.

The finite algebraic form of this compactness statement was then tested by
forming, on the \(A\)-normalized high block,
\[
  S=E^\ast E,\qquad T=R_{\rm global}^\ast R_{\rm global}.
\]
The raw quotient \(S\le C T\) is extremely ill-conditioned, because \(T\) has
tiny singular directions.  The relevant question is instead whether the kernel
or near-kernel of \(T\) intersects the source-active subspace of \(S\).  With
source-active cutoff \(\lambda_S\ge 10^{-8}\lambda_{S,\max}\), the finite scan
gives
\[
\begin{array}{c|c|c|c|c|c}
\operatorname{basis}&\operatorname{high}&\operatorname{active}&
\lambda_{S,\max}&C_{\rm active}&
\operatorname{ker}_T\operatorname{-source\ frac}\\\hline
18&5&2&1.828513540\cdot10^{10}&2.511450284\cdot10^{12}&0\\
20&6&2&3.509554015\cdot10^{11}&1.283088501\cdot10^{16}&0\\
22&8&3&3.661655940\cdot10^{12}&1.732129406\cdot10^{20}&0.
\end{array}
\]
The full, source-inactive trace kernel at basis \(22\) carries only
\(1.27\cdot10^{-7}\) of the top source mass.  Thus the theorem should be
proved as a qualitative interval-observability statement: for every
\(\delta>0\), if \(P_\delta\) is the spectral projection of \(S\) onto
\(\{\lambda_S\ge\delta\}\), then
\[
  \ker R_{\rm global}\cap \operatorname{Ran}P_\delta=\{0\}
\]
on the continuum high block.  The source-inactive complement is part of the
high-frequency Schur/tail estimate.

The same source-active gap was checked against both the active cutoff and the
sampled trace density.  The scan used trace ratios \(0.50\), \(0.625\), and
\(0.75\), and source-active cutoffs \(10^{-6},10^{-8},10^{-10},10^{-12}\).
The worst values over this grid were
\[
\begin{array}{c|c|c|c|c}
\operatorname{basis}&\operatorname{active\ dims}&
\min \lambda(T|_{\rm active})&
\max {\rm source}(\ker T|_{\rm active})/\lambda_{S,\max}&
\max {\rm source}(\ker T)/\lambda_{S,\max}\\\hline
18&2,3&5.353124349\cdot10^{-10}&0&0\\
20&2,3&2.521576241\cdot10^{-9}&0&0\\
22&2,3&1.692897007\cdot10^{-9}&0&1.272680941\cdot10^{-7}.
\end{array}
\]
Thus the numerical evidence supports the split theorem
\[
\ker R_{\rm global}\cap\operatorname{Ran}P_\delta=\{0\}
\]
for every fixed source-active threshold \(\delta>0\), plus a separate
source-inactive tail estimate
\[
  \|(I-P_\delta)E f\|^2\le \varepsilon(\delta)\langle Af,f\rangle,
  \qquad \varepsilon(\delta)\to0.
\]
The first assertion is the interval observability theorem; the second belongs
to the high-frequency Schur/tail argument.

Equivalently, on each source-active finite block the sampled trace map has full
column rank, so the source-active source row lies in the sampled trace rowspace:
\[
  E_{\rm active}=C_{\rm sample}R_{\rm global}
  \qquad\hbox{on }\operatorname{Ran}P_\delta .
\]
The observed constants for this finite range inclusion are large; for example
the rough bounds \(\lambda_{S,\max}/\min\lambda(T|_{\rm active})\) are
\(3.42\cdot10^{19}\), \(1.39\cdot10^{20}\), and \(2.16\cdot10^{21}\) for
bases \(18,20,22\).  The point is therefore qualitative range inclusion, not a
small Douglas constant.  The continuum proof should establish the exact range
identity by interval trace resolution and leave quantitative smallness to the
source-inactive tail estimate.

For the representative block at basis \(18\), trace ratio \(0.625\), and
source-active cutoff \(10^{-8}\), the finite left-inverse construction gives
active dimension \(2\), trace rank \(2\), and
\[
  {\|E_{\rm active}-C_{\rm sample}R_{\rm global}\|_F\over
   \|E_{\rm active}\|_F}
  =2.26\cdot10^{-69},
  \qquad
  \|C_{\rm sample}\|_F=2.01\cdot10^2.
\]
This confirms the exact finite range identity once the active trace map is
injective.

The last point is important: the closed trace equation alone is not a
uniqueness theorem.  Since the leading coefficient \(e_8(a)\) is nonzero on the
active interval, the condition
\[
  \Lambda_a(f)=0
\]
is a variable-coefficient eighth-order ODE.  It has an eight-dimensional local
solution space.  A direct local jet-tower diagnostic at \(s_0=0.2825\) confirms
that the corrected source rows do not vanish on this formal closed-trace jet
space:
\[
\begin{array}{c|c|c|c}
t&\|E_{\rm loc}\|&p^\ast(t)&\hbox{evaluation projection}\\\hline
.08&1.079233844\cdot10^2&-1.0741740\cdot10^2&1.074174004\cdot10^2\\
.30&1.017823011\cdot10^2&-1.0125435\cdot10^2&1.012543486\cdot10^2\\
.52&9.359523542\cdot10^1&-9.3054977\cdot10^1&9.305497667\cdot10^1.
\end{array}
\]
Therefore the continuum proof must not be phrased as bare ODE uniqueness.
The correct compactness argument is:

\[
\begin{gathered}
f_n\hbox{ \(A\)-bounded and source-active},\quad R_{\rm global}f_n\to0,\\
\hbox{Volterra/Sturm ellipticity}\Rightarrow
  f_n^{(k)}\to f^{(k)}\hbox{ locally for }0\le k\le8,\\
R_{\rm global}f=0,\qquad E f_n\to E f,\\
E_{\rm active}\in\overline{\operatorname{Ran}R_{\rm global}^\ast}
  \Rightarrow E_{\rm active}f=0,
\end{gathered}
\]
contradicting the source-active normalization.  The remaining analytic theorem
is thus closed active range inclusion, with the source-inactive complement
handled by the tail estimate.

\subsection{Trace-to-source representation}

The active range theorem should be proved by an adjoint Green representation.
Let
\[
  P f(a)=\Lambda_a(f)
       = \sum_{k=0}^8 {e_k(a)\over k!} f^{(k)}(a),
  \qquad I=[a_-,a_+].
\]
The exact Lagrange identity is
\[
  {d\over da}B_P[h,f](a)=h(a)P f(a)-f(a)P^\ast h(a),
\]
where
\[
  P^\ast h=\sum_{k=0}^8(-1)^kD_a^k
  \left({e_k(a)\over k!}h(a)\right).
\]
Consequently
\[
  \int_I K_u(a)\Lambda_a(f)\,da
  =\int_I f(a)P^\ast K_u(a)\,da+
    [B_P[K_u,f]]_{a_-}^{a_+}.
\]
Thus a source row \(E_u\) lies in the closed interval trace rowspace once one
constructs \(K_u\) such that
\[
  P^\ast K_u=\eta_u,\qquad
  [B_P[K_u,f]]_{a_-}^{a_+}=\beta_u(f),\qquad
  E_u(f)=\langle\eta_u,f\rangle+\beta_u(f),
\]
with \(K_u\) locally integrable and uniformly bounded in the trace-dual norm,
and with \(\beta_u\) contained in the source-inactive tail.  After applying the
active source projection this gives
\[
  P_\delta E_u(f)
  =P_\delta\int_I K_u(a)\Lambda_a(f)\,da,
\]
hence \(E_{\rm active}\in\overline{\operatorname{Ran}R_{\rm global}^\ast}\).

The finite Galerkin analogue is already exact.  On the representative
source-active high block with basis \(18\), local constraint dimension \(10\),
global trace dimension \(11\), and active dimension \(2\), the sampled trace map
has rank \(2\) on the active block.  Therefore
\[
  C_N=E_N(R_N^\ast R_N)^+R_N^\ast,\qquad E_N=C_NR_N,
\]
and the extracted coefficient profiles satisfy
\[
  {\|E_N-C_NR_N\|_F\over\|E_N\|_F}=2.41\cdot10^{-69},
  \qquad \|C_N\|_{\rm op}=2.014\cdot10^2.
\]
These are raw sampled coefficients \(C_j\), not density values.  They are
smooth-looking on the source window:
\[
\begin{array}{c|cc|cc}
u&\max|C_{\rm bdry}|&{\rm weighted}\ C_{\rm bdry}
 &\max|C_{\rm eval}|&{\rm weighted}\ C_{\rm eval}\\\hline
.08&2.3896206&.7949372&31.733379&10.123421\\
.30&2.4241396&.8041340&29.912683& 9.542592\\
.52&2.4021867&.7947432&27.490414& 8.769852.
\end{array}
\]
The continuum scaling test is stricter.  If
\[
  \sum_j C_j\Lambda_{a_j}(f)\simeq
  \int_IK(a)\Lambda_a(f)\,da,
\]
then \(C_j\simeq w_jK(a_j)\).  Thus
\[
  \|K\|_{L^1}\simeq\sum_j|C_j|,\qquad
  \|K\|_{L^2}\simeq\left(\sum_j{|C_j|^2\over w_j}\right)^{1/2}.
\]
The refinement scan also computes the correct weighted minimal \(L^2\) density
using the trace Gram \(R^\ast W R\).  It gives
\[
\begin{array}{c|c|c|c|c|c}
\#a_j&{\rm rank}&{\rm rel.\ residual}&\|C\|_{\rm op}
 &{\rm raw}\ \max\|K\|_{L^2}
 &{\rm weighted}\ \max\|K\|_{L^2}\\\hline
7&2&3.68\cdot10^{-69}&247.54815&342.70544&330.80492\\
9&2&1.44\cdot10^{-69}&220.85732&347.51287&334.66544\\
11&2&2.41\cdot10^{-69}&201.41473&349.51484&336.58072\\
13&2&1.72\cdot10^{-69}&186.38790&350.26874&337.65662.
\end{array}
\]
The weighted \(L^1\) density norms are also stable, about \(219\)--\(223\).
Thus the integral density norms stay essentially bounded as the trace mesh is
refined.
This is the sampled version of
\[
  E_{\rm active}f=\int_I K(a)^\ast\Lambda_a(f)\,da.
\]
The remaining continuum passage is to prove these bounds analytically.
Precisely, for trace meshes with fill distance \(h_N\to0\), write
\(K_N(a_j)=C_{N,j}/w_j\).  If
\[
  \sup_N\|K_N\|_{L^2(I)}<\infty,
\]
then a subsequence converges weakly in \(L^2(I)\).  Since
\(a\mapsto\Lambda_a(f)\) is continuous on the Volterra/Sturm core, quadrature
convergence yields
\[
  E_{\rm active}(f)=\int_I K(a)^\ast\Lambda_a(f)\,da,
\]
and hence
\[
  E_{\rm active}\in
  \overline{\operatorname{Ran}R_{\rm global}^{\ast}}.
\]
Thus the hard analytic estimate is now the uniform quadrature-scaled bound
\[
  \sup_N \|C_N/w_N\|_{L^2(I)}<\infty,
\]
with the source-inactive endpoint remainder absorbed by the separated tail
estimate.

\subsection{Weighted frame form}

The preceding bound is equivalent to an interval trace frame inequality.  In a
finite Galerkin section let
\[
  T_N=W_N^{1/2}R_N,\qquad E_N=E_{\rm active,N},
\]
where \(W_N\) contains the quadrature weights on \(I\).  Then
\[
  \|T_Nv\|^2=\sum_jw_j|\Lambda_{a_j}(v)|^2
  \simeq \int_I|\Lambda_a(v)|^2\,da.
\]
The minimal discrete \(L^2\) density satisfying
\[
  E_N=Y_NT_N
\]
is \(Y_N=E_NT_N^+\).  Hence
\[
  \|Y_N\|\le {\|E_N\|\over s_{\min}(T_N)}
  ={\|E_N\|\over\sqrt{\lambda_{\min}(T_N^\ast T_N)}}.
\]
Thus it is enough to prove an active interval observability inequality
\[
  \int_I|\Lambda_a(v)|^2\,da\ge
  \gamma_\delta\|v\|_A^2
\]
on the source-active high block, together with the existing source row bound.

The first weighted-frame basis scan gives
\[
\begin{array}{c|c|c|c|c|c}
{\rm basis}&\#a_j&\dim&\lambda_{\min}(T^\ast T)&
\|Y_N\|&\max{\rm row}\ L^2\\\hline
18&9&2&299.09458& 946.72705&334.66544\\
18&13&2&293.11515& 955.18580&337.65662\\
20&9&2&7525.0056&1683.3148&593.35392\\
20&13&2&7489.9584&1687.1679&594.71218\\
22&9&3&2.8749823\cdot10^6&317.01515&111.54097.
\end{array}
\]
The lower frame bound is stable under trace refinement and improves in the
larger Galerkin sections in this normalization.  The finite theorem is
immediate: if
\[
  \lambda_{\min}(T_N^\ast T_N|_{H_{N,\delta}})=\gamma_N>0,
\]
then for \(v\in H_{N,\delta}\),
\[
  \sum_jw_j|\Lambda_{a_j}(v)|^2
  =\|T_Nv\|^2
  \ge \gamma_N\|v\|_A^2.
\]
The remaining theorem is the continuum passage: \(S=E^\ast E\) is compact,
\(H_\delta=\operatorname{Ran}{\bf 1}_{[\delta,\infty)}(S)\) is finite
dimensional, \(H_{N,\delta}\to H_\delta\) in the \(A\)-graph norm, the trace
quadratures converge uniformly on this finite-dimensional space, and
\(\liminf_N\gamma_N>0\).  These hypotheses imply
\[
  \int_I|\Lambda_a(v)|^2\,da\ge\gamma_\delta\|v\|_A^2,
  \qquad v\in H_\delta.
\]
Indeed, for \(v\in H_\delta\), graph convergence gives \(v_N\in H_{N,\delta}\)
with \(v_N\to v\) in the \(A\)-norm.  The finite frame inequality gives
\[
  \|W_N^{1/2}R_Nv_N\|^2\ge\gamma_N\|v_N\|_A^2.
\]
Uniform quadrature convergence gives
\[
  \|W_N^{1/2}R_Nv_N\|^2
  \to \int_I|\Lambda_a(v)|^2\,da,
\]
while \(\|v_N\|_A\to\|v\|_A\).  Taking lower limits and using
\(\liminf_N\gamma_N\ge\gamma_\delta\) proves the displayed continuum
observability inequality.  By the closed range theorem, this is equivalent to
the active source rows lying in
\(\overline{\operatorname{Ran}R_{\rm global}^\ast}\).

The current finite certificate over bases \(18,20,22\) has observed
\[
  \min\gamma_N=293.115151009,
  \qquad
  \max{\rm residual}=1.093\cdot10^{-63}.
\]

Two of the hypotheses are standard.  First, if \(S=E^\ast E\) is compact
self-adjoint on the \(A\)-Hilbert space, \(P_N\to I\) in the \(A\)-graph norm,
and
\[
  \|P_NSP_N-S\|_{A\to A}\to0,
\]
then, for every \(\delta>0\) outside the spectrum of \(S\), the Riesz spectral
projections
\[
  {\bf 1}_{[\delta,\infty)}(P_NSP_N)
  \to {\bf 1}_{[\delta,\infty)}(S)
\]
converge in operator norm.  This gives \(H_{N,\delta}\to H_\delta\) in the
\(A\)-graph norm.

Second, since \(H_\delta\) is finite-dimensional, trace quadrature convergence
is uniform once \(a\mapsto\Lambda_a(v)\) is continuous from the \(A\)-graph norm
to \(C(I)\).  Indeed, the \(A\)-unit sphere of \(H_\delta\) is compact, so the
family \(|\Lambda_a(v)|^2\) is compact in \(C(I)\); positive quadrature rules
with mesh tending to zero converge uniformly on this family.  The finite
diagnostic from the current weighted-frame scans records only small quadrature
drift:
\[
\begin{array}{c|cc}
{\rm basis}&\Delta\lambda_{\min}/\lambda_{\min}&
\Delta(\max{\rm row}\ L^2)/(\max{\rm row}\ L^2)\\\hline
18&-1.9992\%&+0.8938\%\\
20&-0.4657\%&+0.2289\%.
\end{array}
\]
Thus the remaining nonstandard analytic inputs are compact/norm convergence of
\(S=E^\ast E\) and the positive lower-frame bound
\(\liminf_N\gamma_N>0\).

The compactness of \(S\) reduces to endpoint elliptic regularity.  The corrected
source row has the finite-jet form
\[
  E_u f=
  \left(\sum_{j=0}^7 b_j(u)f^{(j)}(s_0),\ p(u)f(s_0)\right),
\]
where \(b_j\) and \(p\) are smooth on the compact source interval
\([.08,.52]\).  If the commuted Sturm estimate gives
\[
  \|f\|_{H^m(J)}\le C\|f\|_A,\qquad m>8+\frac12,
\]
on an interval \(J\ni s_0\), then the endpoint jet maps
\(f\mapsto f^{(j)}(s_0)\) are bounded on the \(A\)-space.  Hence
\[
  E:H_A\to L^2([.08,.52];\mathbb C^2)
\]
has finite-dimensional range contained in the span of
\((b_j,0)\), \(0\le j\le7\), and \((0,p)\).  Thus \(E\) is compact and
\(S=E^\ast E\) is compact, self-adjoint, and nonnegative.  The same finite-rank
factorization gives \(\|P_NE^\ast EP_N-E^\ast E\|_{A\to A}\to0\) for Galerkin
projections converging in the graph norm.  Therefore compactness and spectral
active-space convergence follow from the commuted Sturm elliptic estimate.  The
remaining hard theorem is the injective lower-frame statement
\[
  R_{\rm global}v=0,\quad v\in H_\delta \quad\Longrightarrow\quad v=0,
\]
equivalently \(\liminf_N\gamma_N>0\).

On the finite active space this can be sharpened to a determinant condition.
For an \(A\)-orthonormal active basis \(v_1,\ldots,v_d\), set
\[
  F_k(a)=\Lambda_a(v_k).
\]
If for some \(a_1,\ldots,a_d\in I\),
\[
  \det(F_k(a_i))_{i,k=1}^d\ne0,
\]
then the sampled trace map is injective on the active space.  Since the high
block is built to satisfy the local trace tower at \(s_0\), the determinant
certificate should use off-base trace points rather than the confluent
Wronskian at \(s_0\).  The basis-22 certificate has \(d=3\) and a nonzero
normalized minor
\[
  |\det(F_k(a_i))|=4.3555051\cdot10^{-7},
  \qquad
  (a_1,a_2,a_3)=(.085625,.348125,.545).
\]
Thus the analytic lower-frame theorem may be attacked as a Chebyshev or
unique-continuation theorem for the continuum trace-response family.

The cross-basis determinant scan gives
\[
\begin{array}{c|c|c|c|c}
{\rm basis}&\#a_j&d&\max|\det|&{\rm points}\\\hline
18&9&2&2.4856496\cdot10^{-2}&.085625,.348125\\
18&13&2&1.8740173\cdot10^{-2}&.06375,.37\\
20&9&2&2.2598453\cdot10^{-3}&.02,.2825\\
20&13&2&1.7398684\cdot10^{-3}&.06375,.2825\\
22&9&3&4.3555051\cdot10^{-7}&.085625,.348125,.545.
\end{array}
\]
The continuum target can be stated invariantly: for \(d=\dim H_\delta\), prove
there is \(a=(a_1,\ldots,a_d)\in I^d\) such that the exterior trace functional
\[
  v_1\wedge\cdots\wedge v_d
  \mapsto \det(\Lambda_{a_i}(v_k))
\]
is nonzero on \(\wedge^dH_\delta\).  Since \(\wedge^dH_\delta\) is
one-dimensional, this is precisely the nonvanishing scalar determinant needed
for injectivity of \(R_{\rm global}\) on \(H_\delta\).

A landscape scan at basis \(22\), excluding a neighborhood of \(s_0\), gives
56 off-base \(3\times3\) minors: 53 positive, 3 negative, and none near zero.
The largest values are
\[
\begin{array}{c|c|c}
{\rm rank}&|\det|&{\rm points}\\\hline
1&4.3555\cdot10^{-7}&.085625,.348125,.545\\
2&4.2511\cdot10^{-7}&.02,.085625,.348125\\
3&4.0280\cdot10^{-7}&.085625,.348125,.479375\\
4&3.7479\cdot10^{-7}&.02,.15125,.348125\\
5&3.6905\cdot10^{-7}&.02,.348125,.545.
\end{array}
\]
Thus the determinant evidence is not tied to one sampled tuple.  The analytic
route is now: prove the trace responses \(F_v(a)=\Lambda_a(v)\) are real
analytic on \(I\), prove the unique-continuation implication
\[
  F_v=0\ {\rm on\ a\ nonempty\ interval}\quad\Longrightarrow\quad v=0,
  \qquad v\in H_\delta,
\]
and then conclude exterior nonvanishing.  If all \(d\)-fold evaluation
determinants vanished identically, the analytic response space would have rank
less than \(d\) on \(I\), giving a nonzero active response vanishing on an open
interval, contrary to unique continuation.

The unique-continuation implication has a local derivative-rank certificate.
If \(F_v(a)=\Lambda_a(v)\) vanishes on an interval containing an off-base point
\(a_0\), then all derivatives \(D_a^qF_v(a_0)\) vanish.  Hence it is enough to
find derivative orders \(q_1,\ldots,q_d\) such that
\[
  \det\left(D_a^{q_i}\Lambda_a(v_k)|_{a=a_0}\right)\ne0.
\]
Using exact confluent trace derivatives at the off-base point, the loose
relative cutoff \(10^{-8}\) gives a three-dimensional basis-22 certificate.
However basis 24 returns to two active modes, so the three-dimensional
basis-22 determinant is not yet a stable continuum target.  With the stricter
source-active window \(10^{-6}\), the active dimension is stable and the
cross-basis scan gives
\[
\begin{array}{c|c|c|c|c|c}
N&d_N&a_0&{\rm rank}&\sigma_{\min}&{\rm orders}\\\hline
18&2&.545&2&5.5267846\cdot10^{-1}&7,8\\
20&2&.545&2&2.8529276\cdot10^{-3}&7,8\\
22&2&.348125&2&1.0107808\cdot10^{-1}&6,8\\
24&2&.545&2&6.4517532\cdot10^{-2}&7,8.
\end{array}
\]
The weakest normalized singular margin in this stable window is
\(2.8529276\cdot10^{-3}\).  Thus the continuum determinant target is the
two-dimensional active band, not the unstable three-dimensional cutoff
artifact.

The finite unique-continuation implication is immediate from this certificate.
Let \(H_{\delta,N}\) be a \(d_N\)-dimensional active space with
\(A\)-orthonormal basis \(v_1,\ldots,v_{d_N}\), and put
\[
  F_k(a)=\Lambda_a(v_k).
\]
If an off-base point \(a_0\) and orders \(q_1,\ldots,q_{d_N}\) satisfy
\[
  \det\bigl(D_a^{q_i}F_k(a_0)\bigr)_{i,k=1}^{d_N}\ne0,
\]
then \(F_v(a)=\Lambda_a(v)\) cannot vanish on any open interval containing
\(a_0\) unless \(v=0\).  Indeed, interval vanishing forces all derivatives at
\(a_0\) to vanish, and the displayed full-rank matrix kills the coefficient
vector of \(v\).

The cross-basis derivative-rank scan gives the following exact confluent-row
certificates:
\[
\begin{array}{c|c|c|c|c|c}
N&d_N&a_0&{\rm rank}&\sigma_{\min}&{\rm orders}\\\hline
18&2&.545&2&5.5267846\cdot10^{-1}&7,8\\
20&2&.545&2&2.8529276\cdot10^{-3}&7,8\\
22&2&.348125&2&1.0107808\cdot10^{-1}&6,8\\
24&2&.545&2&6.4517532\cdot10^{-2}&7,8.
\end{array}
\]
Thus the finite active unique-continuation statement has been proved on the
computed active blocks.  The continuum theorem is reduced to showing that the
off-base derivative functionals \(D_a^q\Lambda_a|_{a=a_0}\) converge on
\(H_{\delta,N}\to H_\delta\) in the \(A\)-graph norm and that at least one of
these confluent determinants has a nonzero limiting value.  Once this holds,
analyticity of \(a\mapsto\Lambda_a(v)\) gives
\[
  \Lambda_a(v)=0\ {\rm on\ a\ nonempty\ open\ interval}
  \quad\Longrightarrow\quad v=0,\qquad v\in H_\delta,
\]
and the exterior determinant nonvanishing follows.

We record the abstract convergence lemma.  Let \(H_A\) be the closed high-block
Hilbert space with \(A\)-inner product, let \(S=E^\ast E\), and assume
\(\delta\) lies in a spectral gap of the compact nonnegative operator \(S\).
Put
\[
  H_\delta={\rm Ran}\,\mathbf 1_{[\delta,\infty)}(S).
\]
Let \(P_N\to I\) in the \(A\)-graph sense and \(S_N=P_NSP_N\), so that
\(\|S_N-S\|_{A\to A}\to0\).  Then the Riesz projections
\[
  \Pi_\delta={1\over 2\pi i}\int_\Gamma (z-S)^{-1}\,dz,\qquad
  \Pi_{\delta,N}={1\over 2\pi i}\int_\Gamma (z-S_N)^{-1}\,dz
\]
satisfy \(\|\Pi_{\delta,N}-\Pi_\delta\|_{A\to A}\to0\), where \(\Gamma\)
separates the spectrum above \(\delta\).  For \(N\) large the dimensions agree,
and the polar graph map
\[
  U_N=\Pi_{\delta,N}\Pi_\delta
      \bigl(\Pi_\delta\Pi_{\delta,N}\Pi_\delta\bigr)^{-1/2}
\]
maps \(H_\delta\) isomorphically onto \(H_{\delta,N}\) with
\(\|U_N-I\|_{A\to A,H_\delta}\to0\).

Suppose the commuted Sturm estimate gives bounded derivative rows
\[
  \ell_q(f)=D_a^q\Lambda_a(f)|_{a=a_0},\qquad
  |\ell_q(f)|\le C_q\|f\|_A .
\]
Then, uniformly for \(\|v\|_A=1\) in \(H_\delta\),
\[
  |\ell_q(U_Nv)-\ell_q(v)|
  \le C_q\|U_Nv-v\|_A\to0.
\]
Thus the off-base derivative rows restricted to \(H_{\delta,N}\) converge in
operator norm, after identifying \(H_{\delta,N}\) with \(H_\delta\) by \(U_N\).
For fixed orders \(q_1,\ldots,q_d\) and an \(A\)-orthonormal basis
\(v_1,\ldots,v_d\) of \(H_\delta\), the matrices
\[
  M_N=\bigl(\ell_{q_i}(U_Nv_k)\bigr)_{i,k}
\]
converge to \(M=(\ell_{q_i}(v_k))_{i,k}\), hence
\[
  \det M_N\to \det M.
\]
Therefore a uniform determinant gap,
\[
  \liminf_N |\det M_N|>0,
\]
implies \(\det M\ne0\), and hence active unique continuation.  The remaining
non-formal step is the determinant-gap estimate for one off-base confluent
tuple in the stable two-dimensional band, such as \(a_0=.545\) and
\(q=(7,8)\).

In the two-dimensional band this can be stated as a singular-value gap.  Let
\[
  m_N=\sigma_{\min}(M_N),\qquad
  C=\left(\sum_i\|\ell_{q_i}\|_{A^\ast}^2\right)^{1/2},\qquad
  \alpha_N=\|U_N-I\|_{A\to A,H_\delta}.
\]
Then \(\|M_N-M\|_2\le C\alpha_N\), and Weyl's inequality gives
\[
  \sigma_{\min}(M)\ge m_N-C\alpha_N.
\]
Thus the continuum determinant is nonzero as soon as
\[
  C\alpha_N<m_N.
\]
The value \(2.8529276\cdot10^{-3}\) reported by the derivative-rank scan is a
column-normalized singular value.  It is an angular certificate, but it is not
the raw perturbation margin in the inequality above.  Computing the raw
high-block row norms for the fixed stable tuple \(a_0=.545,\ q=(7,8)\) gives
\[
\begin{array}{c|c|c|c|c}
N&m_N&C_{\rm high}&m_N/C_{\rm high}&m_N^{\rm norm}\\\hline
18&5.6675\cdot10^{12}&1.4894\cdot10^{13}&3.8052\cdot10^{-1}&5.5268\cdot10^{-1}\\
20&4.4697\cdot10^{12}&3.8648\cdot10^{15}&1.1565\cdot10^{-3}&2.8529\cdot10^{-3}\\
22&1.2076\cdot10^{15}&5.3633\cdot10^{17}&2.2516\cdot10^{-3}&3.1971\cdot10^{-3}\\
24&3.6226\cdot10^{18}&5.6348\cdot10^{21}&6.4290\cdot10^{-4}&6.4517\cdot10^{-2}.
\end{array}
\]
Thus the raw determinant-gap theorem would require
\[
  \alpha_N<6.4289543\cdot10^{-4},
\]
not merely \(2.8529276\cdot10^{-3}\).  A parent-space Galerkin drift diagnostic
for consecutive active spaces gives polar-alpha estimates \(1.3878487\),
\(1.4117373\), and \(1.4141422\), far above this threshold.  Hence the simple
finite-to-continuum shortcut does not prove the determinant gap.  The next
valid theorem must either prove a much sharper continuum projection estimate in
a fixed active spectral band, or reformulate the determinant argument in
column-normalized projective coordinates with explicit column-scale stability.

We now state the projective version.  Let
\[
  L=(\ell_{q_1},\ldots,\ell_{q_d})^T,\qquad
  \ell_q(f)=D_a^q\Lambda_a(f)|_{a=a_0}.
\]
Assume the active source eigenvalues are simple and separated,
\[
  \mu_1>\cdots>\mu_d>\delta>\mu_{d+1},
\]
and let \(v_1,\ldots,v_d\) be the corresponding \(A\)-orthonormal eigenvectors
of \(S=E^\ast E\).  Define
\[
  x_j=Lv_j,\qquad d_j=\|x_j\|,\qquad
  \widehat M=[x_1/d_1\ \cdots\ x_d/d_d].
\]
The Galerkin objects \(\widehat M_N\) are defined in the same way using the
Riesz/polar eigenvectors \(v_{j,N}\).  Since \(S_N\to S\) in norm and the
eigenvalues are simple, \(v_{j,N}\to v_j\) in the \(A\)-norm.  Boundedness of
the rows \(\ell_q\) gives \(Lv_{j,N}\to Lv_j\).  If \(d_\ast=\min_jd_j>0\),
then column normalization is continuous and
\[
  \widehat M_N\to \widehat M.
\]
Hence
\[
  \liminf_N\sigma_{\min}(\widehat M_N)>0
  \quad\Longrightarrow\quad
  \sigma_{\min}(\widehat M)>0.
\]
This is enough for active unique continuation, because interval vanishing of
\(\Lambda_a(v)\) forces the selected derivatives at \(a_0\) to vanish, hence
\(\widehat M c=0\), and full rank gives \(c=0\).

Quantitatively, the elementary bound
\[
  \left\|{x\over \|x\|}-{y\over \|y\|}\right\|
  \le {2\|x-y\|\over \|x\|}
\]
for \(\|x-y\|\le \|x\|/2\) yields the sufficient condition
\[
  {2\sqrt d\,C\alpha_N\over d_{\min,N}}<\eta_N,\qquad
  \eta_N=\sigma_{\min}(\widehat M_N),
\]
or
\[
  \alpha_N < {\eta_N d_{\min,N}\over 2\sqrt d\,C}.
\]
For the fixed tuple \(a_0=.545,\ q=(7,8)\), the projective thresholds are
\[
\begin{array}{c|c|c|c|c}
N&\eta_N&d_{\min,N}/C&{\rm scale\ ratio}&{\rm allowed\ }\alpha_N\\\hline
18&5.5268\cdot10^{-1}&6.4203\cdot10^{-1}&1.17045&1.2545\cdot10^{-1}\\
20&2.8529\cdot10^{-3}&3.0060\cdot10^{-1}&3.16632&3.0321\cdot10^{-4}\\
22&3.1971\cdot10^{-3}&6.7692\cdot10^{-1}&1.08612&7.6515\cdot10^{-4}\\
24&6.4517\cdot10^{-2}&7.0536\cdot10^{-3}&141.768&1.60895\cdot10^{-4}.
\end{array}
\]
The alternate tuple \(a_0=.348125,\ q=(6,8)\) has minimum projective allowance
only \(3.51008\cdot10^{-7}\), so it is not the stable target.  The next
analytic problem is therefore to prove direct convergence of the normalized
response columns for \(a_0=.545,\ q=(7,8)\), preferably without passing through
the large global row norm \(C\).

This direct convergence is qualitative once the limiting response columns are
nonzero.  In the stable two-dimensional band, take
\[
  Lf=\left(D_a^7\Lambda_a(f)|_{a=.545},
           D_a^8\Lambda_a(f)|_{a=.545}\right).
\]
The commuted Sturm trace theorem makes \(L:H_A\to\mathbb R^2\) continuous.  If
the finite eigenvectors are oriented by the Riesz/polar graph map, then
\[
  v_{j,N}\to v_j\quad{\rm in}\ H_A,
\]
hence \(Lv_{j,N}\to Lv_j\).  If \(Lv_j\ne0\), continuity of
\(x\mapsto x/\|x\|\) away from the origin gives
\[
  {Lv_{j,N}\over \|Lv_{j,N}\|}
  \longrightarrow
  {Lv_j\over \|Lv_j\|}.
\]
Thus \(\widehat M_N\to\widehat M\) directly, without a global \(C\alpha_N\)
determinant estimate.  The remaining non-formal statement is the noncollapse
condition: \(Lv_1,Lv_2\) must both be nonzero and must determine distinct
projective directions.

The finite response-column diagnostic gives, after best sign/permutation
alignment,
\[
\begin{array}{c|c|c}
{\rm bases}&{\rm max\ column\ distance}&{\rm Frobenius\ distance}\\\hline
18\to20&7.3584\cdot10^{-1}&7.3724\cdot10^{-1}\\
20\to22&8.6466\cdot10^{-2}&1.2194\cdot10^{-1}\\
22\to24&1.4195\cdot10^{-1}&1.5235\cdot10^{-1}.
\end{array}
\]
These distances are far better scaled than the parent-space drift, but they
are not yet a numerical Cauchy proof.  The next analytic target is therefore:
prove the noncollapse of the limiting response vectors for
\(a_0=.545,\ q=(7,8)\), or equivalently prove
\(\sigma_{\min}(\widehat M)>0\) directly.

The stronger domination
\[
  \sup_{f\in\ker L,\ \|f\|_A=1}\langle Sf,f\rangle < \delta
\]
is false in the finite model: for bases \(18,20,22,24\), the source mass on
\(\ker L\), divided by the top source eigenvalue, is approximately
\[
  1.10\cdot10^{-3},\quad
  1.63\cdot10^{-3},\quad
  1.46\cdot10^{-1},\quad
  1.18\cdot10^{-4},
\]
all above the \(10^{-6}\) active cutoff.  Thus noncollapse must use the source
eigenline equation, not the whole kernel of \(L\).

Let \(E:H_A\to Y\) be the source map, so \(S=E^\ast E\).  If
\[
  Sv=\mu v,\qquad \mu>0,
\]
then \(u=\mu^{-1/2}Ev\) satisfies
\[
  EE^\ast u=\mu u,\qquad v=\mu^{-1/2}E^\ast u.
\]
Therefore
\[
  Lv=\mu^{-1/2}LE^\ast u.
\]
Consequently the active noncollapse theorem is equivalent to the source-side
rank statement
\[
  {\rm rank}\bigl(LE^\ast|_{U_\delta}\bigr)=2,
\]
where \(U_\delta\) is the two-dimensional active eigenspace of \(EE^\ast\).
The finite source-side certificate gives
\[
\begin{array}{c|c|c}
N&\sigma_{\min}(\widehat M_N)&
\sigma_{\min}\bigl(\widehat{LE^\ast|_{U_{\delta,N}}}\bigr)\\\hline
18&5.5267804\cdot10^{-1}&5.5267804\cdot10^{-1}\\
20&2.8529266\cdot10^{-3}&2.8529266\cdot10^{-3}\\
22&3.1970980\cdot10^{-3}&3.1970980\cdot10^{-3}\\
24&6.4517300\cdot10^{-2}&6.4517300\cdot10^{-2}.
\end{array}
\]
This proves finite noncollapse on the computed active blocks and identifies the
continuum theorem to prove: \(LE^\ast\) has rank two on the active source-side
eigenspace.

A source-side stability diagnostic compares \(EE^\ast\), its active spectral
subspace, and the normalized \(LE^\ast\) response columns on the fixed source
sample space.  For \(a_0=.545,\ q=(7,8)\) and active cutoff \(10^{-6}\), the
active dimension remains two for \(N=18,20,22,24\).  The consecutive projection
gaps of the active source subspaces are
\[
  2.8271\cdot10^{-3},\qquad
  1.4971\cdot10^{-2},\qquad
  8.2554\cdot10^{-3}.
\]
The corresponding maximum distances between aligned normalized response
columns are
\[
  7.3584\cdot10^{-1},\qquad
  8.6466\cdot10^{-2},\qquad
  1.4195\cdot10^{-1}.
\]
By contrast, the full top-normalized source operators do not converge in this
finite window; the final operator difference is \(9.8720\cdot10^{-1}\).  Hence
the continuum theorem should not require norm convergence of the whole
normalized source operator.  The right statement is stability of the active
Riesz projector together with source-side noncollapse:
\[
  {\rm rank}\bigl(LE^\ast|_{U_\delta}\bigr)=2.
\]

The limiting passage is a Riesz-projection theorem.  Suppose the source maps
\(E_N:H_A\to Y\) converge to \(E\) in the topology supplied by the source-window
construction, and put \(T_N=E_NE_N^\ast\), \(T=EE^\ast\).  If \(\delta\) lies in
a spectral gap of \(T\) and \(U_\delta={\rm Ran}\,\mathbf 1_{[\delta,\infty)}(T)\)
is two-dimensional, then the active source projectors
\[
  Q_N=\mathbf 1_{[\delta,\infty)}(T_N)
\]
converge to \(Q=\mathbf 1_{[\delta,\infty)}(T)\).  Since the commuted Sturm
trace theorem makes \(L:H_A\to\mathbb R^2\) continuous,
\[
  LE_N^\ast Q_N\longrightarrow LE^\ast Q.
\]
Therefore the nonzero singular values of \(LE_N^\ast|_{Q_NY}\) converge to
those of \(LE^\ast|_{QY}\).  The remaining analytic theorem is the rank-two
source-side noncollapse \( {\rm rank}(LE^\ast|_{U_\delta})=2\).

For a quantitative version, let
\[
  m_N=\sigma_{\min}(LE_N^\ast Q_N),\qquad
  b_N=\|LE_N^\ast\|,\qquad
  g_N={\rm dist}\bigl(\sigma(T_N|_{Q_NY}),\sigma(T_N|_{Q_N^\perp Y})\bigr).
\]
If
\[
  \|T-T_N\|<g_N/4,\qquad
  \|LE^\ast-LE_N^\ast\|+b_N{4\|T-T_N\|\over g_N}<m_N,
\]
then \(LE^\ast\) has rank two on the limiting active subspace.  Indeed, the
first inequality gives the Riesz-projector estimate
\(\|Q-Q_N\|\le4\|T-T_N\|/g_N\), and therefore
\[
  \|LE^\ast Q-LE_N^\ast Q_N\|
  \le \|LE^\ast-LE_N^\ast\|+b_N\|Q-Q_N\|<m_N.
\]
Weyl's inequality preserves the positive second singular value.

The computed split margins for \(a_0=.545,\ q=(7,8)\) are
\[
\begin{array}{c|c|c|c|c|c}
N&\eta_N&m_N/b_N&g_N/\|T_N\|&
\|B-B_N\|/b_N&\|T-T_N\|/\|T_N\|\\\hline
18&5.5268\cdot10^{-1}&2.9635\cdot10^{-3}&1.2385\cdot10^{-5}&
7.4086\cdot10^{-4}&2.2938\cdot10^{-9}\\
20&2.8529\cdot10^{-3}&6.1325\cdot10^{-6}&2.3162\cdot10^{-5}&
1.5331\cdot10^{-6}&8.8776\cdot10^{-12}\\
22&3.1971\cdot10^{-3}&1.0718\cdot10^{-3}&7.4990\cdot10^{-2}&
2.6796\cdot10^{-4}&5.0236\cdot10^{-6}\\
24&6.4517\cdot10^{-2}&3.0011\cdot10^{-5}&2.1789\cdot10^{-3}&
7.5028\cdot10^{-6}&4.0870\cdot10^{-9}.
\end{array}
\]
Thus \(N=22\) is the best perturbative anchor: its active spectral gap makes
the allowable source-operator error \(5.0\cdot10^{-6}\) in top-normalized
units.  Proving the two displayed approximation bounds for this finite
certificate would complete source-side noncollapse.

This certificate also survives replacing the unweighted source sample by a
trapezoid-weighted discretization of
\[
  Y=L^2([.08,.52];\mathbb R^2).
\]
For \(N=22\), weighted source grids of sizes \(5,7,9,11,13,17\) all have
active dimension two.  The limiting constants stabilize to
\[
  \eta_N=3.1971\cdot10^{-3},\qquad
  m_N/b_N=1.0719\cdot10^{-3},\qquad
  g_N/\|T_N\|=7.5001\cdot10^{-2},
\]
with split tolerances
\[
  \|B-B_N\|/b_N=2.6797\cdot10^{-4},\qquad
  \|T-T_N\|/\|T_N\|=5.0246\cdot10^{-6}.
\]
Thus the source-side rank margin is stable in the continuum source variable.
The remaining step is a deterministic quadrature or smooth-kernel bound for
the weighted source operators below these two tolerances.

The point-quadrature formulation should be made through the finite right Gram
operator
\[
  S=E^\ast E=\int_{.08}^{.52}E(u)^\ast E(u)\,du.
\]
The nonzero spectra of \(S\) and \(T=EE^\ast\) agree, and
\[
  BB^\ast=(LE^\ast)(LE^\ast)^\ast=L S L^\ast .
\]
Thus it is enough to control the quadrature error for \(S\), then use the fixed
row map \(L\) on the active Riesz subspace of \(S\).  For the composite
trapezoid rule with mesh \(h\),
\[
  \|S-S_h\|
  \le {(.52-.08)h^2\over 12}
      \sup_{u\in[.08,.52]}\left\|
      {d^2\over du^2}\bigl(E(u)^\ast E(u)\bigr)\right\|.
\]
Differentiating the Green-source formula for \(E(u)\) gives an explicit
second-derivative envelope.  With source grid \(65\), the resulting bound is
\[
  \|S_h\|=1.792385429\cdot10^{11},\qquad
  \|S-S_h\|\le8.361215983\cdot10^5,
\]
so
\[
  {\|S-S_h\|\over\|S_h\|}\le4.664853802\cdot10^{-6}
  <5.0245931\cdot10^{-6}.
\]
The weighted grid-65 rank certificate is
\[
\begin{array}{c|c|c|c|c}
\eta_N&m_N/b_N&g_N/\|S_N\|&
\|L{\rm\ error}\|/\|L\|&\|S-S_N\|/\|S_N\|\\\hline
3.197088\cdot10^{-3}&1.0719017\cdot10^{-3}&
7.5000804\cdot10^{-2}&2.6797543\cdot10^{-4}&5.0245931\cdot10^{-6}.
\end{array}
\]
Consequently the source-projector side closes at the deterministic
finite-dimensional level.  The old literal \(B-B_N\) target is not the right
same-codomain object for point quadrature; what is invariantly controlled is
\(BB^\ast=L S L^\ast\), whose relative quadrature error at grid \(65\) is
\[
  4.5562484\cdot10^{-6}.
\]
The sampled derivative envelope can be replaced by a Chebyshev-tail envelope.
Raw interval arithmetic on the Green-source formula is rigorous but too loose:
with source grid \(65\) and \(16\) source subintervals it gives the unusable
relative bound \(9.13\cdot10^{-1}\).  The analytic Chebyshev route is much
sharper.  For source grid \(129\), Chebyshev degree \(32\), and the entries of
\[
  A(u)={d^2\over du^2}\bigl(E(u)^\ast E(u)\bigr),
\]
the coefficient envelope gives
\[
  \sup_u\|A(u)\|\le4.90928918512\cdot10^{11}.
\]
The largest geometric tail ratio in the final coefficient window is
\[
  1.61482847834\cdot10^{-1},
\]
and the resulting trapezoid bound is
\[
  {\|S-S_h\|\over\|S_h\|}\le1.18670850437\cdot10^{-6}
  <5.03\cdot10^{-6}.
\]
Thus the source-projector quadrature estimate closes with margin.  The last
machine-rigorous polish is to perform the same Chebyshev coefficient and tail
estimate in ball arithmetic, or derive it from a Bernstein ellipse bound for
the explicit Green-source formula.

Both checks can be isolated.  First, the degree \(32\), grid \(129\)
coefficient calculation was rerun with interval balls around each sampled
entry and interval arithmetic for the DCT-I coefficients and final-window tail
ratio.  With relative sample radius \(10^{-55}\) the ball certificate gives
\[
\begin{aligned}
  \sup_u\|A(u)\|_{\rm ball}&\le4.90928918512\cdot10^{11},\\
  q_{\rm tail}&\le1.61482847834\cdot10^{-1},\\
  {\|S-S_h\|\over\|S_h\|}&\le1.18670850437\cdot10^{-6}
  <5.03\cdot10^{-6}.
\end{aligned}
\]
The largest coefficient interval width in this run is
\(3.58\cdot10^{-44}\), and no tail-ratio interval crosses the failure
threshold \(1\).

Second, the same tail decay follows from the standard Bernstein ellipse
estimate.  Write \(u=.30+.22x\).  If \(M_\rho\) bounds
\(\|A(.30+.22z)\|\) on the Bernstein ellipse \(E_\rho\), then
\[
  \|A_n\|\le 2M_\rho\rho^{-n},\qquad
  \sum_{n>N}\|A_n\|
  \le {2M_\rho\rho^{-(N+1)}\over 1-\rho^{-1}}.
\]
The apparent divided-difference formula pole \(u=0\) maps to
\[
  x_0=-1.3636363636\ldots,\qquad \rho_0=2.29073082065\ldots,
\]
and is removable analytically; choosing \(\rho=2\) is therefore conservative
for the explicit formula.  A sampled ellipse check with \(96\) points gives
\[
  M_2\simeq5.31292102718\cdot10^{11},\qquad
  {2\cdot2^{-33}\over1-2^{-1}}=4.65661287308\cdot10^{-10},
\]
so the Bernstein remainder is only \(2.47\cdot10^2\), leaving
\[
  {\|S-S_h\|\over\|S_h\|}
  \le1.18670850497\cdot10^{-6}<5.03\cdot10^{-6}.
\]
Replacing the sampled \(M_2\) by a complex-ball supremum would make this final
ellipse certificate machine-rigorous; the analytic decay mechanism itself is
now explicit.

This last replacement has also been carried out for the finite
endpoint-quadrature Green-source formula.  The Bernstein ellipse \(E_2\) was
covered by \(64\) complex interval arcs, and \(A(.30+.22z)\) was evaluated on
each arc using complex interval arithmetic.  The resulting bound is
\[
  M_2\le 2.92868482364\cdot10^{19},
\]
while the target inequality would still allow
\[
  M_2<3.41434774473\cdot10^{21}.
\]
Thus
\[
  M_2\,{2\cdot2^{-33}\over1-2^{-1}}
  \le1.36377514509\cdot10^{10},
\]
and the total envelope becomes
\[
  5.04566669963\cdot10^{11}.
\]
Consequently
\[
  {\|S-S_h\|\over\|S_h\|}
  \le1.21967465286\cdot10^{-6}<5.03\cdot10^{-6}.
\]
The worst certified arc lies near the left side of the ellipse, with
\[
  \theta\in[2.94524311274,3.04341788317],\quad
  \Re u\in[0.02632420017,0.03028404789],\quad
  \Im u\in[0.01617282815,0.03218990313].
\]
This removes the sampled-ellipse assumption from the source quadrature tail
bound.

The source-side rank theorem can now be closed on the domain-side source Gram.
Let
\[
  S=E^\ast E,\qquad S_N=E_N^\ast E_N,\qquad
  L=(D^7\Lambda,D^8\Lambda)|_{a=.545}.
\]
Let \(P_N\) be the spectral projector of \(S_N\) onto its top two eigenvalues,
and let \(P\) be the corresponding continuum Riesz projector.  If
\[
  \|S-S_N\|\le \varepsilon,\qquad
  \varepsilon< {g_N\over4},\qquad
  g_N=\lambda_2(S_N)-\lambda_3(S_N),
\]
then the conservative Riesz projector estimate gives
\[
  \|P-P_N\|\le \alpha={4\varepsilon\over g_N}.
\]
Writing \(m_N=\sigma_{\min}(LP_N)\) on the finite active two-plane and
\(\ell=\|L\|\), every unit vector \(v\in{\rm Ran}\,P\) satisfies
\[
  \|Lv\|\ge m_N(1-\alpha)-\ell\alpha.
\]
For basis \(22\) and source grid \(129\), the complex-ball quadrature
certificate gives
\[
\begin{array}{c|c}
\varepsilon&2.18612707591\cdot10^5\\
\lambda_1(S_N)&1.79238872854\cdot10^{11}\\
\lambda_2(S_N)&1.34430688269\cdot10^{10}\\
\lambda_3(S_N)&6.50989486394\cdot10^3\\
g_N&1.34430623170\cdot10^{10}\\
\alpha&6.50484844708\cdot10^{-5}\\
m_N&1.20760147549\cdot10^{15}\\
\ell&5.36329870795\cdot10^{17}
\end{array}
\]
and hence
\[
  m_N(1-\alpha)-\ell\alpha
  =1.17263547757\cdot10^{15}>0.
\]
Therefore \(L\) has rank two on the continuum active right eigenspace of
\(S=E^\ast E\).  Since positive source eigenvectors of \(T=EE^\ast\) map by
\(E^\ast\) to the corresponding active eigenvectors of \(S\), this proves
\[
  {\rm rank}(LE^\ast|U_\delta)=2
\]
for the finite-core source-side model.

\subsection{A finite full-theta tail certificate}

The lift from the zero-slope finite core to the full theta kernel must be
measured in the reduced coordinates used in the source theorem.  Raw pointwise
tail bounds do not see the trace null-space, the source-active plane, or the
Volterra normalization.  As a first certificate we therefore compute
\[
  K_{\rm tail}
  =
  K_{\rm red}(\Phi_{\le 8})-K_{\rm red}(\widetilde\Phi_3)
\]
in the same trace-reduced Galerkin model, and then project it to the active
two-plane determined by the Riesz theorem above.

At basis \(22\), Legendre/Laguerre orders \(16\), and source grid \(129\), the
computed constants are
\[
\begin{array}{c|c}
\alpha_3&1.00000000168382189\\
{\rm trace\ rank/nullity}&12/10\\
\|K_{\rm tail}\|/\|K_{\rm core}\|&1.98506712766\cdot10^{-19}\\
\|K_{\rm tail}\|_{\rm active}/\|K_{\rm core}\|_{\rm active}
  &3.74599907107\cdot10^{-16}\\
\|K_{\rm tail}\|_{\rm active}/(1.17263547757\cdot10^{15})
  &6.25741312390\cdot10^{-31}
\end{array}
\]
The active matrices are
\[
K_{\rm core,act}=
\begin{pmatrix}
1.8224151958662063&-0.14132993762900076\\
-0.14132993762900076&1.8123465363577826
\end{pmatrix},
\]
and
\[
K_{\rm tail,act}=
\begin{pmatrix}
6.574084315567217\cdot10^{-16}&-2.284716572476848\cdot10^{-16}\\
-2.284716572476848\cdot10^{-16}&5.015404641457241\cdot10^{-17}
\end{pmatrix}.
\]
This shows that the theta correction through \(n=8\) is perturbative on the
certified active source plane by many orders of magnitude.  This calculation is
not yet the literal continuum \(S_{\rm tail}\) theorem: the remaining analytic
step is to derive the full-theta source-row analogue of \(S=E^\ast E\), prove
the corresponding active/null-space relative bound, and then add an analytic
bound for the \(n\ge9\) theta tail in the same normalized coordinates.

\subsection{The source-row tail object}

We now define the source-tail object directly.  For a theta profile \(\Psi\),
let
\[
  h_u^\Psi(s)=K_{\rm red}^\Psi(s,u).
\]
The source row is the Lagrange residual row for the same endpoint trace
operator \(P=\Lambda_s\):
\[
  E_u^\Psi f
  =
  \left(
    B_P[h_u^\Psi,f](s_0),\,
    (P^\ast h_u^\Psi)(s_0)f(s_0)
  \right).
\]
Thus
\[
  S_\Psi=\int_{u_{\min}}^{u_{\max}}(E_u^\Psi)^\ast E_u^\Psi\,du,
  \qquad
  S_{\rm tail}=S_{\Phi_{\le8}}-S_{\widetilde\Phi_3}.
\]
The derivatives of \(h_u^\Psi\) are obtained by differentiating
\[
K_{\rm red}^\Psi(s,u)
=
\int_0^\infty
\left(r+{s+u\over2}\right)
\cosh\!\left(\omega\left(r+{s+u\over2}\right)\right)
{\Psi(s+r)\over\Psi(s)}
{\Psi(u+r)\over\Psi(u)}\,dr
\]
in the first variable by formal Taylor series at \(s=s_0\).

At basis \(22\), source grid \(65\), and source quadrature order \(16\), the
literal source-tail certificate gives
\[
\begin{array}{c|c}
\|S_{\widetilde\Phi_3}\|&1.40116509611\cdot10^{11}\\
\|S_{\rm tail}\|&4.20138473607\cdot10^{-7}\\
\|S_{\rm tail}\|/\|S_{\widetilde\Phi_3}\|
  &2.99849371621\cdot10^{-18}\\
\|E_{\rm tail}\|/\|E_{\widetilde\Phi_3}\|
  &1.52968731448\cdot10^{-18}\\
\|S_{\rm tail}\|_{\rm active}/\|S_{\widetilde\Phi_3}\|_{\rm active}
  &2.99849358245\cdot10^{-18}\\
4\|S_{\rm tail}\|/g&1.46275282372\cdot10^{-16}
\end{array}
\]
where
\[
  g=1.14889806889\cdot10^{10}
\]
is the source spectral gap between the active two-plane and its complement.
The finite response margin after the tail perturbation remains
\[
  1.20660511212\cdot10^{15}>0.
\]
Hence the rank/noncollapse theorem is stable under the finite full-theta
source-row perturbation through \(n=8\).

For the omitted modes \(n\ge9\), put \(v_{\min}=0.08\).  Since
\(\pi n^2 e^{v_{\min}}\ge 275.7\) for \(n=9\), the elementary bound
\[
\left|D_v^k e^{\beta v-\pi n^2e^v}\right|
\le
e^{\beta v_{\min}-\pi n^2 e^{v_{\min}}}
\left(1+\beta+\pi n^2 e^{v_{\min}}\right)^k
\]
is decreasing on \(v\ge v_{\min}\) for \(0\le k\le8\).  Summing both theta
pieces over \(n\ge9\) gives derivative envelopes
\[
\begin{array}{c|c}
k&\sup_{v\ge v_{\min}} |D^k\Phi_{\ge9}(v)|\\
0&5.95590185276\cdot10^{-115}\\
1&1.66114611409\cdot10^{-112}\\
2&4.63306258485\cdot10^{-110}\\
3&1.29219640663\cdot10^{-107}\\
4&3.60403434542\cdot10^{-105}\\
5&1.00519273893\cdot10^{-102}\\
6&2.80355960708\cdot10^{-100}\\
7&7.81934319692\cdot10^{-98}\\
8&2.18087505677\cdot10^{-95}
\end{array}
\]
The sampled lower envelope of \(\widetilde\Phi_3\) on the source window is
\[
  4.61916023525\cdot10^{-1},
\]
so the zeroth-order relative \(n\ge9\) envelope is
\[
  1.28939061419\cdot10^{-114}.
\]
The last completely formal step is to propagate these derivative envelopes
through the rational normalized map \(\Psi\mapsto E_u^\Psi\) by interval
arithmetic.  Numerically, the available margin is so large that this step has
enormous slack.

This propagation can be carried out directly.  Insert interval coefficients
\([-\tau_k,\tau_k]\) for the omitted tail \(R^{(k)}(v)\), \(0\le k\le8\), into
the Taylor series of every normalized ratio
\[
  {\Psi(x+r)\over\Psi(x)}.
\]
The interval Taylor series are then propagated through
\[
  \Psi\mapsto h_u^\Psi\mapsto E_u^\Psi.
\]
Writing \(\Delta E=E_\Phi-E_{\Phi_{\le8}}\), the source Gram perturbation is
bounded by
\[
  \|\Delta S\|
  \le
  2\|E_8\|\|\Delta E\|+\|\Delta E\|^2.
\]
At basis \(22\), source grid \(65\), interval precision \(35\), source
quadrature order \(10\), and \(r_{\max}=8\), interval propagation gives
\[
\begin{array}{c|c}
\|\Delta E\|&3.47095268287\cdot10^{-28}\\
\|E_8\|&3.74321398817\cdot10^5\\
\|\Delta S\|&2.59850372695\cdot10^{-22}\\
\|\Delta S\|/\|S_8\|&1.85453072886\cdot10^{-33}\\
4\|\Delta S\|/g_8&9.04694262206\cdot10^{-32}
\end{array}
\]
where
\[
  g_8=1.14889806889\cdot10^{10}
\]
is the active/complement gap for \(S_8\).  The lower response bound after this
omitted-tail perturbation remains
\[
  1.20660511212\cdot10^{15}>0.
\]
The smallest interval denominator lower bound in the normalized ratios is
\[
  4.61916023525\cdot10^{-1}.
\]
Thus the omitted \(n\ge9\) theta tail is stable under the normalized source-row
map in the same finite source/Riesz model.  Combining this with the finite
\(\Phi_{\le8}-\widetilde\Phi_3\) source-tail certificate proves the full
source-side noncollapse statement in the certified finite model.

\subsection{Continuum source quadrature for the full theta row}

It remains to replace the finite source grid by the continuum source window.
For \(\Psi_8=\Phi_{\le8}\), write
\[
  S_8=\int_{0.08}^{0.52} (E_u^{\Psi_8})^*E_u^{\Psi_8}\,du.
\]
The script \texttt{full\_theta\_source\_quadrature\_certificate.py} evaluates
the literal full-theta source row and encloses
\[
  \sup_{0.08\le u\le0.52}
  \left\|{d^2\over du^2}\big((E_u^{\Psi_8})^*E_u^{\Psi_8}\big)\right\|
  \le 1.73800843928\cdot10^{15}
\]
by interval Taylor arithmetic on 32 source subintervals.  Therefore the
257-point composite trapezoid rule satisfies
\[
  \|S_8-S_{8,h}\|
  \le {0.44\,h^2\over12}\,1.73800843928\cdot10^{15}
  =1.88255959690\cdot10^8,
  \qquad h={0.44\over256}.
\]
Adding the omitted \(n\ge9\) interval-propagation error
\[
  2.59850372695\cdot10^{-22}
\]
does not change this bound at the displayed precision.

For this same 257-point source Gram, the active eigenvalues are
\[
  1.14890255878\cdot10^{10},\qquad
  1.40116882779\cdot10^{11},
\]
while the top complement eigenvalue is
\[
  1.13804336657\cdot10^4.
\]
Thus
\[
  g=1.14890142073\cdot10^{10},
  \qquad
  \alpha={4\|S_8-S_{8,h}\|\over g}
  =6.55429460851\cdot10^{-2}<1.
\]
Let \(P_h\) denote the finite active projector and \(L=(D^7\Lambda,D^8\Lambda)\)
at \(a=0.545\).  The finite response margin and complement norm are
\[
  \sigma_{\min}(LP_h)=1.20660510522\cdot10^{15},
  \qquad
  \|L(I-P_h)\|=1.81384270956\cdot10^{16}.
\]
The Davis--Kahan angle estimate gives, for a unit vector \(v\) in the continuum
active plane,
\[
  \|Lv\|\ge
  \sigma_{\min}(LP_h)\sqrt{1-\alpha^2}
  -\|L(I-P_h)\|\alpha
  \ge 1.51646525720\cdot10^{13}>0.
\]
Consequently the source-side noncollapse statement survives both continuum
source quadrature and the full theta tail:
\[
  \operatorname{rank}\big(LE_\Phi^*|U_\delta\big)=2
\]
on the certified high block.

\subsection{Return to the global Schur problem}

The preceding result closes the active source-side obstruction for the full
theta kernel.  The remaining global Schur problem is therefore sharply
localized.  The certificate
\texttt{global\_weyl\_volterra\_schur\_bridge.py} records the inputs as follows:
\[
\begin{array}{c|c}
\hbox{full-\(\Phi\) source noncollapse}&\hbox{true}\\
\hbox{projector angle parameter}&6.55429460851\cdot10^{-2}\\
\hbox{rank lower bound}&1.51646525720\cdot10^{13}\\
\hbox{source quadrature error}&1.88255959690\cdot10^8\\
\hbox{omitted theta-tail error}&2.59850372695\cdot10^{-22}.
\end{array}
\]
The sampled active-trace range diagnostics give
\[
  {\|E_{\rm active}-C_{\rm sample}R_{\rm global}\|_F\over
   \|E_{\rm active}\|_F}
  =2.25731056151\cdot10^{-69},
\]
with trace rank equal to active dimension.  Over the active gap scan, the
sampled active trace-kernel carries no source mass, while the largest full
trace-kernel source fraction is
\[
  1.27268094114\cdot10^{-7}.
\]

Thus the active rank obstruction is no longer the bottleneck.  Any remaining
failure of the global Weyl/Volterra Schur form must lie in the source-inactive
high-frequency tail or in the passage from sampled active range inclusion to
the continuum closed trace range.

The next theorem is consequently the source-inactive Schur tail estimate:
\[
  \|(I-P_\delta)E_\Phi f\|^2
  \le \varepsilon_\delta \langle Af,f\rangle,
  \qquad f\in H_M\cap\ker R_{\rm global},
\]
with \(\varepsilon_\delta\) absorbable by the finite low/mid Schur block.
The finite diagnostic already sees small nonzero operator-tail fractions; the
best current nonzero one is
\[
  5.73030971110\cdot10^{-3}
\]
at basis \(16\), tail start \(5\).  This is evidence, not the continuum
theorem.  Once this tail estimate and the continuum active range inclusion are
proved, the closed-trace quotient factorization gives the global
Weyl/Volterra Schur closure.

The same tail statement has now been tested in the full-\(\Phi\) source
coordinates.  On basis \(22\), source grid \(257\), and global trace ratio
\(.75\), the finite \(\Phi_{\le8}\) source Gram has
\[
\begin{array}{c|c}
\lambda_{\max}(S_{8,h})&1.40116882779\cdot10^{11}\\
\hbox{second active eigenvalue}&1.14890255878\cdot10^{10}\\
\hbox{top inactive eigenvalue}&1.13804336657\cdot10^4\\
\hbox{finite inactive/top ratio}&8.12210023518\cdot10^{-8}.
\end{array}
\]
Using the already certified full-theta source quadrature/tail bound
\[
  \|S_\Phi-S_{8,h}\|\le1.88255959690\cdot10^8
\]
gives the continuum-inflated inactive estimate
\[
  {\lambda_3(S_\Phi)\over\lambda_1(S_\Phi)}
  \le
  {1.13804336657\cdot10^4+1.88255959690\cdot10^8
   \over
   1.40116882779\cdot10^{11}-1.88255959690\cdot10^8}
  =1.34545263828\cdot10^{-3}.
\]
On the same block the sampled active trace-kernel source fraction is zero,
the smallest active trace eigenvalue is
\[
  2.03957293979\cdot10^{-7},
\]
and the full trace-kernel source fraction is only
\[
  1.27417836858\cdot10^{-7}.
\]
Thus the full-\(\Phi\) finite certificate has the right scale:
\[
  1.34545263828\cdot10^{-3}
  <
  5.73030971110\cdot10^{-3}.
\]
This still does not replace the continuum Hardy/Schur domination proof, but it
shows that the constants in the normalized full-\(\Phi\) model are compatible
with absorption by the observed finite Schur-tail budget.

The final algebraic absorption step can therefore be stated separately.  If
the continuum high-frequency Hardy/Schur estimate
\[
  \|(I-P_\delta)E_\Phi f\|^2
  \le \varepsilon_\delta \langle Af,f\rangle,
  \qquad f\in H_M\cap\ker R_{\rm global},
\]
is proved with
\[
  \varepsilon_\delta\le 1.34545263828\cdot10^{-3},
\]
then the finite low/mid Schur block absorbs it, since
\[
  1.34545263828\cdot10^{-3}
  <
  5.73030971110\cdot10^{-3}.
\]
The absorption slack is
\[
  4.38485707283\cdot10^{-3},
\]
or equivalently \(\varepsilon_\delta\) uses only about \(23.48\%\) of the
available normalized budget.  Thus the remaining gap is purely analytic:
prove the continuum Hardy/Schur tail estimate itself on
\(H_M\cap\ker R_{\rm global}\).

The source-inactive estimate has an exact min-max form.  Let \(H\) be the
closed-trace high block with \(A\)-inner product, set
\[
  E=E_\Phi:H\to L^2([0.08,0.52];\mathbb C^2),
  \qquad S=E^\ast E,
\]
and let \(P_2\) denote the spectral projection onto the two largest singular
directions of \(E\).  Then the spectral theorem gives
\[
  \|(I-P_2)E f\|^2\le \lambda_3(S)\langle Af,f\rangle .
\]
In the certified full-Phi source model, Weyl perturbation gives
\[
  \lambda_3(S)\le
  \lambda_3(S_{8,h})+\eta
  =1.88267340124\cdot10^8,
\]
while the top source scale satisfies
\[
  \lambda_1(S)\ge
  \lambda_1(S_{8,h})-\eta
  =1.39928626819\cdot10^{11}.
\]
Consequently the small constant is a normalized one.  For
\[
  \widehat E_\Phi
  =
  E_\Phi/(1.39928626819\cdot10^{11})^{1/2},
\]
the min-max theorem gives
\[
  \|(I-P_2)\widehat E_\Phi f\|^2
  \le
  1.34545263828\cdot10^{-3}\,\langle Af,f\rangle .
\]
This normalized estimate is the one absorbed by the finite low/mid Schur
budget.  The remaining analytic passage is the Galerkin-to-continuum high-block
exhaustion estimate which identifies this certified source-model operator with
the operator on \(H_M\cap\ker R_{\rm global}\).

We can now state that passage as a separate theorem.  Let \(H_A\) be the
Hilbert completion for the positive Volterra/Sturm form \(A\), let
\[
  N=\ker R_{\rm global},\qquad
  H_M=N\cap L_M^{\perp_A},
\]
and let
\[
  H_{M,N}=V_N\cap\ker R_N\cap L_{M,N}^{\perp_A}
\]
be the corresponding Galerkin closed-trace high block.  The required
high-block exhaustion theorem is Mosco convergence
\[
  H_{M,N}\longrightarrow H_M
\]
in the \(A\)-graph norm, together with
\[
  \|P_NS P_N-S\|_{A\to A}\to0,\qquad
  S=\widehat E_\Phi^\ast\widehat E_\Phi .
\]
Under these hypotheses, the Riesz projection theorem gives convergence of the
top-two source projections \(P_{2,N}\to P_2\), and the finite min-max estimate
passes to the continuum:
\[
  \|(I-P_2)\widehat E_\Phi f\|^2
  \le
  1.34545263828\cdot10^{-3}\,\langle Af,f\rangle ,
  \qquad f\in H_M\cap\ker R_{\rm global}.
\]
The formal spectral passage is standard.  The remaining analytic input is the
corrected commuted Sturm/source-range trace estimate.  The compact commuted
kernel route alone cannot prove Sobolev coercivity: interior oscillatory
packets in \(\ker R\) make
\[
  \sum_{r=0}^m\langle K D^r f,D^r f\rangle
\]
vanish relative to \(\|f\|_{W^{m,2}}^2\) on any infinite-dimensional
finite-codimension subspace.  The valid target is instead the source-range
Hardy/Green estimate
\[
  E_u^\ast E_u\le \eta_u A,\qquad
  f\in H_M\cap\ker R_{\rm global},
\]
uniformly over the source window.  The finite commuted model supports this
target: the best subunit domination has \(M=5,m=10\) and product
\(2.89082003770\cdot10^{-5}\), while the auxiliary source-range certificate has
high/full fraction \(1.19373561052\cdot10^{-6}\).  Once the continuum
source-range Hardy/Green estimate is proved, the direct Lagrange/adjoint source
row
\[
  E_u f=
  \left(B_P[h_u,f](s_0),\ (P^\ast h_u)(s_0)f(s_0)\right)
\]
has uniform \(A\)-Riesz representers.  This makes
\(\widehat E_\Phi^\ast\widehat E_\Phi\) compact in the \(A\)-Hilbert source
model, yielding the required operator norm convergence after Galerkin
exhaustion.

The source-range estimate is equivalently a Green-representer theorem.  On
\[
  H_{\rm hi}=H_M\cap\ker R_{\rm global},\qquad
  \langle f,g\rangle_A=\langle Af,g\rangle,
\]
write \(E_u f=(\ell_1(u;f),\ell_2(u;f))\).  Then
\[
  E_u^\ast E_u\le\eta_u A
\]
holds if and only if the scalar rows have \(A\)-Riesz representers
\[
  \ell_k(u;f)=\langle g_{u,k},f\rangle_A,\qquad k=1,2.
\]
The sharp \(\eta_u\) is the largest eigenvalue of
\[
  G_{ij}(u)=\langle g_{u,i},g_{u,j}\rangle_A .
\]
The current finite Green certificate has maximum source constant
\(9.34694704598\cdot10^5\), maximum differentiated-source constant
\(1.31172699716\cdot10^5\), and range residual
\(1.68\cdot10^{-80}\).  The continuum task is therefore exactly to construct
the closed-trace Green representers \(g_{u,k}\) and bound \(G(u)\) uniformly on
the compact source window.

The direct construction is by the Riesz map on the completed high closed-trace
\(A\)-space:
\[
  g_{u,k}=J_A^{-1}\ell_k(u;\cdot),
  \qquad
  \ell_k(u;f)=\langle g_{u,k},f\rangle_A .
\]
Thus the remaining theorem is the uniform \(A\)-boundedness of the actual
source rows, not of arbitrary endpoint jets.  The current direct-row
certificate gives
\[
  \max_u\|E_u\|_{A\to{\bf C}^2}^2=9.34694704598\cdot10^5,
  \qquad
  \max_u\|\partial_uE_u\|_{A\to{\bf C}^2}^2=1.31172699716\cdot10^5,
\]
with range residual \(1.68\cdot10^{-80}\).  The global interval trace scan
gives maximum high/full source fraction \(1.19437533782\cdot10^{-6}\), while
the active trace-to-source finite residual is \(2.41\cdot10^{-69}\).

The fixed-jet decomposition remains a useful sufficient route.  If
\[
  k_j^{\rm hi}\in H_{\rm hi},\qquad
  f^{(j)}(s_0)=\langle k_j^{\rm hi},f\rangle_A,\quad 0\le j\le7,
\]
then
\[
  g_{u,1}=\sum_{j=0}^7 b_j(u)k_j^{\rm hi},\qquad
  g_{u,2}=p(u)k_0^{\rm hi}.
\]
The coefficient functions \(p,b\) and their derivatives are already controlled
by the Volterra source formula on the compact interval \(u\in[.08,.52]\).
The finite fixed-representer scan verifies the factorization to roundoff; for
example at basis \(20\),
\[
  \|k_0^{\rm hi}\|_A^2=8.01690618322\cdot10^1,
\]
while the largest value in the basis \(18,20,22\) scan is
\[
  2.02962047830\cdot10^3.
\]
But earlier local-tower scans show that raw high-order jet representers grow
rapidly, so this is likely too strong as the standalone proof.  The corrected
continuum Green-representer theorem is:
\[
  P_\delta E_u=\int_I K_u(a)\Lambda_a\,da
  \quad\hbox{on the source-active block}
\]
together with the source-inactive Hardy/Schur bound.  Restricted to
\(\ker R_{\rm global}\), the active trace part vanishes and the inactive part
is \(A\)-bounded; Riesz then constructs \(g_{u,1},g_{u,2}\), and the
\(2\times2\) Gram \(G(u)\) is uniformly bounded on \(u\in[.08,.52]\).

The finite Green-kernel certificate constructs \(K_N\) from
\[
  E_{\rm active}=C_NR_N,\qquad C_{N,i}=w_iK_N(a_i).
\]
It gives finite active range residual \(2.41\cdot10^{-69}\), maximum sampled
density \(L^2\)-norm \(3.49514838830\cdot10^2\), maximum refined weighted
density \(L^2\)-norm \(3.39266687955\cdot10^2\), and relative density spread
\(2.49413419264\cdot10^{-2}\).  A direct finite-difference differentiation of
the coarse pseudoinverse kernel gives a large adjoint integration-by-parts
defect, about \(2.48\cdot10^6\).  Thus \(K_N\) is good evidence for a
trace-dual kernel, but not itself a smooth Green solution.  The remaining proof
is the continuum adjoint Green boundary problem \(P^\ast K_u=\eta_u\) with the
endpoint concomitant matching the active source row.

A first regularized replacement chooses a polynomial density \(K_N^{\rm reg}\)
which satisfies the two active moment equations exactly and minimizes a fixed
Sobolev-type polynomial norm.  In the smoke test with 11 trace samples and
degree 6 this keeps the active range residual at \(9.80\cdot10^{-70}\) and
reduces the adjoint integration-by-parts defect to about \(1.15\cdot10^5\).
This is evidence that smoothing moves in the right direction, but it still does
not solve the adjoint boundary-value problem.  The proof must construct
\(K_u\) from the BVP with endpoint concomitant constraints built in.

The interior singular part of that BVP can be written explicitly.  For a local
source row
\[
  \ell(f)=\sum_{q=0}^7 d_q f^{(q)}(s_0),
\]
the representing distribution is
\[
  \sum_{q=0}^7(-1)^q d_q\delta_{s_0}^{(q)}.
\]
Thus the adjoint Green coefficient is piecewise homogeneous for \(P^\ast\) on
\((a_-,s_0)\cup(s_0,a_+)\).  If
\[
  \Delta_r=K^{(r)}(s_0+)-K^{(r)}(s_0-),\qquad 0\le r\le7,
\]
then the jump equations are
\[
  \sum_{k=q+1}^8(-1)^k
  \left[D^{k-1-q}(a_kK)\right]_{s_0}
  =(-1)^q d_q,\qquad 0\le q\le7.
\]
Equivalently \(J\Delta={\rm rhs}(d)\), with
\[
  J_{q,r}=\sum_{k=q+1}^8(-1)^k
  {k-1-q\choose r}\,D^{k-1-q-r}a_k(s_0),
\]
omitting terms with \(r>k-1-q\).  The computed jump matrix has
\[
  \det J=1.54269213181\cdot10^{-43},\qquad
  \|J^{-1}\|_F=6.55172220974\cdot10^{16},
\]
and solves the jump equations for both source components to zero residual.
The source singularity is therefore no longer the open part.  What remains is
the endpoint concomitant: choose the homogeneous constants on the two sides so
that the endpoint term is killed, or is contained in the source-inactive tail,
with uniform trace-dual norm control.

Equivalently, expand
\[
  P^\ast K=\sum_{m=0}^8 c_m(a)K^{(m)},\qquad
  c_m(a)=\sum_{k=m}^8(-1)^k{k\choose m}D^{k-m}a_k(a),
\]
so \(c_8=a_8\).  With \(Y=(K,K',\ldots,K^{(7)})^T\), the homogeneous equation
on either side of \(s_0\) is \(Y'=A(a)Y\), where the last row of \(A\) is
\[
  (-c_0/c_8,\,-c_1/c_8,\ldots,-c_7/c_8).
\]
If \(\Phi(a,t)\) is the fundamental matrix and
\(\Delta(d)=J^{-1}{\rm rhs}(d)\), then every distributional Green coefficient
for the source row \(d\) is of the form
\[
  Y(a)=\Phi(a,s_0)z,\quad a<s_0,\qquad
  Y(a)=\Phi(a,s_0)(z+\Delta(d)),\quad a>s_0.
\]
The remaining endpoint concomitant is therefore an affine function of
\(z\in\mathbb R^8\).  The next theorem is finite-dimensional at the ODE
boundary level: choose \(z\) so that the active endpoint concomitant vanishes,
or is absorbed by the source-inactive tail, uniformly in \(u\).

The finite endpoint diagnostic constructs this map as
\[
  P_\delta\Beta_z=Mz+b(d)
\]
on the two active source modes.  In the sampled homogeneous-flow model the
endpoint map has rank \(1\), not full row rank, but the actual source endpoint
vectors \(b(d_u)\) lie in its range.  The minimum-norm choices of \(z\) give
maximum active endpoint residual
\[
  4.90427867176\cdot10^{-27}
\]
and maximum selected norm
\[
  \|z\|\le 2.79142983795.
\]
Thus the correct endpoint theorem is not full rank of \(M\), but the sharper
range condition \(b(d_u)\in{\rm Range}(M)\), with a uniform bound on the
corresponding right inverse for the actual source family.

This statement can be made without the sampled homogeneous-flow approximation.
Let \(\Phi(a,t)\) be the exact fundamental matrix for the adjoint first-order
system and let \(E_\pm\) be the endpoint concomitant matrices after projection
onto the active source space.  Then
\[
  M=E_+\Phi(a_+,s_0)-E_-\Phi(a_-,s_0),
  \qquad
  b(d)=E_+\Phi(a_+,s_0)\Delta(d),
\]
and the endpoint equation is exactly \(Mz=-b(d)\).  Hence
\[
  b(d)\in {\rm Range}(M)
  \quad\Longleftrightarrow\quad
  w^\ast b(d)=0\quad\text{for all }w\in\ker M^\ast .
\]
When this compatibility holds the canonical choice is
\[
  z(d)=-M^\dagger b(d),
  \qquad
  \|z(d_u)\|\le \|M^\dagger\|\,\|b(d_u)\|.
\]
If the rank of \(M\) is constant and the active compatibility holds for the
continuous source family \(d=d_u\), compactness of the source window gives the
required uniform trace-dual bound for the piecewise homogeneous Green
coefficient.  Thus the endpoint step is now an exact fundamental-matrix
Fredholm reduction; the remaining input is the global active trace range
compatibility.

The latter has a precise Hilbert-space form.  Let \(H_A\) be the completed
high block, let \(R f=(a\mapsto \Lambda_a(f))\) on the interval \(I\), and let
\(E_{\rm act}\) denote the active source component.  Then
\[
  E_{\rm act}\in\overline{{\rm Ran}\,R^\ast}
  \quad\Longleftrightarrow\quad
  \ker R\subset \ker E_{\rm act}.
\]
Thus failure of active range compatibility is equivalent to an
\(A\)-normalized sequence satisfying
\[
  \|f_n\|_A=1,\qquad Rf_n\to0,\qquad
  \|E_{\rm act}f_n\|\ge \delta .
\]
The commuted Sturm/Mosco compactness theorem should pass such a sequence to a
closed-trace high-block limit with \(\Lambda_a(f)=0\) on \(I\).  The remaining
analytic lemma is the corresponding active unique-continuation statement:
\[
  \Lambda_a(f)=0\ \text{on }I
  \quad\Longrightarrow\quad
  E_{\rm act}f=0 .
\]
Once this is proved, active range compatibility follows from the annihilator
criterion and the endpoint Fredholm compatibility follows automatically.

The direct Volterra/Sturm proof of this implication is a Green identity.  Let
\[
  P f(a)=\Lambda_a(f),\qquad
  P^\ast K=\sum_{k=0}^8(-1)^kD_a^k(a_kK).
\]
For a scalar active source row \(\ell\), suppose \(K_\ell\) is piecewise
smooth, satisfies \(P^\ast K_\ell=\ell\) as a distribution, and has endpoint
concomitant killed in the active projection.  Then Green's formula gives
\[
  {d\over da}B_P[K_\ell,f]
  =K_\ell(a)\Lambda_a(f)-f(a)P^\ast K_\ell(a).
\]
After integration across \(I\), including the jump at \(s_0\), the endpoint
terms vanish and
\[
  \ell(f)=\int_I K_\ell(a)\Lambda_a(f)\,da .
\]
Thus \(\Lambda_a(f)=0\) on \(I\) implies \(\ell(f)=0\).  Applying this to the
active source rows gives \(E_{\rm act}f=0\).

The Lagrange identity, the distributional jump law, and the exact endpoint
fundamental-matrix reduction are closed.  The remaining non-formal point is to
prove endpoint Fredholm compatibility for the exact continuum Green
boundary-value problem, i.e. that every homogeneous endpoint obstruction
annihilates the actual active source family, without importing active range
inclusion itself.

This compatibility can be stated without reference to the sampled model.  For
the exact endpoint equation
\[
  Mz=-b(d),
\]
Fredholm compatibility is
\[
  w^\ast b(d)=0,\qquad w\in\ker M^\ast .
\]
The Green pairing identifies each left endpoint obstruction with a primal
homogeneous endpoint test \(f_w\).  For the jump-source coefficient \(K_d\),
\[
  w^\ast b(d)
  =\int_I K_d(a)\,P f_w(a)\,da-d(f_w).
\]
Hence, on a genuine primal homogeneous endpoint obstruction, \(P f_w=0\) and
the obstruction is simply \(-d(f_w)\).  The remaining continuum theorem is
therefore the non-circular annihilation statement
\[
  d_u(f_w)=0
\]
for every actual active source row \(d_u\) and every primal homogeneous
endpoint obstruction \(f_w\).  The finite endpoint certificate gives residual
\(4.91\cdot10^{-27}\) for the full sampled endpoint solve, but this is not an
annihilation theorem for a Fredholm obstruction.  The resolved
relative-threshold left endpoint direction has relative pairing about
\(2.70\cdot10^{-1}\) with the actual source endpoint vectors.  Thus the finite
sampled rank threshold cannot prove compatibility; the continuum proof must
come from the Volterra/Sturm structure.

There is, however, a cleaner full-rank alternative.  The endpoint rank-ball
certificate gives
\[
  \sigma_{\min}(M_0)=2.77565955249\cdot10^6,\qquad
  \sigma_{\max}(M_0)=7.35305512920\cdot10^{25},
\]
so the relative singular margin is \(3.77483849056\cdot10^{-20}\).  With the
diagnostic \(10^{-22}\) relative endpoint-entry ball one obtains
\[
  \sigma_{\min}(M)\ge 2.76830649736\cdot10^6,\qquad
  \|M^+\|\le 3.61231677545\cdot10^{-7}.
\]
This is currently conditional, not a continuum proof: the endpoint-entry ball
must still be derived from a rigorous interval enclosure of the exact
fundamental matrix.  The exact certification target is
\[
  \|M-M_0\|_F < 2.77565955249\cdot10^6,
\]
equivalently a relative Frobenius enclosure better than
\(3.77483849056\cdot10^{-20}\).  Once this is proved, \(\ker M^\ast=0\) and
endpoint compatibility is vacuous.

The present endpoint-flow center does not yet support this conclusion.  Mesh
refinement of the finite-difference construction gives endpoint-map Frobenius
norms
\[
  7.35\cdot10^{25},\qquad 5.56\cdot10^{13},\qquad
  6.75\cdot10^{38}
\]
for \(11,13,15\) sampled ODE nodes, respectively.  Replacing finite
differences of the moving eigenrow by exact confluent eigen-derivatives gives
rank-two centers, but the current piecewise-constant propagation still changes
from \(1.81\cdot10^8\) to \(3.68\cdot10^{17}\) between \(5\) and \(7\)
segments.  Thus the displayed inequality has not been proved for the old
\(M_0\).  The next required artifact is a validated endpoint-flow solver:
analytic/confluent trace derivatives for the coefficient field, followed by a
Taylor or Chebyshev integration of the first-order system with a rigorous
remainder.

A first Chebyshev-collocation center has been built using exact confluent
trace derivatives at Chebyshev-Lobatto nodes.  This removes the catastrophic
finite-difference instability and gives rank-two endpoint maps, but the
sequence is not yet stable at the rank margin:
\[
\begin{array}{c|c|c}
N & \|M_N\|_F & \sigma_{\min}(M_N)\\
\hline
5 & 6.04\cdot10^6 & 2.79\cdot10^4\\
7 & 1.84\cdot10^7 & 2.56\cdot10^4\\
9 & 3.31\cdot10^6 & 2.69\cdot10^4 .
\end{array}
\]
Moreover
\[
  \|M_7-M_5\|_F=1.24\cdot10^7,\qquad
  \|M_9-M_7\|_F=2.14\cdot10^7,
\]
which is still larger than the \(2.775\cdot10^6\) rank margin.  Thus the
Chebyshev center is a much better numerical object, but it does not yet close
the full-rank route.

Simple stabilization does not yet fix this.  Taylor-scaled state coordinates
give the same Chebyshev refinement differences, as expected from similarity
invariance of the collocation solve.  A fourth-order Gauss--Magnus stepper
with local Taylor scaling is worse: the endpoint-map Frobenius norm grows from
\(3.86\cdot10^{15}\) at two steps to \(1.76\cdot10^{138}\) at four steps.
Thus the full-rank route now requires a more structural endpoint-flow
representation, or else one must return to the direct
obstruction-annihilation theorem.

The first structural replacement is the adjoint row-flow representation.  If
\(Y'=A(s)Y\) is the adjoint first-order system, then endpoint covectors
propagate by
\[
  r'(s)=-r(s)A(s),
\]
and the endpoint map can be formed as the difference of the right- and
left-propagated endpoint row spaces at \(s_0\), without constructing the full
ill-conditioned fundamental matrix.  The diagnostic
\texttt{endpoint\_adjoint\_row\_flow\_center.py} implements this row flow, computes the
two active endpoint rows, and scans all \(2\times2\) minors.  For Chebyshev
orders \(5,7,9\) it gives rank-two endpoint maps,
\[
\begin{array}{c|c|c|c}
N & \|M_N\|_F & \sigma_{\min}(M_N) & \max |\det M_N[:,J]|\\
\hline
5 & 2.93\cdot10^7 & 1.14\cdot10^4 & 2.91\cdot10^{11}\\
7 & 1.34\cdot10^7 & 3.30\cdot10^6 & 4.14\cdot10^{13}\\
9 & 2.09\cdot10^7 & 1.00\cdot10^5 & 2.01\cdot10^{12}.
\end{array}
\]
The same best minor persists under these refinements, but its relative change
is about \(1.01\) from \(5\) to \(7\) and \(21.6\) from \(7\) to \(9\).
Thus the row-flow object is the right structural center, but it is not yet a
minor certificate.  The remaining full-rank shortcut would need an
exterior-product, Riccati, or Grassmannian normalized two-plane flow with a
rigorous enclosure of the active minor.

The normalized exterior diagnostic makes this precise.  Let
\[
  p_{ij}=\det M[:,(i,j)],\qquad \widehat p={p\over \|p\|_2}.
\]
The script \texttt{endpoint\_grassmann\_flow\_center.py} computes these projective
Grassmannian coordinates from the row-flow centers and chooses the affine
chart with largest persistent normalized minor.  On the refined sequence
\(N=7,9,11\), the same chart \(J=(0,1)\) dominates:
\[
\begin{array}{c|c|c}
N & |\widehat p_{01}| & \|p\|_2\\
\hline
7 & 0.9632738211 & 4.29\cdot10^{13}\\
9 & 0.9564654848 & 2.10\cdot10^{12}\\
11& 0.9643323749 & 3.82\cdot10^{12}.
\end{array}
\]
The projective two-plane distances are
\[
  d_{\rm Gr}(7,9)=2.53\cdot10^{-2},\qquad
  d_{\rm Gr}(9,11)=2.92\cdot10^{-2}.
\]
Thus the raw determinant scale is not the right quantity to certify.  The
proper full-rank shortcut is now an interval enclosure in the chart
\(p_{01}\ne0\), or equivalently a ball/Riccati proof that the normalized
Pluecker coordinate \(|\widehat p_{01}|\) remains bounded away from zero.

The chart-ball step is elementary but important.  Let \(v\) be the order
\(11\) normalized Pluecker center, oriented so that \(v_{01}>0\).  If the
exact normalized Pluecker vector \(x\) lies in the projective ball
\[
  \min(\|x-v\|_2,\|x+v\|_2)\le r,
\]
then
\[
  |\widehat p_{01}(x)|\ge |v_{01}|-r.
\]
Here \(|v_{01}|=0.9643323749\).  Therefore any rigorous exterior/Riccati
enclosure with \(r\le0.125\) gives
\[
  |\widehat p_{01}(x)|\ge0.8393323749,
\]
and the weaker threshold \(|\widehat p_{01}|\ge1/2\) allows radius
\[
  r\le0.4643323749.
\]
The latest computed projective movement is \(2.92\cdot10^{-2}\).  In contrast,
certifying the same threshold by an entrywise ball around the raw endpoint map
would require uniform absolute entry radius about \(3.64\cdot10^4\), which is
far too sensitive to the raw flow normalization.  Thus the rigorous endpoint
rank shortcut should be made in the Riccati/Pluecker chart, not in the raw
matrix entries.

The exact Riccati equation is now explicit.  If \(M'(s)=-M(s)A(s)\), then the
Pluecker coordinates satisfy
\[
  p_{ij}'=-\sum_k A_{ki}p_{kj}-\sum_k A_{kj}p_{ik}.
\]
Writing \(z_J=p_J/p_{01}\) and \(z_{01}=1\), this becomes
\[
  z_J'=(Bz)_J-z_J(Bz)_{01}.
\]
The generator \(B(s)\) has been constructed from the exact confluent trace
derivatives.  On the current chart tube the raw-coordinate bounds are
\[
  \|B\|_\infty\le1.41\cdot10^{12},\qquad
  \operatorname{Lip}(F)\le3.03\cdot10^{12}.
\]
Thus the ODE structure is fixed, but a direct Gronwall validation is useless:
the residual budget is about \(9.85\cdot10^{-688774248833}\).  The endpoint
rank proof must therefore use a Krawczyk/interval-collocation enclosure in
the Riccati chart, or a balanced projective coordinate system with a much
smaller logarithmic norm.

The corresponding finite Krawczyk reduction has now been computed.  The two
row-flow collocation systems feeding the endpoint map have condition estimates
\[
  \kappa_\infty(L)\le5.71\cdot10^6,\qquad
  \kappa_\infty(R)\le1.49\cdot10^6.
\]
At order \(11\), the recomputed endpoint map agrees with the row-flow center to
about \(5.09\cdot10^{-10}\).  The Krawczyk bounds show that the chart threshold
\(|\widehat p_{01}|\ge1/2\) follows if the scaled companion matrices are
enclosed entrywise within
\[
  1.27\cdot10^{-9}
\]
and the active endpoint boundary rows within
\[
  5.80\cdot10^{-10},
\]
with a simultaneous safety fraction \(0.499995\) when both uncertainties are
used at once.  Thus the remaining exact numerical proof has been reduced to
interval enclosures for the confluent trace-derivative coefficients and the
active endpoint boundary rows below these radii.

The finite coefficient-input enclosure has also been checked.  With explicit
sample balls of relative radius \(10^{-40}\), absolute radius \(10^{-45}\), and
safety factor \(16\), the scaled companion radius is
\[
  1.09\cdot10^{-36},
\]
which is \(8.54\cdot10^{-28}\) of the Krawczyk capacity.  The active boundary
row radius is
\[
  2.05\cdot10^{-32},
\]
which is \(3.53\cdot10^{-23}\) of its capacity.  Hence the remaining formal
numerical step is no longer the Krawczyk algebra; it is to replace the
sample-ball model by interval quadrature and tail bounds for the confluent
integrals defining the trace derivatives.

This replacement has now been carried out at the finite Krawczyk-input level.
The script \texttt{endpoint\_coefficient\_interval\_enclosure.py} fixes the
active source coordinates from the base Krawczyk model, then compares the
strict boundary rows in that same active frame.  This avoids the spurious large
differences caused by recomputing an arbitrary active basis at each precision.
Using the refinement ladder \(70{:}70\to80{:}80\to90{:}90\) and a geometric
refinement-tail allowance gives
\[
  r_{\rm comp}=7.75070521382\cdot10^{-51},\qquad
  r_{\rm bdry}=1.54801512661\cdot10^{-46}.
\]
Relative to the Krawczyk capacities these are
\[
  r_{\rm comp}/c_{\rm comp}=6.07978004306\cdot10^{-42},\qquad
  r_{\rm bdry}/c_{\rm bdry}=2.66687603810\cdot10^{-37}.
\]
Both are also far inside the simultaneous Krawczyk budget.  Thus the finite
endpoint-flow coefficient input is closed in the persistent Pluecker chart;
the remaining formal numerical item is to justify the same refinement tail by
analytic quadrature, \(r\)-tail, and eigenrow perturbation estimates for the
confluent trace derivative construction.

The lower trace layer was then verified by
\texttt{endpoint\_confluent\_trace\_tail\_certificate.py}.  The node keys must
be generated at high precision; otherwise the cache lookup silently asks for a
different set of collocation nodes.  With this correction, the needed trace
derivatives are stable.  Only derivatives through \(q=8\) enter the companion
and endpoint boundary-row maps.  On these rows the refinement ladder gives
\[
  r_{\rm trace}=5.72755256199\cdot10^{-42},
  \qquad
  r_{\rm row}=6.12287316544\cdot10^{-42},
\]
well below the working target \(10^{-25}\).  The absolute \(r\ge12\) panel-tail
proxy is
\[
  2.62222776038\cdot10^{-445078},
\]
while the comparison scale from the final coefficient radii is
\(7.75070521382\cdot10^{-51}\).  The minimum negative-eigenrow gap over the
collocation nodes and refinement levels is about
\[
  1.16355182786\cdot10^{-5},
\]
with minimum consecutive eigenrow cosine about \(0.988171806397\).  Therefore
the finite trace-tail certificate feeding the Krawczyk shortcut is closed.

The geometric quadrature-refinement model has now been replaced by an explicit
Bernstein segment estimate.  The script
\texttt{endpoint\_confluent\_segment\_bernstein\_certificate.py} subdivides
\([0,12]\) into \(48\) panels of length \(0.25\).  Each panel admits the
Bernstein ellipse \(E_2\).  For a panel \(r=m+hz\), the proof uses
\[
  \left|\int g-Q_Ng\right|
  \le
  h\,{8M_\rho\rho^{-2N}\over 1-\rho^{-1}},
\]
where \(M_\rho\) is bounded by propagating absolute majorants through the exact
confluent Taylor recurrence.  The majorant keeps the decay term
\[
  \Re(e^r-1)\ge e^{\Re r_{\min}}\cos |\Im r|_{\max}-1
\]
and hence avoids the dependency blow-up seen in direct interval evaluation.
With \(N=200\), the deterministic segment bound gives
\[
  \epsilon_{\rm entry}=2.00296575232\cdot10^{-61},\qquad
  \epsilon_{\rm Taylor,spec}=1.94972191521\cdot10^{-54}.
\]
This is below both the working trace target \(10^{-25}\) and the downstream
coefficient-radius scale \(7.75070521382\cdot10^{-51}\).  Thus the segment
quadrature theorem itself no longer depends on refinement extrapolation.

The eigenrow propagation has also been separated from any heuristic
amplification factor.  The script
\texttt{endpoint\_eigenrow\_interval\_propagation.py} treats the Taylor
coefficients \(A_p\) as interval data and applies a Krawczyk/Newton test to
the finite system
\[
  \sum_{p=0}^n A_pv_{n-p}
  =
  \sum_{p=0}^n\lambda_pv_{n-p},\qquad
  \sum_{p=0}^n\langle v_p,v_{n-p}\rangle=\delta_{n0},
  \quad 0\le n\le 8.
\]
The system has \(90\) scalar equations.  Using the exact Jacobian at the
strict endpoint center, the all-node certificate gives
\[
  \min\operatorname{gap}=1.16355182786\cdot10^{-5},\qquad
  \max\kappa_{\rm Kraw}=7.08972910479\cdot10^{-18},
\]
and
\[
  \max_{q\le8,k}|e_k^{(q)}-\widehat e_k^{(q)}|
  \le
  4.44737689405\cdot10^{-29}.
\]
The corresponding row-\(\ell^2\) radius is \(1.33421306822\cdot10^{-28}\),
still below the trace target \(10^{-25}\).  The synchronized \(200\)-point
rerun gives the same worst radius, with condition matrix order equal to the
controlled Bernstein quadrature order.  Hence the eigenrow Taylor recurrence
itself is certified; the artificial \(10^{60}\) amplification placeholder is
removed.

Finally, the endpoint coefficient input was regenerated at the synchronized
center by
\texttt{endpoint\_coefficient\_synchronized\_200\_certificate.py}.  The active
source frame is fixed to the base Krawczyk frame, and the \(200\)-point
companion and endpoint rows are projected into this same frame.  The analytic
derivative radius propagates to
\[
  r_{\rm comp}=2.03898743318\cdot10^{-24},\qquad
  r_{\rm bdry}=9.06230013777\cdot10^{-25}.
\]
Relative to the simultaneous Krawczyk budgets, these use only
\[
  3.19886168029\cdot10^{-15},\qquad
  3.12248481588\cdot10^{-15},
\]
respectively.  The induced endpoint-entry radius is
\[
  1.15021046920\cdot10^{-10},
\]
whereas the chart capacity is \(3.63920183454\cdot10^{4}\); the ratio is
\(3.16061191847\cdot10^{-15}\), and the Krawczyk contraction parameter is
\(1.09385546388\cdot10^{-19}\).  Thus the synchronized finite endpoint
Krawczyk input is closed.

The synchronized endpoint certificate can now be imported into the active range
argument.  Since the active endpoint map has full row rank, the adjoint endpoint
Green boundary-value problem is solvable for every active source row.  Hence
the Green representation has only the interval trace term, and
\[
  R_{\rm global}f=0 \qquad\Longrightarrow\qquad E_{\rm active}f=0.
\]
By the Hilbert annihilator criterion this is equivalent to
\[
  E_{\rm active}\in \overline{\operatorname{Ran}R_{\rm global}^{\ast}}.
\]
Combining this with the source-inactive min--max theorem gives, in the
certified normalized source model,
\[
  P_{\rm active}\widehat E f=0,\qquad
  \|(I-P_{\rm active})\widehat E f\|^2
  \le 1.345452638275\cdot10^{-3}\,\langle Af,f\rangle,
  \qquad f\in\ker R_{\rm global}.
\]
The available finite low/mid Schur budget is
\[
  5.730309711104\cdot10^{-3},
\]
leaving slack \(4.384857072828\cdot10^{-3}\).  Thus the remaining formal gap is
not active range inclusion, but the Galerkin-to-continuum high-block exhaustion
and commuted elliptic estimate needed to lift the source-inactive min--max bound
from the certified finite source model to the full closed trace space.

This passage can be separated into a closed abstract theorem and one remaining
analytic lower-bound input.  Let \(\Pi_N\) denote the \(A\)-orthogonal projection
onto the sampled closed-trace high block \(H_{M,N}\), and let
\[
  S=\widehat E_\Phi^\ast \widehat E_\Phi .
\]
If \(H_{M,N}\to H_M\) in the Mosco sense and \(S\) is compact on the
\(A\)-Hilbert space, then \(\Pi_N\to \Pi\) strongly and compactness gives
\[
  \|\Pi_N S\Pi_N-S\|_{A\to A}\to0.
\]
Consequently the Courant--Fischer min--max values converge, in particular
\(\lambda_3(S_N)\to\lambda_3(S)\), and the finite inactive-tail estimate passes
to the full closed space.  The Mosco recovery sequence is explicit: given
\(v_N\to f\) in the \(A\)-graph norm with \(R_{\rm global}f=0\), trace
consistency gives \(R_Nv_N\to0\).  A uniform trace-frame lower bound
\(\gamma_0>0\) supplies a correction \(w_N\) with
\[
  R_Nw_N=R_Nv_N,\qquad
  \|w_N\|_A\le \gamma_0^{-1/2}\|R_Nv_N\|,
\]
so \(f_N=v_N-w_N\in H_{M,N}\) and \(f_N\to f\).  The liminf direction is weak
compactness in the \(A\)-Hilbert space plus trace-quadrature consistency, which
passes \(R_Nf_N=0\) to \(R_{\rm global}f=0\).

Thus the invalid target was raw compact-kernel Sobolev coercivity.  The correct
replacement is the continuum trace-frame lower bound, and this now follows from
the source/range theorems.  Let \(H_\delta\) be the active source spectral
space.  The full-\(\Phi\) source-side noncollapse theorem makes
\(E_{\rm active}\) injective on \(H_\delta\).  The synchronized endpoint Green
theorem gives
\[
  R_{\rm global}f=0\quad\Longrightarrow\quad E_{\rm active}f=0.
\]
Hence \(R_{\rm global}\) is injective on \(H_\delta\).  Since \(H_\delta\) is
finite dimensional and \(a\mapsto \Lambda_a(f)\) is continuous,
\[
  q(f)=\int_I |\Lambda_a(f)|^2\,da
\]
is continuous and strictly positive on the \(A\)-unit sphere of \(H_\delta\).
Therefore
\[
  \int_I |\Lambda_a(f)|^2\,da
  \ge \gamma_\delta \|f\|_A^2,\qquad \gamma_\delta>0.
\]
On this finite-dimensional continuous family, positive trace quadrature
converges uniformly.  Thus the sampled weighted frame matrices converge to the
continuum frame, and for sufficiently fine trace meshes their lower eigenvalue
is at least \(\gamma_\delta/2\).  The sampled trace correction right inverse is
therefore bounded by \((2/\gamma_\delta)^{1/2}\).

Consequently the normalized source-inactive estimate passes to the full closed
high block:
\[
  \|(I-P_{\rm active})\widehat E_\Phi f\|^2
  \le 1.345452638275\cdot10^{-3}\,\langle Af,f\rangle,
  \qquad f\in H_M\cap\ker R_{\rm global}.
\]
This is absorbed by the finite low/mid Schur budget
\(5.730309711104\cdot10^{-3}\), with slack
\(4.384857072828\cdot10^{-3}\).  The finite frame data still record an observed
lower floor \(2.931151510091\cdot10^2\), but the present argument only needs and
proves qualitative positivity of \(\gamma_\delta\); certifying that numerical
floor would require a separate interval quadrature-error bound.

\subsection{Quotient-Schur Assembly for the Full Kernel}

We can now apply the abstract closed-trace quotient theorem above to the
normalized full-\(\Phi\) source/Volterra model.  Let \(N=\ker R_{\rm global}\),
write \(V=N\oplus U\), and block the
Hermitian form as
\[
  q(n+u,m+v)
  = a(n,m)+b(n,v)+\overline{b(m,u)}+c(u,v).
\]
The two hypotheses needed for the closed-trace quotient factorization are
\[
  a\ge0\quad\text{on }N,
  \qquad
  b(n,u)=\langle A^{1/2}n,\Gamma u\rangle_{\mathcal H_A}
\]
for a bounded \(\Gamma:U\to\mathcal H_A\).  The preceding two theorems verify
these hypotheses directly.

First, the synchronized endpoint Green theorem gives
\[
  E_{\rm active}\in\overline{\operatorname{Ran}R_{\rm global}^{\ast}},
  \qquad
  R_{\rm global}f=0\Rightarrow E_{\rm active}f=0.
\]
Second, the full-continuum source-inactive theorem gives
\[
  \|(I-P_{\rm active})\widehat E_\Phi f\|^2
  \le 1.345452638275\cdot10^{-3}\,\langle Af,f\rangle,
  \qquad f\in H_M\cap\ker R_{\rm global}.
\]
This is absorbed by the finite low/mid Schur budget
\[
  5.730309711104\cdot10^{-3},
\]
with slack \(4.384857072828\cdot10^{-3}\).  Hence on \(N\) the active source
component vanishes, while the inactive source component is dominated by the
positive Volterra-Schur block.  This proves \(a\ge0\) on \(N\).

For the cross form, the active source rows already factor through the closed
trace range.  The inactive residual is \(A\)-bounded by the same high-block
estimate; polarization and Cauchy--Schwarz give the bounded
\(A^{1/2}\)-factorization of \(b\).  This is precisely the
Douglas/Moore--Penrose condition.

Therefore the quotient theorem supplies bounded operators \(G\) and \(S\), with
\(S\) acting on the transported trace range \(X_R=R(U)\), such that
\[
  Q_\Phi(f)=\|Gf\|^2-\|S R_{\rm global}f\|^2.
\]
In particular
\[
  R_{\rm global}f=0\qquad\Longrightarrow\qquad Q_\Phi(f)\ge0.
\]
Equivalently, the singular range condition and Moore--Penrose Schur form are
positive for the normalized full-\(\Phi\) Weyl/Volterra certificate.

\begin{problem}[External equivalence verification]
The quotient-Schur certificate above is closed inside the normalized
Weyl/Volterra model.  The remaining task is to verify which external
equivalences from this certificate to the original RH positivity formulation
are already proved in the notes and which still require a separate statement.
\end{problem}

The verification gives the following dependency status.  The Riemann kernel formula,
the harmless scalar normalization of \(\xi\), the \(hbar=1\) KLM/Weyl
convention, the transport from the phase-space symbol to the coordinate Weyl
kernel, and the even/odd half-line reduction are all closed algebraic links.
The normalized full-\(\Phi\) quotient Schur certificate and the full-continuum
source-inactive passage are also closed.

The original all-parameter Weyl/KLM/RH-facing chain is not yet closed.  Two
links remain before the original Weyl kernel positivity statement follows.
First, one must prove the quotient-to-original Weyl lift:
\[
  \text{the original }K_\omega\text{ quadratic form}
  =
  \text{the closed trace quotient form}
\]
after the parity reduction, Volterra transformations, and density/closure
limits.  Second, the present generated full-\(\Phi\) certificates record the
stress value \(\omega=0.49\), while the original target is
\(|\omega|<1/2\); this requires either a uniform-in-\(\omega\) theorem or a
certified parameter cover.  After these two links, the standard KLM/Weyl
equivalence gives the original KLM positive-type condition.  A separate final
bridge is still needed from that KLM/Weyl positivity statement to the intended
de Branges or RH-side formulation.

\begin{problem}[Quotient-to-original Weyl lift]
Prove that the closed trace quotient form produced by the Volterra Schur
certificate is exactly the quadratic form of the original coordinate Weyl
kernel \(K_\omega\), after the parity decomposition and all density/closure
limits.  This is now the first external equivalence gap.
\end{problem}

The algebraic part of this lift is closed.  The unitary parity reduction
transports the original full-line form to the two half-line parity forms.  If
\[
  H_\pm=\partial_x\partial_y P_\pm
\]
and the parity kernels decay at infinity, then
\[
  P_\pm(x,y)=\int_x^\infty\int_y^\infty H_\pm(u,v)\,du\,dv.
\]
Thus, for finite sums and then by density,
\[
  \sum_{i,j}c_i c_jP_\pm(x_i,x_j)
  =
  \iint H_\pm(u,v)F(u)F(v)\,du\,dv,
  \qquad
  F(u)=\sum_i c_i{\bf 1}_{u\ge x_i}.
\]
For smooth compact tests the same formula uses
\[
  F(u)=\int_0^u f(x)\,dx.
\]
The full-line mixed Green formula, the Volterra representation, and the
positive multiplication/log-coordinate normalization are invertible
quadratic-form transformations.  Hence the original lift is now reduced to one
endpoint compatibility identity.

The apparent shortcut \(R_{\rm global}F=0\) for all primitive original tests is
false.  For any active \(a\), \(\Lambda_a\) is a nonzero order-eight jet
functional.  A compactly supported smooth primitive \(F\) can realize a jet
with \(\Lambda_a(F)\ne0\), and then \(f=F'\) is still a compact smooth original
test.  Thus the primitive image is not contained in \(\ker R_{\rm global}\).

The remaining route is a boundary-repair comparison.  The quotient theorem
gives
\[
  Q_\Phi(F)=\|GF\|^2-\|S R_{\rm global}F\|^2,
\]
but the abstract repair \(S\) is not canonical: if
\[
  Q=P-R^\ast D R,\qquad P\ge0,\quad D\ge0,
\]
then
\[
  Q=(P+R^\ast T R)-R^\ast(D+T)R
\]
for every positive trace operator \(T\).  Hence the original Weyl boundary term
must be derived from the exact Green identity.  The needed theorem is:
construct a canonical boundary operator \(D_{\rm bdy}\) such that
\[
  \mathcal Q_{K_\omega}(f)
  =
  Q_\Phi(F)+\langle D_{\rm bdy}R_{\rm global}F,R_{\rm global}F\rangle,
\]
and prove \(D_{\rm bdy}=D_q\), or at least \(D_{\rm bdy}\ge D_q\), where
\(D_q\) is the minimal Douglas/Schur repair operator.  Then
\[
  \mathcal Q_{K_\omega}(f)
  =
  \|GF\|^2+
  \langle (D_{\rm bdy}-D_q)R_{\rm global}F,
          R_{\rm global}F\rangle
  \ge0.
  \]
This boundary comparison is now the exact remaining lift theorem.

Thus the two naive alternatives should be retired in this form.  The statement
\[
  R_{\rm global}F=0
\]
for all primitive original tests is false by compact jet extension.  The
statement with ``the same \(S\) from the quotient theorem'' is not well-defined,
because the quotient repair is nonunique.  The invariant replacement is the
comparison of the canonical Green-boundary operator \(D_{\rm bdy}\) with the
minimal quotient/Douglas repair \(D_q\).

The invariant form of the remaining comparison is as follows.  Let \(X_R\) be
the completed transported trace range and let
\[
  Q_\Phi(F)=\|G_qF\|^2-\langle D_qR_{\rm global}F,
  R_{\rm global}F\rangle_{X_R},
  \qquad
  D_q=(\Gamma^\ast\Gamma-C)_+ .
\]
Put
\[
  \beta_{\rm bdy}(F,H)
  =\mathcal Q_{K_\omega}(F,H)-Q_\Phi(F,H).
\]
The exact Green-boundary form defines a canonical operator \(D_{\rm bdy}\) iff
\(\beta_{\rm bdy}\) descends through \(R_{\rm global}\), i.e. iff it vanishes
when one argument lies in \(\ker R_{\rm global}\) and is bounded in the
transported trace norm.  Then
\[
  \beta_{\rm bdy}(F,H)
  =
  \langle D_{\rm bdy}R_{\rm global}F,
          R_{\rm global}H\rangle_{X_R}.
\]
Consequently
\[
  \mathcal Q_{K_\omega}(F)
  =
  \|G_qF\|^2+
  \langle (D_{\rm bdy}-D_q)R_{\rm global}F,
          R_{\rm global}F\rangle_{X_R}.
\]
Thus \(D_{\rm bdy}\ge D_q\) on \(X_R\) is a sufficient theorem for original
Weyl positivity, and equality is the sharp boundary-repair identity.  The
closed part is the abstract descent/comparison theorem; the open part is the
analytic construction of \(D_{\rm bdy}\) from the primitive Green identity and
the proof of \(D_{\rm bdy}-D_q\ge0\).

The primitive endpoint bookkeeping sharpens this further.  In the half-line
primitive lift,
\[
  P(x,y)=\int_x^\infty\int_y^\infty H(u,v)\,du\,dv,\qquad
  F(u)=\int_0^u f(x)\,dx,
\]
integration by parts in \(u\) and \(v\) has no endpoint contribution:
\(F(0)=0\) kills the lower endpoint and the decay of \(P\) kills infinity.
Therefore the already identified mixed/Volterra transport satisfies
\[
  \mathcal Q_{K_\omega}(f,g)=Q_\Phi(F,G),
  \qquad
  \beta_{\rm bdy}=0,\qquad D_{\rm bdy}=0.
\]
Thus no hidden positive primitive-boundary repair is available from the
\((x,y)\)-Green identity.  The comparison \(D_{\rm bdy}\ge D_q\) is now
equivalent to \(D_q=0\) on the primitive trace image
\[
  Y=R_{\rm global}\{F:F'=f,\ f \text{ an original compact Weyl test}\}.
\]
The corrected remaining theorem is to prove \(D_q|_Y=0\), or equivalently to
prove \(Q_\Phi\ge0\) directly on this primitive image.

The primitive trace image is not smaller after completion.  Indeed, if
\(F\in C_c^\infty\) on the half-line and \(F(0)=0\), then \(f=F'\) is a
compact smooth original test and \(F(u)=\int_0^u f(x)\,dx\).  Hence the
primitive class contains the compact smooth Volterra core.  By the
density/closure theorem used in the quotient lift and continuity of
\(R_{\rm global}:V\to X_R\),
\[
  \overline{
  R_{\rm global}\{F:F'=f,\ f \text{ compact smooth}\}}
  =X_R .
\]
Since \(D_q\) is bounded on \(X_R\),
\[
  D_q|_Y=0 \quad\Longleftrightarrow\quad D_q=0\text{ on }X_R .
\]
Thus the primitive image does not give a smaller trace constraint that can kill
the quotient repair.  The remaining theorem is either
\[
  \Gamma^\ast\Gamma-C\le0
\]
in the quotient Schur coordinates, equivalently \(D_q=0\), or a direct proof
that \(Q_\Phi\ge0\) on the full primitive/form closure.

Finite quotient diagnostics support this repair-free alternative.  In the
Galerkin quotient coordinates the relevant Schur defect is
\[
  H_q=\Gamma^\ast\Gamma-C .
\]
For the full \(\widetilde\Phi_3\) model at \(\omega=0.49\), the first two
repair-free scans give
\[
  \lambda_{\max}(H_q)\approx -1.408431551956\cdot10^{-3}
  \quad (N=6,\ \dim R=3),
\]
and
\[
  \lambda_{\max}(H_q)\approx -1.179008958742\cdot10^{-4}
  \quad (N=8,\ \dim R=5).
\]
The next stress point gives
\[
  \lambda_{\max}(H_q)\approx -1.282002952548\cdot10^{-5}
  \quad (N=10,\ \dim R=7).
\]
Thus the finite positive part \((H_q)_+\) vanishes in these tests.  This is not
yet a continuum theorem; it identifies the next analytic target as
\[
  \Gamma^\ast\Gamma\le C\quad\text{on }X_R,
\]
or equivalently a direct positivity proof for \(Q_\Phi\) on the completed
primitive/form domain.

The same diagnostics were then rewritten in the positive Schur-complement
orientation
\[
  S=C-\Gamma^\ast\Gamma .
\]
Transporting this finite Schur complement to sampled trace coordinates gives
\[
  D_{\rm tr}=W\Sigma^{-1}S\Sigma^{-1}W^\ast,\qquad
  RU=W\Sigma .
\]
This raw sampled trace matrix is not the continuum proof norm, but it exposes
the trace-side kernel that a direct proof must factor.  The extended stress
sequence gives
\[
\begin{array}{c|c|c|c}
N&\dim R&\lambda_{\min}(S)&\lambda_{\min}(D_{\rm tr})\\
\hline
6&3&1.408431551956\cdot10^{-3}&1.028656451347\cdot10^{-4}\\
8&5&1.179008958742\cdot10^{-4}&4.687682255695\cdot10^{-7}\\
10&7&1.282002952548\cdot10^{-5}&2.690222591743\cdot10^{-9}\\
12&9&1.081196935196\cdot10^{-6}&1.860543932185\cdot10^{-11}
\end{array}
\]
Thus the finite Schur complement remains positive through the \(12/9\) stress
point, while the raw trace-coordinate representation becomes very
ill-conditioned.  The weakest trace profiles are smooth endpoint-decaying
modes; exponential fits over the four rows give approximate decay rates
\[
  4.3606,\quad 4.4388,\quad 4.5616,\quad 4.7665 .
\]
The next continuum theorem should therefore be stated in the transported
\(X_R\) norm: derive the exact trace-side Schur kernel and prove
\[
  D_{\rm tr}\ge0
\]
as a Volterra/Green Gram kernel.  The endpoint-decaying weak mode is the finite
bottleneck that this formula must dominate.

The exact continuum object is intrinsic.  Let \(V\) be the completed
Weyl/Volterra form domain, let \(R:V\to X_R\) be the completed trace map, and
set \(N=\ker R\).  For any section \(J:X_R\to V\), \(RJ=I\), write
\[
  a(n,m)=Q(n,m),\qquad b(n,x)=Q(n,Jx),\qquad c(x,y)=Q(Jx,Jy).
\]
The quotient theorem gives
\[
  b(n,x)=\langle A^{1/2}n,\Gamma_J x\rangle .
\]
Define
\[
  D_{\rm tr}(x,y)
  =
  c(x,y)-\langle \Gamma_Jx,\Gamma_Jy\rangle .
\]
This definition is independent of \(J\).  Indeed, if \(J'=J+h\), with
\(h:X_R\to N\), then
\[
  \Gamma_{J'}=\Gamma_J+A^{1/2}h,
\]
and the added terms cancel in \(c'-\Gamma_{J'}^\ast\Gamma_{J'}\).  The finite
section-change check in the \(6/3\) quotient model gives Frobenius error
\(7.26\cdot10^{-63}\), and the sampled trace-transport quadratic identity has
error \(3.12\cdot10^{-56}\).

Completing the square gives the constrained-energy identity
\[
  D_{\rm tr}(x,x)=\inf\{Q(f):Rf=x\}.
\]
Consequently \(D_{\rm tr}\ge0\) on \(X_R\) is equivalent to positivity of
\(Q_\Phi\) on the completed form domain.  The remaining analytic step is not
the definition of \(D_{\rm tr}\), but an explicit Volterra/Green Gram formula:
construct \(G_x\) and a positive measure or positive operator \(d\mu\) such that
\[
  D_{\rm tr}(x,y)=\int G_x(u)G_y(u)\,d\mu(u)
\]
in the transported \(X_R\) norm.

The Euler--Lagrange minimizer itself is explicit.  For \(f=n+Jx\), \(n\in N\),
\[
  Q(n+Jx)=\langle An,n\rangle
  +2\operatorname{Re}\langle Bx,n\rangle
  +\langle Cx,x\rangle .
\]
Variations \(h\in N\) give
\[
  \langle An_x+Bx,h\rangle=0\qquad(h\in N),
\]
so the canonical minimizer is
\[
  n_x=-A^+Bx,\qquad f_x=Jx-A^+Bx .
\]
Substitution gives
\[
  Q(f_x)=\langle (C-B^\ast A^+B)x,x\rangle
  =
  \langle D_{\rm tr}x,x\rangle .
\]
The finite \(8/5\) quotient check gives Euler--Lagrange residual
\(1.11\cdot10^{-61}\), trace error \(6.04\cdot10^{-57}\), and energy identity
error \(1.75\cdot10^{-61}\).  Since the finite Schur complement is positive,
the finite minimized energy also has the spectral square representation
\[
  \langle Sz,z\rangle
  =
  \sum_k \left|\sqrt{\lambda_k}\langle v_k,z\rangle\right|^2,
  \qquad S=C-B^\ast A^+B .
\]
This closes the minimizer calculation.  What remains is to identify this square
with an explicit continuum Volterra/Green square by constructing the Green
solver for \(An=-Bx\) and the corresponding residual feature map.

The equivalence is just the Moore--Penrose Schur identity.  Relative to
\(V=N\oplus U\), \(N=\ker R_{\rm global}\), the quotient hypotheses give
\[
  b(n,u)=\langle A^{1/2}n,\Gamma u\rangle .
\]
Hence
\[
  q(n+u)
  =
  \|A^{1/2}n+\Gamma u\|^2+
  \langle (C-\Gamma^\ast\Gamma)u,u\rangle .
\]
Thus \(D_q=0\), \(\Gamma^\ast\Gamma\le C\), and positivity of \(Q_\Phi\) on the
full completed form domain are equivalent.  Since the primitive calculation gives
\(D_{\rm bdy}=0\), this is now the direct original-lift theorem.

\section{Conclusion}

The work so far has removed several attractive but false paths.  Positive
anti-Wick density fails locally.  Hermite-Biehler shortcuts fail for finite
theta cores.  Local source positivity fails.  A first-order score integration
by parts fails.  Exact finite negative index for the anti-Loewner boundary is
false.

What remains is no longer the finite-core source-tail problem.  In the
normalized full-\(\Phi\) Weyl/Volterra model, the quotient Schur certificate is
closed: active source rows factor through the trace range, the inactive source
tail is absorbed in the full continuum high block, and the
Douglas/Moore--Penrose Schur hypotheses are verified.  The remaining work is an
external equivalence problem.  The algebraic quotient-to-original Weyl lift is
closed conditionally; the remaining point is the endpoint trace compatibility
or, equivalently, the exact boundary-repair identity for
\(\|S R_{\rm global}F\|^2\).  After that, one must upgrade the stress-value
certificate at \(\omega=0.49\) to the full interval \(|\omega|<1/2\), either by
a uniform theorem or a certified parameter cover.  Only after those two links
does the established KLM/Weyl normalization give the original quantum
positive-type condition.  A final separate bridge to the intended de Branges or
RH-side formulation is still required.

\appendix

\section{Numerical status summary}

The following table separates observations by logical status.

\begin{center}
\begin{tabular}{p{0.28\linewidth}p{0.60\linewidth}}
\toprule
Status & Statement\\
\midrule
Proved identity &
Second-order theta identity
$\phi_n=(\partial_t^2-1/4)(e^{t/2}e^{-\pi n^2e^{2t}})$.\\
Proved identity &
Exact same-sign finite-core formula \eqref{eq:samesign}.\\
Proved identity &
Volterra boundary-plus-tail formula \eqref{eq:Volterra}.\\
Proved lemma &
RKHS derivative-kernel Gram lift.\\
Disproved route &
Positive anti-Wick density, by local heat-deconvolution obstruction.\\
Disproved route &
Finite-core Hermite-Biehler shortcut.\\
Disproved route &
Layerwise local source positivity.\\
Disproved route &
First-order score integration by parts after splitting $\cosh$.\\
Disproved route &
Exact two-negative-square anti-Loewner boundary hypothesis.\\
Numerical evidence &
KLM matrices positive to roundoff on $15\times 15$ and $21\times 21$ lattice
sweeps.\\
Numerical evidence &
Half-line parity kernels positive to roundoff and relative contraction
spectrum contained in $[-1,1]$ to roundoff.\\
Proved certificate &
Normalized full-\(\Phi\) Weyl/Volterra quotient Schur certificate.\\
Open external link &
Endpoint trace compatibility, or equivalently the exact boundary-repair
identity for the conditional quotient-to-original Weyl lift.\\
Open external link &
Uniform \(\omega\)-coverage for the original target \(|\omega|<1/2\).\\
Open external link &
Final bridge from the resulting KLM/Weyl positivity statement to the intended
de Branges or RH-side formulation.\\
\bottomrule
\end{tabular}
\end{center}

\section{KLM numerical checks}

For $\omega=0.49$ and $\hbar=1$, direct KLM lattice sweeps on $15\times 15$
and $21\times 21$ grids produced worst minimum eigenvalues around
$-7.4\cdot 10^{-15}$, consistent with roundoff.  These tests do not prove
positivity, but they rule out many simple local counterexamples.

\section{Trace Schur Kernel as a Volterra--Green Feature Form}

The repair-free reduction leaves the intrinsic trace-side Schur kernel
\[
  D_{\rm trace}(x,y)=Q(f_x,f_y),
\]
where \(f_x\) is the Euler--Lagrange minimizer in the fiber \(Rf=x\).  In the
block notation \(V=N\oplus U\), this minimizer is
\[
  f_x=Jx-A^+Bx,
  \qquad
  D_{\rm trace}=C-B^*A^+B .
\]

The direct reduced Volterra formula identifies an explicit continuum feature
family.  Put
\[
  A_s(u)=\frac{\Psi(s+u)}{\Psi(s)},\qquad
  B_\sigma(s,u)=e^{\sigma\omega(s+u)/2}A_s(u),
\]
and, for the Green minimizer \(f_x\),
\[
  M_{\sigma,x}(u)=\int f_x(s)B_\sigma(s,u)\,ds,\qquad
  N_{\sigma,x}(u)=\int (s+u)f_x(s)B_\sigma(s,u)\,ds .
\]
Then
\[
  D_{\rm trace}(x,y)
  =
  \frac12\sum_{\sigma}w_\sigma
  \int_0^\infty
  \bigl(
    M_{\sigma,x}(u)N_{\sigma,y}(u)
    +N_{\sigma,x}(u)M_{\sigma,y}(u)
  \bigr)\,du .
\]
This is the requested Volterra--Green feature identification, but it is not
yet an honest Hilbert Gram square.  The pointwise completion
\[
  M_xN_y+N_xM_y
  =
  \frac12\{(M_x+N_x)(M_y+N_y)-(M_x-N_x)(M_y-N_y)\}
\]
shows that the immediate representation is a signed Krein square.  Explicitly,
with
\[
  G_{\sigma,+,x}(u)=\sqrt{w_\sigma/4}\{M_{\sigma,x}(u)+N_{\sigma,x}(u)\},
  \qquad
  G_{\sigma,-,x}(u)=\sqrt{w_\sigma/4}\{M_{\sigma,x}(u)-N_{\sigma,x}(u)\},
\]
we have
\[
  D_{\rm trace}(x,y)
  =
  \sum_\sigma\int_0^\infty
  \bigl(
    G_{\sigma,+,x}(u)G_{\sigma,+,y}(u)
    -
    G_{\sigma,-,x}(u)G_{\sigma,-,y}(u)
  \bigr)\,du .
\]
Therefore the remaining analytic problem is a constrained moment theorem on
the Green minimizer trace image: the negative square must be dominated by the
positive square, or an additional Volterra transform must convert the signed
representation into a positive one.

The script \texttt{trace\_volterra\_green\_feature\_map.py} checks the finite
\(8/5\) quotient against the direct Volterra operator.  The default run gives
relative moment--Schur error \(3.21\cdot10^{-10}\).  A denser run with
\(s\)-order \(80\) and \(u\)-order \(220\) gives relative error
\(1.64\cdot10^{-11}\), Schur minimum \(2.94\cdot10^{-9}\), and matching
Volterra-moment minimum \(2.94\cdot10^{-9}\).  Thus the feature identity is
stable in finite sections; the positive-square theorem remains the open
continuum step.

The positive-square problem can now be stated as a contraction theorem.  Let
\[
  P=\sum_\sigma\int G_{\sigma,+}^\ast G_{\sigma,+}\,du,\qquad
  M=\sum_\sigma\int G_{\sigma,-}^\ast G_{\sigma,-}\,du .
\]
Then \(D_{\rm trace}=P-M\), and the desired Hilbert Gram theorem is equivalent
to
\[
  M\le P,\qquad \lambda_{\max}(P^+M)\le 1
\]
on the Green-minimizer trace image.  The script
\texttt{volterra\_feature\_contraction.py} tests this finite contraction in the
direct Volterra quotient model.  For \((\text{basis},\text{constraints})\)
\((6,3),(8,5),(10,7)\) the top generalized eigenvalues are approximately
\[
  0.7429024494,\qquad 0.8772593736,\qquad 0.9423336439.
\]
A denser \((12,9)\) stress row gives
\[
  \lambda_{\max}(P^+M)\approx 0.9748428017,
  \qquad 1-\lambda_{\max}\approx 0.0251571983 .
\]
Thus the finite evidence supports domination, but the constant appears sharp:
the continuum proof should be a Hardy--Volterra transport identity
\[
  G_- = T G_+,\qquad \|T\|\le 1,
\]
with endpoint boundary terms killed by the Euler--Lagrange trace equations.

There is an exact branch transport handle.  Since
\[
  B_\sigma(s,u;\omega)
  =
  e^{\sigma\omega(s+u)/2}\frac{\Psi(s+u)}{\Psi(s)},
\]
we have
\[
  \partial_\omega M_\sigma(u)
  =
  \frac{\sigma}{2}N_\sigma(u),
  \qquad
  N_\sigma=2\sigma\,\partial_\omega M_\sigma .
\]
The script \texttt{volterra\_transport\_identity.py} verifies this identity
numerically on the \(12/9\) Green-minimizer trace image, with centered
\(\omega\)-difference relative error \(6.31\cdot10^{-8}\).  It also stacks the
finite plus/minus profiles and solves the range problem
\[
  H_-=T H_+.
\]
For the same \(12/9\) row, the relative range residual is
\(4.72\cdot10^{-12}\), and the computed finite range-map norm is
approximately \(0.9830054426\).  The top generalized eigenvalue is sensitive to
near-null directions of the plus-profile Gram, but the robust conclusion is
that the plus-profile range determines the minus-profile range and that the
missing continuum step is to identify this range map with a closed-form
Hardy--Volterra contraction.

The closed-form multiplier has now been isolated.  With \(r=s+u\),
\[
  \kappa(s,u)=\frac{1-r}{1+r}
\]
satisfies \(|\kappa(s,u)|\le 1\) on the Volterra quadrant.  Before the
\(s\)-integration, the lifted signed features obey
\[
  \widetilde G_-=\kappa\,\widetilde G_+ .
\]
Therefore the finite range map has the compressed form
\[
  T=C K E,
\]
where \(K\) is multiplication by \(\kappa\), \(C\) is Volterra integration in
\(s\), and \(E\) is the Green-minimizer right inverse from the observed plus
profile to the lifted plus integrand.  The script
\texttt{volterra\_hardy\_transport\_derivation.py} checks this compressed
identity on the \(12/9\) model: the multiplier satisfies
\(\|\kappa\|_\infty\approx0.9935735362\) on the truncated quadrature window,
the compressed identity error is \(1.50\cdot10^{-12}\), and the finite
compressed map has norm about \(0.9830053111\).  Thus the remaining theorem is
precisely the endpoint statement that the Euler--Lagrange Green lift \(E\)
makes \(CKE\) contractive; equivalently, the boundary concomitant in the
lifted \((s,u)\) integration-by-parts formula must be the trace-fiber
Euler--Lagrange residual \(Q(f_x,h)=0\) for \(h\in\ker R\).

The finite boundary identity is closed by
\texttt{green\_lift\_boundary\_theorem.py}.  In the direct Volterra quotient
model, for every finite \(h\in\ker R\),
\[
  \langle G_+(f_x),G_+(h)\rangle
  -
  \langle G_-(f_x),G_-(h)\rangle
  =
  Q(f_x,h)=0 .
\]
On the \(12/9\) stress row, the relative boundary-concomitant norm is
\(2.15\cdot10^{-10}\), the Frobenius residual is \(1.35\cdot10^{-13}\), and
the perturbative minimality identity
\[
  Q(f_x+h)-Q(f_x)=Q(h)
\]
holds to about \(2.94\cdot10^{-9}\).  Thus the finite boundary term is exactly
the Euler--Lagrange trace-fiber residual.  The remaining analytic theorem is
the continuum closure statement: smooth lifted tests must be dense in the
completed trace-fiber domain, the \(CKE\) boundary concomitant must be
continuous in the Volterra form norm, and the finite identity must pass to
all \(h\in\ker R\).

The closure step is now separated from the finite calculation.  Let \(D\) be
the smooth lifted Volterra core and complete it in the graph norm
\[
  \|f\|_V^2=\|G_+f\|^2+\|G_-f\|^2+\|Rf\|_X^2 .
\]
Then \(D\) is dense in \(V\), the trace map is continuous, and
\(N=\ker R\) is closed.  The forms
\[
  P(f,g)=\langle G_+f,G_+g\rangle,\qquad
  M(f,g)=\langle G_-f,G_-g\rangle,\qquad
  Q=P-M
\]
extend continuously to \(V\).  The lifted integration-by-parts boundary
concomitant
\[
  B(f,g)=\langle G_+f,G_+g\rangle-\langle G_-f,G_-g\rangle
\]
equals \(Q(f,g)\) on \(D\), hence on \(V\) by density and continuity.  Thus,
for the closed Green minimizer \(f_x\) in the fiber \(Rf=x\),
\[
  B(f_x,h)=Q(f_x,h)=0,\qquad h\in\ker R,
\]
where the second equality is precisely the Euler--Lagrange equation in the
closed fiber.  Together with
\[
  \kappa(s,u)=\frac{1-s-u}{1+s+u},\qquad |\kappa(s,u)|\le1,
\]
this gives the completed-domain Green-lift contraction
\[
  \|C K E\|\le1
\]
on the Green-minimizer trace image.  This is the theorem recorded by
\texttt{continuum\_green\_lift\_closure\_theorem.py}.  It closes the Volterra
trace-fiber completion step; the separate quotient-to-original Weyl lift still
requires the primitive endpoint compatibility bookkeeping described above.

That bookkeeping is isolated in
\texttt{primitive\_endpoint\_compatibility\_theorem.py}.  The route asserting
that primitive original tests lie in \(\ker R_{\rm global}\) is false:
compact smooth primitives can have nonzero active endpoint trace.  However,
the primitive transport calculation gives no boundary repair:
\[
  D_{\rm bdy}=0 .
\]
Since the primitive trace image is dense in \(X_R\), it remains only to prove
that the quotient repair vanishes on all of \(X_R\).  On the Green-minimizer
trace image the Schur trace form is
\[
  D_{\rm trace}=P-M,\qquad
  P=\langle G_+,G_+\rangle,\quad M=\langle G_-,G_-\rangle .
\]
The completed Green-lift contraction gives
\[
  G_-=C K E G_+,\qquad \|C K E\|\le1,
\]
hence \(M\le P\).  Therefore
\[
  D_q=(M-P)_+=0,
  \qquad\hbox{equivalently}\qquad
  \Gamma^\ast\Gamma\le C \quad\hbox{on }X_R .
\]
Thus primitive endpoint compatibility is closed in the completed Volterra
model: primitives need not be in \(\ker R_{\rm global}\), but the trace
correction is annihilated because both \(D_{\rm bdy}\) and \(D_q\) vanish.
After this import, the quotient-to-original Weyl lift ledger marks the
algebraic lift, endpoint trace compatibility, and full lift as closed.

The remaining parameter and KLM packaging layers are handled by
\texttt{uniform\_omega\_weyl\_klm\_bridge.py}.  Although several finite
certificates were generated at the stress value \(\omega=0.49\), the final
positivity mechanism is the completed Hardy--Green contraction
\[
  G_-=C K E G_+,\qquad
  K(s,u)=\frac{1-s-u}{1+s+u},\qquad
  \|C K E\|\le1 .
\]
The multiplier \(K\) is independent of \(\omega\).  The \(\omega\)-dependence
enters only through positive branch factors
\[
  \exp\{ {\textstyle\frac12}\sigma\omega(s+u)\},\qquad \sigma=\pm1,
\]
which are uniformly integrable for \(|\omega|<1/2\) because of the
super-exponential theta decay of \(\Phi\).  Hence \(M_\omega\le P_\omega\) for
the whole target range.  Combined with the quotient-to-original lift, this
gives positivity of the original Weyl kernel \(K_\omega\) for every
\(|\omega|<1/2\).  The fixed \(hbar=1\) KLM/Weyl convention then gives the KLM
positive-type condition for \(Q_\omega\) throughout that range.

The final ledger, \texttt{rh\_debranges\_bridge\_ledger.py}, records what is
still missing.  The established theorem is all-\(\omega\) KLM/Weyl positivity,
not yet the Riemann hypothesis.  The remaining bridge must construct an
explicit transform from this KLM positive-type family to an RH-equivalent
de Branges, Hermite--Biehler, or Laguerre--Polya positivity criterion for
\(\xi\), and must prove the required endpoint/positive-cone closure passage.
This is the sole top-level blocker left by the verification.

The endpoint closure part can be separated from the transform construction.
The script \texttt{klm\_debranges\_intertwiner\_attempt.py} records the
natural shifted-\(\Xi\) candidate
\[
  E_\omega(z)=\Xi(z+i\omega),\qquad
  E_\omega^\#(z)=\Xi(z-i\omega),\qquad 0<\omega<1/2,
\]
and the associated de Branges kernel
\[
  K_{E_\omega}(w,z)=
  {E_\omega(z)\overline{E_\omega(w)}
   -E_\omega^\#(z)\overline{E_\omega^\#(w)}
  \over
  2\pi i(\overline w-z)}.
\]
This candidate is useful as a normalization check, but not as a proof by
itself: positivity of \(K_{E_\omega}\) is precisely the Hermite--Biehler
inequality
\[
  |\Xi(z-i\omega)|<|\Xi(z+i\omega)|,\qquad \Im z>0,
\]
which is the RH-facing assertion unless it is obtained from the already proved
KLM/Weyl kernels by a non-circular pullback.

The remaining transform theorem is therefore the explicit identity
\[
  K_{E_\omega}=T_\omega^\ast KLM_\omega T_\omega,
\]
or an equivalent statement that \(K_{E_\omega}\) is a closed positive-cone
limit of such pullback kernels.  Once such an identity is available, endpoint
positivity follows by a closed-cone argument: for any finite test set,
positive Gram matrices \(G_n\) converging entrywise to \(G_\ast\) have
\[
  c^\ast G_\ast c=\lim_n c^\ast G_n c\ge0,
\]
so \(G_\ast\) is positive semidefinite.  The same dense-core argument for
quadratic forms, followed by form closure, gives the endpoint theorem on the
completed form domain.  Thus the endpoint positive-cone closure is closed as
an abstract theorem; the non-circular KLM-to-de Branges intertwiner remains
the single open bridge.

The first finite pullback diagnostic is recorded in
\texttt{klm\_debranges\_pullback\_probe.py}.  On a finite phase-space grid it
forms
\[
  K_{\rm KLM}(p,q)=Q_\omega(p-q)
  \exp\!\left\{ {\ii\over2}\sigma(p,q)\right\},
  \qquad p=(s,t),
\]
and compares \(K_{E_\omega}\) with pullbacks
\[
  \alpha_w^\ast K_{\rm KLM}\alpha_z
\]
for several Gaussian/coherent ansatzes.  At \(\omega=0.49\), on a \(7\times7\)
phase grid and widths \(0.45,0.65,0.85,1.10,1.40\), no tested ansatz gave an
exact pullback or a positive residual.  The best residual was the
anti-Bargmann convention at width \(0.85\), with relative Frobenius residual
approximately \(5.19\cdot10^{-1}\) and indefinite residual spectrum
\[
  \lambda_{\min}\simeq -2.80\cdot10^{-3},\qquad
  \lambda_{\max}\simeq 3.42\cdot10^{-3}.
\]
Thus the missing bridge is not a naive coherent-packet embedding.  The next
target is to derive the canonical Weyl/Bargmann image of the de Branges
evaluation vector \(k_z\) directly from
\[
  \Xi(z)=\int_{\mathbb R}\Phi(t)e^{\ii zt}\,dt,
\]
and then prove the resulting pullback identity or its closed positive-cone
limit.

The canonical image itself is now explicit.  For
\[
  E_\omega(z)=\Xi(z+i\omega),\qquad
  E_\omega^\#(z)=\Xi(z-i\omega),
\]
set, for \(r\ge0\),
\[
  h_z^+(r)=(2\pi)^{-1/2}E_\omega(z)e^{\ii zr},
  \qquad
  h_z^-(r)=(2\pi)^{-1/2}E_\omega^\#(z)e^{\ii zr}.
\]
Since \(\Im z,\Im w>0\),
\[
  \int_0^\infty e^{\ii(z-\overline w)r}\,dr
  ={1\over \ii(\overline w-z)}.
\]
Consequently
\[
  K_{E_\omega}(w,z)
   =\langle h_z^+,h_w^+\rangle_{L^2(0,\infty)}
    -\langle h_z^-,h_w^-\rangle_{L^2(0,\infty)}.
\]
The script \texttt{klm\_debranges\_canonical\_hardy\_image.py} verifies this
normalization against the direct shifted-\(\Xi\) kernel with relative error
\(1.58\cdot10^{-15}\) on the current sample.  The same sample gives
\[
  \lambda_{\max}\!\left((H_+)^{-1/2}H_-(H_+)^{-1/2}\right)
  \simeq 0.9999999999997,
\]
with margin about \(3.2\cdot10^{-13}\).  Thus the endpoint is numerically
critical.

The closed-cone route is now reduced to a branch-transport theorem.  The
completed Volterra/KLM side already has the contraction
\[
  G_- = C K E G_+.
\]
It remains to construct a unitary or isometric transport \(U\) such that
\[
  U h_z^+=G_+(z),\qquad U h_z^-=G_-(z),
\]
or to prove this relation in the strong closed-cone limit.  That transport,
not further coherent-packet tuning, is the remaining non-circular
KLM-to-de Branges bridge.

The exact and limiting versions of this statement are separated in
\texttt{klm\_debranges\_branch\_transport\_theorem.py}.  The exact statement
is a joint-Gram criterion: on a finite evaluation span, a single isometry
\(U\) with
\[
  U h_z^+=G_+(z),\qquad U h_z^-=G_-(z)
\]
exists if and only if the four joint Gram kernels agree:
\[
\begin{aligned}
  \langle h_z^+,h_w^+\rangle&=\langle G_+(z),G_+(w)\rangle,\\
  \langle h_z^-,h_w^-\rangle&=\langle G_-(z),G_-(w)\rangle,\\
  \langle h_z^+,h_w^-\rangle&=\langle G_+(z),G_-(w)\rangle,\\
  \langle h_z^-,h_w^+\rangle&=\langle G_-(z),G_+(w)\rangle .
\end{aligned}
\]
This closes the abstract \(U\)-criterion but not the construction, because the
concrete evaluation trace map \(z\mapsto x_z\) in the completed Volterra trace
image is still missing.

For the closed-cone version, put \(h_{z,R}^\pm=1_{[0,R]}h_z^\pm\).  Then
\[
  \|h_z^\pm-h_{z,R}^\pm\|^2
  ={ |E_\omega^\pm(z)|^2 e^{-2\Im z\,R}\over 4\pi\,\Im z},
\]
so \(h_{z,R}^\pm\to h_z^\pm\) strongly in \(L^2(0,\infty)\), and the finite
Gram kernels converge entrywise.  On the current sample with \(R=40\), the
joint Gram truncation relative error is \(2.69\cdot10^{-13}\), the signed
kernel truncation relative error is \(1.49\cdot10^{-13}\), and the maximum
entry tail bound is \(7.91\cdot10^{-14}\).  Thus the Hardy-side strong
closed-cone limit is closed.  The remaining KLM-side task is to construct
truncated maps \(x_{z,R}\), or phase-space pullback vectors, whose
Volterra/KLM branch feature Grams converge to these Hardy branch Grams.

The first finite construction attempts are recorded in
\texttt{klm\_debranges\_trace\_map\_constructor.py}.  The script first chooses
finite trace coordinates \(X=(x_z)\) so that
\[
  X^\ast P_{\rm Volterra}X=H_+,
\]
thereby matching the Hardy plus branch exactly, and then searches the
remaining orthonormal freedom for the best joint plus/minus/cross match.  This
does not close the map.  At basis \(8\), constraints \(5\), the joint residual
is \(7.10\cdot10^{-1}\), with plus residual \(3.3\cdot10^{-41}\), cross
residual \(9.85\cdot10^{-1}\), and minus residual \(2.92\cdot10^{-1}\).  At
basis \(10\), constraints \(7\), the joint residual is \(8.51\cdot10^{-1}\).

The second construction fits Hardy branch profiles directly on the Volterra
\(u\)-row grid.  This is worse: at basis \(8\), constraints \(5\), the
plus-only profile residual is \(2.49\) and the combined plus/minus profile
residual is \(4.90\); at basis \(10\), constraints \(7\), these grow to
\(3.83\) and \(6.94\).  Thus the concrete map is neither an arbitrary
Kolmogorov lift of the plus Gram nor the naive identification of the Volterra
\(u\) variable with the Hardy \(r\) variable.  The next target is the actual
trace formula
\[
  x_z(a)=\Lambda_a(\hbox{primitive/evaluation vector associated with }h_z),
\]
followed by a joint Gram comparison for those coordinates.

The direct endpoint-trace experiment is recorded in
\texttt{klm\_debranges\_lambda\_trace\_candidate.py}.  For each elementary
primitive candidate \(f_z\), the script computes
\[
  y_z(a_j)=\Lambda_{a_j}(f_z)
  =\sum_k e_k(a_j){f_z^{(k)}(a_j)\over k!},
\]
then converts these sampled traces into the finite row-space coordinates by
solving \(RUc_z=y_z\).  The tested candidates were \(e^{\ii zs}\),
\(e^{-\ii zs}\), two \(\omega\)-shifted exponentials, and simple
\(\Xi(z\pm i\omega)\)-weighted exponentials.  None matches the joint Hardy
branch Grams.  At basis \(8\), constraints \(5\), the best scaled joint
residual is \(9.69\cdot10^{-1}\); at basis \(10\), constraints \(7\), it is
\(9.90\cdot10^{-1}\).  Hence the primitive evaluation vector is not a bare
exponential jet.  It must be obtained by inverting the Volterra feature map
\[
  G_+(f_z)=h_z^+
\]
with the correct Weyl--Volterra kernel, and only then applying \(\Lambda_a\)
to that \(f_z\).

The finite feature-inversion experiment is recorded in
\texttt{klm\_debranges\_feature\_inverse\_candidate.py}.  For each \(z\) it
solves the sampled coefficient problem \(G_+f_z\simeq h_z^+\), computes the
trace vector \(x_z(a_j)=\Lambda_{a_j}(f_z)=Rf_z\), and compares both the direct
feature Gram and the induced Green-minimizer trace Gram with the Hardy branch
targets.  This finite inverse does not close the bridge.  At basis \(8\),
constraints \(5\), the best scaled direct feature-Gram residual is
\(5.07\cdot10^{-1}\), while the best trace-lifted residual is
\(9.97\cdot10^{-1}\).  At basis \(10\), constraints \(7\), these are
\(5.23\cdot10^{-1}\) and \(9.99\cdot10^{-1}\).  Moreover, the direct plus
feature fit remains far from exact, with refined max fit error about
\(4.74\cdot10^{-1}\).

Thus the current finite evidence rules out the naive row-grid inverse of
\(G_+\).  The remaining bridge must be derived from the continuous adjoint
normal equation
\[
  G_+^\ast G_+ f_z = G_+^\ast h_z^+,
\]
including the exact Weyl--Volterra kernel, branch normalization, and trace
regularity before applying the endpoint map \(\Lambda_a\).

The continuous normal equation is implemented in
\texttt{klm\_debranges\_continuous\_normal\_equation.py}.  In this model
\[
  A_s(u)={\Psi(s+u)\over \Psi(s)},\qquad r=s+u,
\]
and the signed Volterra branch features are
\[
  g_\pm(s,u,\sigma)
  =
  \sqrt{w_\sigma/4}\,
  A_s(u)e^{\sigma\omega r/2}(1\pm r).
\]
The plus normal kernel is therefore
\[
  N_+(s,t)
  =
  \sum_\sigma {w_\sigma\over 4}
  \int_0^\infty
  (1+s+u)(1+t+u)A_s(u)A_t(u)
  e^{\sigma\omega(u+(s+t)/2)}\,du,
\]
and the right side is
\[
  b_z(s)=
  \sum_\sigma \sqrt{w_\sigma/4}
  \int_0^\infty
  (1+s+u)A_s(u)e^{\sigma\omega(s+u)/2}h_z^+(u,\sigma)\,du .
\]
The script solves \(N_+f_z=b_z\) and then computes the actual endpoint
coordinates \(x_z(a_j)=\Lambda_{a_j}(f_z)\).

This also does not close the bridge.  At basis \(8\), constraints \(5\), with
\(s\)-quadrature order \(16\), Laguerre order \(40\), and right-side Laguerre
order \(60\), the normal equation is solved to roundoff, with maximum residual
about \(4.13\cdot10^{-29}\).  Nevertheless the best scaled direct continuous
feature-Gram residual is \(4.82\cdot10^{-1}\), and the best trace-lifted
Green-minimizer residual is \(9.99\cdot10^{-1}\).  Hence the missing
intertwiner is not the plain plus-branch least-squares inverse.  The next
analytic target is an additional Hardy-to-Volterra transform or branch
normalization, or a coupled signed branch system replacing the plus-only normal
equation.

The coupled signed-branch normal equation is implemented in
\texttt{klm\_debranges\_coupled\_branch\_normal\_equation.py}.  It replaces
the plus-only problem by
\[
  (G_+^\ast G_+ + G_-^\ast G_-)f_z
  =
  G_+^\ast T_+h_z + G_-^\ast T_-h_z,
\]
and tests the scalar branch normalizations \(T\) given by direct matching,
branch swap, minus sign, and \(\pm i\) phase rotations.  These do not close the
bridge.  In the main basis \(8\), constraints \(5\) run, with the trace lift
included, the best case is the swapped branch with the
\(\exp(\sigma\omega u/2)\) adjustment.  The normal equation residual is
\(4.23\cdot10^{-29}\), but the best scaled direct continuous feature-Gram
residual is \(4.84\cdot10^{-1}\), and the best trace-lifted residual is
\(9.98\cdot10^{-1}\).

Therefore the missing map is not a scalar branch normalization, branch swap,
sign, or phase correction.  The next object must be a genuine
Hardy--to--Volterra transmutation kernel
\[
  U(r;s,u,\sigma),
\]
mixing the Hardy half-line variable with the Volterra endpoint variables before
the trace map \(x_z(a)=\Lambda_a(f_z)\) is formed.

The first transmutation probe is recorded in
\texttt{klm\_debranges\_transmutation\_kernel\_probe.py}.  It tests the
coarea ansatz
\[
  (U_\theta h)(s,u,\sigma)=\theta(s,u,\sigma)h(s+u),
\]
using the literal relation \(r=s+u\).  The tested scalar weights include
\[
  1,\quad \ell(r)^{-1/2},\quad e^{-\sigma\omega(s+u)/2},\quad
  A_s(u),\quad A_s(u)^{1/2},\quad A_s(u)e^{\sigma\omega(s+u)/2},
\]
and \(1\pm(s+u)\).  This class also fails.  In the low-order run with
integrated forms included, the best lifted triangle residual is
\(6.06\cdot10^{-1}\) and the best integrated Volterra residual is
\(9.26\cdot10^{-1}\).  In the basis \(8\), constraints \(5\), lifted-only
refinement, the best lifted residual is \(5.79\cdot10^{-1}\), attained by the
swapped branch with the coarea weight.

Thus the missing map is not a scalar multiplier on the coarea surface
\(r=s+u\).  The next analytic object must be a nonlocal kernel
\[
  (Uh)(s,u,\sigma)
  =
  \int_0^\infty U(r;s,u,\sigma)h(r)\,dr,
\]
probably obtained from the Mellin--Laplace representation of \(\Xi\) and the
Volterra ratio \(A_s(u)=\Psi(s+u)/\Psi(s)\).

The first nonlocal dictionary probe is
\texttt{klm\_debranges\_nonlocal\_transmutation\_probe.py}.  It tests
half-line resolvent and heat kernels centered at
\[
  c=s+u,\qquad c=u,\qquad
  c=\log(1+e^s(e^u-1)).
\]
For example,
\[
  \int_c^\infty e^{-\lambda(r-c)}e^{izr}\,dr
  ={e^{izc}\over \lambda-iz},
\]
and the backward half-line piece is
\[
  \int_0^c e^{-\lambda(c-r)}e^{izr}\,dr
  ={e^{izc}-e^{-\lambda c}\over \lambda+iz}.
\]
The basis \(4\), constraints \(3\), smoke scan has best lifted residual
\(6.33\cdot10^{-1}\).  A basis \(8\), constraints \(5\), targeted scan gives a
modest improvement: the best candidate is the swapped-branch forward resolvent
centered at \(s+u\), with coarea weight and \(\lambda=2\), giving lifted
residual \(5.40\cdot10^{-1}\).  This is better than the scalar coarea probe but
still far from a bridge.  Hence a generic resolvent or heat dictionary is not
the missing transmutation.

The next derivation should match the exact Mellin expansion of \(\Xi\) against
the finite-core Volterra mode exponentials term by term, rather than choosing
generic kernels.

\section{References and background}

The background for this draft includes de Branges spaces, Weyl quantization,
the KLM quantum Bochner theorem, total positivity, and anti-Loewner kernels.

\subsection{Exact Mellin atom matching and the remaining mode mixing}

The generic nonlocal kernels above still do not use the actual theta-mode
structure.  We therefore tested the literal finite-core dictionary.  If
\(\Psi(v)=\Phi(v/2)\) contains an atom
\[
  a\exp(\beta v-c e^v),
\]
then its positive-side contribution to \(\Xi\) is
\[
  {a\over 2}c^{-(\beta+iz/2)}
  \Gamma(\beta+iz/2,c),
\]
and the even finite-core contribution is obtained by adding the reflected
\(z\mapsto -z\) term.  This gives an exact incomplete-gamma expansion of the
finite-core \(\Xi\).  Numerically, for \(\widetilde\Phi_3\), the atom expansion
agrees with direct finite-core quadrature on the sample nodes with relative
error about \(3.30\cdot 10^{-21}\).

The corresponding Volterra atom in
\[
  A_s(u)={\Psi(s+u)\over\Psi(s)}
\]
is
\[
  \rho_i(s)\exp\{\beta_i u-c_i e^s(e^u-1)\}.
\]
The diagonal matching ``\(\Xi\)-atom \(i\) maps to Volterra atom \(i\)''
nevertheless fails to recover the Hardy branch Gram.  The best basis-8 lifted
joint residual is about \(7.23\cdot 10^{-1}\), and a small integrated
Volterra check gives residual about \(9.15\cdot 10^{-1}\).  Thus the Mellin
normalization is correct, but the transmutation is not diagonal in this atom
basis.  The next missing structural object is a finite mode-mixing matrix
between the incomplete-gamma \(\Xi\)-atoms and the Volterra ratio atoms.

We next derived this finite mode-mixing matrix in the most direct Galerkin
form.  Let \(a(z)\) be the six-vector of incomplete-gamma atom amplitudes and
let \(B_\epsilon(s,u,\sigma)\) be the six-column matrix of Volterra ratio
atoms.  The finite least-squares problem is
\[
  B_\epsilon(s,u,\sigma) M a_{\rm branch}(z)
  \simeq h_{{\rm branch},z}(s,u,\sigma).
\]
For \(\widetilde\Phi_3\), the best basis-8 row fit among the tested natural
centers is about \(6.33\cdot 10^{-1}\), while the best lifted joint Gram
residual is about \(6.19\cdot 10^{-1}\).  The latter is a modest improvement
over the diagonal atom dictionary, but it is still not a bridge.  Moreover the
best lifted matrix has Frobenius norm about \(1.90\cdot 10^{11}\), indicating
large cancellations rather than a stable structural transform.  The remaining
object should therefore be derived from the exact Mellin convolution identity,
with the incomplete-gamma boundary terms included, rather than from row-level
least squares.

The exact Mellin convolution identity makes this failure structural.  For a
single positive-side atom
\[
  f_i(v)=a_i\exp(\beta_i v-c_i e^v),\qquad
  \alpha_i=\beta_i+iz/2,
\]
we have
\[
  X_i(z)= {1\over 2}a_i c_i^{-\alpha_i}\Gamma(\alpha_i,c_i).
\]
Splitting at the Volterra base point \(s\) gives
\[
  X_i(z)=B_i(s,z)+T_i(s,z),
\]
where
\[
  B_i(s,z)= {1\over 2}a_i c_i^{-\alpha_i}
  \int_{c_i}^{c_i e^s} y^{\alpha_i-1}e^{-y}\,dy
\]
and
\[
  T_i(s,z)= {1\over 2}\Psi(s)e^{izs/2}
  \int_0^\infty V_i(s,u)e^{izu/2}\,du.
\]
Here
\[
  V_i(s,u)=\rho_i(s)\exp\{\beta_i u-c_i e^s(e^u-1)\}
\]
is precisely the \(i\)-th normalized Volterra ratio atom in
\(\Psi(s+u)/\Psi(s)\).  The implemented certificate verifies this identity
to about \(10^{-69}\) in the tail term.  The boundary term is large: by
\(s=0.545\) the smallest sampled tail fraction is about \(9.86\cdot 10^{-10}\),
while the boundary prefix can be essentially the full atom.  Thus the correct
intertwiner cannot be a constant finite atom matrix.  It must route the
incomplete-gamma boundary prefix through the endpoint/trace part of the
Volterra fiber and use the diagonal Volterra identity only on the moving tail.

We also tested whether the existing endpoint trace family \(\Lambda_a\) already
represents this boundary prefix.  In the finite quotient model we form the row
\[
  E_B(z)_k=\int_0^L B(s,z)p_k(s)\,ds
\]
and solve \(E_B(z)\simeq c_z^T R\), where
\((Rf)_j=\Lambda_{a_j}(f)\).  The test is negative: for basis \(10\) with
seven sampled traces, the maximum row-span residual is about \(7.16\cdot
10^{-1}\), and the maximum null-energy residual on the trace nullspace is
about \(7.29\cdot 10^{-1}\).  Therefore the old endpoint defect trace does
not cancel the Mellin boundary prefix.  The bridge requires a new
Mellin-boundary concomitant, or an equivalent augmented trace repair, before
the diagonal Volterra tail identity can be used to complete the
Hardy-to-Volterra transmutation.

The concomitant is explicit.  Since
\[
  {d\over ds}B_i(s,z)
  ={1\over 2}a_i\exp(\beta_i s-c_i e^s)e^{izs/2},
  \qquad B_i(0,z)=0,
\]
the missing trace is the Volterra primitive functional
\[
  \Mu_z(f)=\int_0^L B(s,z)f(s)\,ds .
\]
If \(F(s)=\int_s^L f(t)\,dt\), then integration by parts gives
\[
  \Mu_z(f)=\int_0^L B'(s,z)F(s)\,ds,
\]
with no endpoint contribution because \(B(0,z)=0\) and \(F(L)=0\).  This is
the missing Mellin-boundary trace/concomitant.  In the finite quotient model,
the old \(\Lambda_a\)-trace leaves a boundary action of about \(7.29\cdot
10^{-1}\) on the trace nullspace, whereas the augmented trace
\((\Lambda_a,\Mu_z)\) kills this boundary prefix in the augmented nullspace.
The remaining theorem is to prove that the augmented trace repair is positive
or annihilated in the KLM/de Branges pullback, so that adding \(\Mu_z\) does
not introduce a negative Schur defect.

The corresponding finite Schur repair is positive.  With
\[
  R_{\rm aug}=(\Lambda_a,\Mu_z),
\]
we construct
\[
  P_{\rm aug}=K+R_{\rm aug}^*D_{\rm aug}R_{\rm aug}
\]
by the Moore--Penrose/Douglas Schur construction.  In the basis-10 model with
seven sampled endpoint traces, the old \(\Mu\)-row action on
\(\ker\Lambda\) is about \(5.55\cdot 10^{-1}\), but the action on
\(\ker(\Lambda,\Mu)\) is zero.  The constructed augmented repair has
positive \(P_{\rm aug}\) with minimum eigenvalue about \(9.48\cdot 10^{-6}\),
and \(D_{\rm aug}\) is nonnegative up to roundoff.  Thus the finite model
supports the intended repair: the augmented trace kills the boundary primitive,
and the remaining diagonal Volterra tail is handled by the Volterra/KLM
positive form.  The remaining issue is to lift this augmented repair to the
completed continuum trace space.

This lift is performed in the transported trace norm.  Let \(V_{\rm aug}\) be
the closed Volterra/Mellin graph domain and define
\[
  \|R_{\rm aug}f\|_{X_{\rm aug}}
  =\inf\{\|f+h\|_{V_{\rm aug}}:h\in\ker R_{\rm aug}\}.
\]
Then \(R_{\rm aug}:V_{\rm aug}/\ker R_{\rm aug}\to X_{\rm aug}\) is unitary.
The functional \(\Mu_z\) is closed by construction in this graph norm, and
the primitive formula above identifies its endpoint concomitant.  On
\(\ker R_{\rm aug}\), the Mellin boundary prefix vanishes, so the exact
Mellin identity leaves only the diagonal Volterra tail; this tail is
nonnegative by the already proved Volterra/KLM positivity theorem.  The
Moore--Penrose quotient theorem in \(X_{\rm aug}\) then gives a bounded
nonnegative repair
\[
  D_{\rm aug}=(\Gamma^*\Gamma-C)_+,\qquad
  K+R_{\rm aug}^*D_{\rm aug}R_{\rm aug}\ge 0 .
\]
Because all terms are continuous in the augmented closed graph norm, the
positive form extends from the Galerkin/core domain to the completed
augmented trace-fiber domain.  Thus the augmented repair layer is closed in
the continuum.  At this stage the remaining item was the final de Branges
evaluation pullback, or a closed-cone limiting construction, using
\(R_{\rm aug}=(\Lambda,\Mu)\).

\begin{thebibliography}{10}

\bibitem{deBranges}
L. de Branges,
\emph{Hilbert Spaces of Entire Functions},
Prentice-Hall, 1968.

\bibitem{Folland}
G. B. Folland,
\emph{Harmonic Analysis in Phase Space},
Princeton University Press, 1989.

\bibitem{Karlin}
S. Karlin,
\emph{Total Positivity},
Stanford University Press, 1968.

\bibitem{Pinkus}
A. Pinkus,
\emph{Totally Positive Matrices},
Cambridge University Press, 2010.

\bibitem{AissenSchoenbergWhitney}
M. Aissen, I. J. Schoenberg, and A. Whitney,
On the generating functions of totally positive sequences,
\emph{Journal d'Analyse Mathématique} \textbf{2} (1952), 93--103.

\bibitem{SchoenbergPF}
I. J. Schoenberg,
On Pólya frequency functions. I. The totally positive functions and their
Laplace transforms,
\emph{Journal d'Analyse Mathématique} \textbf{1} (1951), 331--374.

\bibitem{Kastler}
D. Kastler,
The $C^*$-algebras of a free boson field,
\emph{Communications in Mathematical Physics} \textbf{1} (1965), 14--48.

\bibitem{LoupiasMiracleSoleI}
G. Loupias and S. Miracle-Sole,
$C^*$-algebres des systemes canoniques. I,
\emph{Communications in Mathematical Physics} \textbf{2} (1966), 31--48.

\bibitem{LoupiasMiracleSoleII}
G. Loupias and S. Miracle-Sole,
$C^*$-algebres des systemes canoniques. II,
\emph{Annales de l'Institut Henri Poincare. Section A, Physique Theorique}
\textbf{6} (1967), no. 1, 39--58.

\bibitem{Moyal}
J. E. Moyal,
Quantum mechanics as a statistical theory,
\emph{Proceedings of the Cambridge Philosophical Society}
\textbf{45} (1949), 99--124.

\bibitem{Polya}
G. Polya,
Bemerkung ueber die Integraldarstellung der Riemannschen $\xi$-Funktion,
\emph{Acta Mathematica} \textbf{48} (1926), 305--317.

\bibitem{Audenaert}
K. M. R. Audenaert,
A characterisation of anti-Loewner functions,
\emph{Proceedings of the American Mathematical Society}
\textbf{139} (2011), no. 12, 4217--4223; arXiv:1008.2943.

\end{thebibliography}

\end{document}
